\documentclass[11pt,a4paper]{article}

\usepackage[utf8]{inputenc}
\usepackage{amsmath,amssymb}
\usepackage{mathtools}
\usepackage{graphicx}
\usepackage{booktabs}
\usepackage[margin=1in]{geometry}
\usepackage{authblk}
\usepackage{siunitx}
\usepackage[hidelinks]{hyperref}
\usepackage{caption}
\usepackage{xcolor}
\usepackage{microtype}
\usepackage{csquotes}
\usepackage[numbers,sort&compress]{natbib}

% ---- float placement (same rationale as Paper IV) ----
\usepackage[section]{placeins}
\renewcommand{\topfraction}{0.92}
\renewcommand{\bottomfraction}{0.70}
\renewcommand{\textfraction}{0.05}
\renewcommand{\floatpagefraction}{0.72}
\setcounter{topnumber}{3}
\setcounter{bottomnumber}{2}
\setcounter{totalnumber}{5}

\allowdisplaybreaks
\setlength{\parindent}{15pt}
\setlength{\parskip}{0.3em}
\linespread{1.03}

\captionsetup{font=small,labelfont=bf}

\graphicspath{{.}{./figures/}{./fig/}{./images/}}

% ---- shared macros (consistent with Papers I--IV) ----
\newcommand{\gmean}{\langle g\rangle}
\newcommand{\RA}{R_{A}}
\newcommand{\RV}{R_{V}}
\newcommand{\Ptwo}{P_{2}}
\newcommand{\Pthree}{P_{3}}
\newcommand{\geff}{\mathbf{g}_{\mathrm{eff}}}
\newcommand{\nhat}{\hat{\mathbf{n}}}
\newcommand{\Upot}{U}
\newcommand{\Ubar}{\bar{U}}
\newcommand{\potvar}{\sigma_{U}}
\newcommand{\potvarn}{\tilde{\sigma}_{U}}
\newcommand{\omegarot}{\omega}
\newcommand{\dens}{\varrho}
\newcommand{\Nf}{N_{f}}

% ---- Paper V macros: two-region density model ----
\newcommand{\densR}{\dens_{R}}          % density of the region
\newcommand{\densrest}{\dens_{\mathrm{rest}}}
\newcommand{\denshead}{\dens_{\mathrm{head}}}
\newcommand{\densbody}{\dens_{\mathrm{body}}}
\newcommand{\Ddens}{\Delta\dens}        % density difference
\newcommand{\Dmass}{\Delta M}           % excess mass
\newcommand{\VR}{V_{R}}
\newcommand{\Vrest}{V_{\mathrm{rest}}}
\newcommand{\Vtot}{V_{\mathrm{tot}}}
\newcommand{\Mtot}{M_{\mathrm{tot}}}
\newcommand{\xsplit}{x_{\mathrm{split}}}
\newcommand{\Ufull}{U_{\mathrm{full}}}
\newcommand{\UR}{U_{R}}
\newcommand{\dCOM}{\Delta x_{\mathrm{COM}}}
\newcommand{\improve}{\mathcal{I}}       % improvement over uniform
% ---- added for ver10: statistics and derived quantities ----
\newcommand{\ampfac}{\mathcal{A}}        % rho_rest V_tot / Delta M
\newcommand{\Sstat}{\mathcal{S}}         % displacement statistic
\newcommand{\CV}{\mathrm{CV}}            % coefficient of variation
\newcommand{\densbulk}{\dens_{\mathrm{bulk}}}
\newcommand{\densneck}{\dens_{\mathrm{neck}}}
\newcommand{\geomslope}{g_{\mathrm{geom}}}

\newcommand{\Dipole}{D}      % mass dipole, D = \Delta M (x_R - x_cf)

\begin{document}
\emergencystretch=2em

% ==================== TITLE BLOCK ====================
\begin{center}
{\LARGE\bfseries What potential-dispersion minimization constrains in
small-body interiors\par}
\vspace{0.8em}
{\itshape Prapon Kanjanatarayont\par}
\vspace{0.2em}
{\small\color{gray} Independent Researcher \textperiodcentered\ Thailand\par}
\vspace{0.1em}
{\small\color{gray} \texttt{prapon.kanjana@gmail.com}\par}
\end{center}
\vspace{1.0em}

%=======================================================================
%  Paper V — Section 1 (Abstract + Introduction)   revised
%  Lines marked  % VERIFY  contain values to confirm against your tables
%=======================================================================
%\documentclass[preprint,12pt]{elsarticle}

%\usepackage[T1]{fontenc}
%\usepackage[utf8]{inputenc}
%\usepackage{amsmath,amssymb}
%\usepackage{graphicx}
%\usepackage{booktabs}
%\usepackage{natbib}
%\usepackage{siunitx}
%\sisetup{per-mode=symbol, range-phrase={--}, range-units=single}
%\usepackage[hidelinks]{hyperref}

%------------------------- notation macros -----------------------------
%\newcommand{\potvarn}{\tilde{\sigma}_{U}}
%\newcommand{\denshead}{\varrho_{\mathrm{head}}}
%\newcommand{\densrest}{\varrho_{\mathrm{rest}}}
%\newcommand{\densR}{\varrho_{R}}
%\newcommand{\densbulk}{\varrho_{\mathrm{bulk}}}
%\newcommand{\Ddens}{\Delta\varrho}
%\newcommand{\Dmass}{\Delta M}
%\newcommand{\VR}{V_{R}}
%\newcommand{\Vtot}{V_{\mathrm{tot}}}
%\newcommand{\Mtot}{M_{\mathrm{tot}}}
%\newcommand{\ampfac}{\mathcal{A}}
%\newcommand{\Ufull}{U^{\mathrm{grav}}_{\mathrm{full}}}
\newcommand{\Ureg}{U^{\mathrm{grav}}_{R}}
\newcommand{\Urot}{U^{\mathrm{rot}}}

%\journal{Icarus}

%\begin{document}

%\begin{frontmatter}

%\title{What potential-dispersion minimization constrains in the interior of a small body}

%\author{...}   % add author block

%=======================================================================
%  ABSTRACT  --  target 250--300 words (this version: 296)
%=======================================================================
{\noindent\large\bfseries Abstract}\\[0.4em]
{\noindent
\citet{richardson2014} argued that on a small body whose surface has
relaxed under downslope regolith flow the bulk density can be estimated
by minimizing the dispersion of the surface geopotential, and
\citet{kanamaru2019} extended the argument from a single mean density to
a two-region interior with the total mass held fixed. Paper~IV
established the property that makes such a minimization meaningful: the
normalized geopotential dispersion $\potvarn$ is far less sensitive to
mesh resolution than the mean surface slope. We ask what the extension
to interior structure actually constrains.

Regions are built by exact plane clipping, verified against closed-form
solids, a complement closure, a vertex predicate and a manifold criterion,
with a fifth containment test where a shape model has more than one
component, and externally against the published lobe volume of 67P. Of five bodies, only Itokawa gives a dispersion
improvement above the $0.9$--$5.9\%$ mesh-resolution sensitivity
envelope of Paper~IV: $16.9\%$ at the detected neck. Eros, Bennu, 67P
and Arrokoth return $<10^{-3}\%$ to $1.3\%$, at or below that envelope.
At the dividing plane $x=\SI{150}{\metre}$ of \citet{kanamaru2019} we
recover $\denshead=2725$ and $\densrest=\SI{1858}{\kilo\gram\per\cubic\metre}$,
a contrast of $1.47$ against their $1.46$, and a centre-of-mass offset of
$16.5\pm\SI{2.9}{\metre}$, consistent with the $\SI{16}{\metre}$ from
equipotential fitting and with the $21\pm\SI{12}{\metre}$ from YORP spin-up
\citep{lowry2014}---though that determination is itself only $1.75$ standard
deviations from zero, and its interval contains every plane of our sweep.

The mass constraint makes the excess mass $\Dmass=\Ddens\VR$ identically
$\Mtot-\densrest\Vtot$, so a statement about one is a statement about the
other. Over the eleven planes of a thirteen-plane sweep on Itokawa that lie
below the grain-density ceiling, $\denshead$ varies by $\pm13.2\%$ and
$\densrest$ by $\pm0.43\%$. In coefficients of variation the ratio is $32$;
the identity contributes $11.8$ of it and the subtraction that defines
$\Ddens$ removes $2.8$, leaving a factor of $7.5$ that reflects genuine
compensation between contrast and region volume. Including the curvature sensitivity interval
gives $\densrest=1858\pm\SI{27}{\kilo\gram\per\cubic\metre}$.
Eight uniform synthetic contact binaries show the invariance emerges
only as the signal strengthens. Dispersion minimization does not
constrain its own dividing plane, and Itokawa fails the first of the
four relaxation conditions. Conservation of the mass dipole is excluded at twenty-one standard errors,
so the agreement with YORP is an independent check and not a second reading
of the same quantity. Dispersion minimization does not constrain its own
dividing plane: left to itself it runs to a head density standing $11.2$
standard deviations above the LL grain density. Itokawa fails the first of
the four relaxation conditions, but that condition governs the conditioning
of a variable-mass inversion and has no counterpart here, and a control run
with the rotational term removed confirms that the fixed-mass objective
retains an interior minimum. On Arrokoth, which has no measured mass, the
sense of the recovered contrast reverses within the published range of bulk
density, between $235$ and $\SI{250}{\kilo\gram\per\cubic\metre}$.
%\end{abstract}


%-----------------------------------------------------------------------

}

\medskip
\noindent\textbf{Keywords:} 
asteroids,
interior structure,
small bodies,
surface geopotential,
polyhedral gravity,
inverse problems,
identifiability,
shape models,
25143 Itokawa,
486958 Arrokoth,
67P/Churyumov-Gerasimenko,
101955 Bennu,
433 Eros

%\end{frontmatter}

%=======================================================================
\section{Introduction}
\label{sec:intro}

The interior of a small body is not directly observable. Where a
spacecraft has tracked the body, radio science fixes the total mass and,
on the most favourable targets, the low-degree gravity field
\citep{yeomans2000,patzold2016,scheeres2020}; but the mass alone is one
number, and the gravity coefficients that would resolve internal
structure are degenerate with the shape unless the spacecraft flies very
close for a long time. Where no spacecraft has visited, even the mass is
absent. Indirect routes are therefore of interest, and one of them uses
the surface itself.

%-----------------------------------------------------------------------
\subsection{Relaxation as an inference tool}
\label{sec:intro:relaxation}

\citet{richardson2014} set out the argument in full. On a rotating
irregular body covered by mobile regolith, downslope flow triggered by
seismic shaking or volatile activity acts as a negative-feedback system,
eroding local topography toward a flat equipotential surface. If that
process has run to completion, the surface geopotential should be nearly
uniform, and the bulk density can be estimated by finding the value that
minimizes its dispersion. They state four conditions under which the
argument holds: the mean rotational force must be a significant fraction
of the mean gravitational force; a sufficiently thick, low-cohesion,
mobile regolith layer must exist over most of the surface; a downslope
flow disturbance source must be present; and sufficient time must have
passed since the last major surface alteration. They are explicit that
the result is neither a direct nor an indirect density measurement, but
the density corresponding to the most eroded state of the body given its
present shape, topography and spin. We return in
Section~\ref{sec:discussion} to the standing of these four conditions on
the bodies examined here, and note at the outset that the single body
that responds to the inversion, Itokawa, does not satisfy the first of
them.

\citet{kanamaru2019} extended the argument from a single mean density to
an interior distribution. Holding the total mass at its measured value
and dividing Itokawa into a head and a body at $x=\SI{150}{\metre}$,
they minimized the area-weighted potential variance over the global
surface with respect to the one remaining free parameter, and recovered
a head density of $2730$ and a body density of
$\SI{1870}{\kilo\gram\per\cubic\metre}$, a contrast of $1.46$. A
companion study using the smooth terrains rather than the whole surface
\citep{kanamaru2019pss} obtained $2450$ and
$\SI{1930}{\kilo\gram\per\cubic\metre}$, a contrast of $1.27$, and,
importantly, found that a model in which the two lobes have comparable
densities separated by a compressed neck fits marginally better than one
in which the lobes differ, and that a three-region model treating head,
neck and body separately is largely unconstrained; only when the
independently estimated centre-of-mass offset from YORP spin-up is
imposed do their models require the head to exceed the global average.
\citet{motooka2012} had earlier attempted to fit the global shape of
Itokawa to a single equipotential surface. An independent line of
evidence comes from \citet{lowry2014}, whose measured YORP spin-up
requires a centre-of-mass displacement of
$\SI{21\pm12}{\metre}$ along the long axis. Because this is the only
constraint on Itokawa's interior not derived from an equipotential
argument, it is the only genuinely external check available to us, and
we treat it as such in Section~\ref{sec:itokawa}.

Paper~IV of this series supplied a property on which any such
minimization depends. Comparing meshes of the same body decimated from
roughly $42{,}000$ to $2{,}000$ facets, the area-weighted mean slope
varied by $1.8\%$ to $25.6\%$ across the test set while the normalized
geopotential dispersion varied by only $0.9\%$ to $5.9\%$, reaching a
plateau between the two finest meshes. The endpoints of these two ranges
are attained on different bodies and should not be divided against one
another; the per-body ratios are given in Paper~IV. A quantity that
changes with the arbitrary resolution of the mesh cannot support an
inference; the dispersion, over the tested range, does not.

Section~\ref{sec:discussion:ellipsoid} tests that comparison against an
analytic reference and finds it, if anything, understated. We refer to
the $0.9$--$5.9\%$ range throughout as the \emph{mesh-resolution
sensitivity envelope}. It measures how far $\potvarn$ moves when the
discretization changes; it is neither a statistical noise level nor a
detection threshold, and we use it only as a scale against which an
improvement may be judged large or small. Constructing a proper null
distribution would require either a noise model for the shape data or a
placebo ensemble of randomly placed dividing surfaces; neither is
attempted here, and the absence is a limitation we state plainly rather
than a gap we conceal.

%-----------------------------------------------------------------------
\subsection{The question this paper asks}
\label{sec:intro:question}

That the extension works on Itokawa has been reported. What it
constrains has not been examined. The dividing plane is chosen by the
investigator, and no published study varies it and reports how the
recovered density responds. The question matters because the density of
a region is not an observable in the way the total mass is: it is the
ratio of an inferred excess mass to a volume that the investigator has
defined.

The point can be made before any computation. With $\Mtot$ and $\Vtot$
held fixed, conservation of mass gives
%
\begin{equation}
\Mtot \;=\; \densR\VR + \densrest\!\left(\Vtot-\VR\right)
      \;=\; \densrest\Vtot + \Ddens\,\VR ,
\label{eq:massbalance}
\end{equation}
%
with $\Ddens\equiv\densR-\densrest$, so that the excess mass
$\Dmass\equiv\Ddens\VR$ satisfies the exact identity
%
\begin{equation}
\Dmass \;=\; \Mtot-\densrest\Vtot
       \;=\; \left(\densbulk-\densrest\right)\Vtot ,
\qquad \densbulk\equiv\Mtot/\Vtot .
\label{eq:excessidentity}
\end{equation}
%
Equation~\eqref{eq:excessidentity} involves no property of the objective
function. It says that $\Dmass$ and $\densrest$ are related by a
bijective affine map, that a two-region model with fixed total mass has
exactly one free parameter, and that the pair $(\densR,\VR)$ is
degenerate along the hyperbola $\Ddens\VR=\mathrm{const}$. Any claim
that one of these quantities is ``more stable'' than another must
therefore be accompanied by the amplification factor relating them,
%
\begin{equation}
\ampfac \;\equiv\; \frac{\densrest\Vtot}{\Dmass}
        \;=\; \frac{\densrest}{\densbulk-\densrest}
        \;=\; \frac{1}{f\,(C-1)},
\qquad f\equiv\VR/\Vtot,\quad C\equiv\densR/\densrest ,
\label{eq:amplification}
\end{equation}
%
which on Itokawa at the detected neck takes the value
$\ampfac=12.0$.
We show in Section~\ref{sec:constrains} that reporting the relative
scatter of $\densrest$ without $\ampfac$ overstates the precision of the
inversion by exactly this factor, and that $\ampfac$ itself is amplified
by $\densbulk/(\densbulk-\densrest)\simeq13$ relative to errors in
$\densrest$, so that a $0.13\%$ shift in the background density moves
$\ampfac$ by $1.7\%$.

We therefore treat the method as an object of study rather than a tool.
Five bodies are analysed---Itokawa, Eros, Bennu, 67P and
Arrokoth---together with eight synthetic contact binaries constructed so
that the strength of the signal can be varied while everything else is
controlled: uniform in density by construction, volumes within $7\%$ of
one another, and every plane sweep spanning the same factor $1.90$ in
region volume. The synthetic experiment is necessary because no further
real body can settle the question: we cannot choose in advance whether
it will carry a signal, and no body's true interior is known, so
agreement would establish only that something had happened twice. We
note the corresponding limitation: because the synthetic bodies are
uniform, they bound the false-positive rate but do not measure recovery
of a known contrast, and the detection limit of the method is therefore
not established here.

%-----------------------------------------------------------------------
\subsection{A construction error and how it was found}
\label{sec:intro:construction}

One result of this work is negative and concerns the region construction
itself, and we state it before the contributions because it conditions
all of them.

An earlier implementation closed each region by fanning its boundary
edges to an interior point. That construction is watertight, has an
exact tetrahedral volume, and passes a volume-closure test---and it
encloses the wrong solid. Fanning to an interior point produces a radial
cone over the surface patch, not the plane-cut lobe intended. On a unit
sphere cut at $x=0.5$ the cone is a spherical sector of volume $1.0472$
and centroid $0.5625$, against $0.6545$ and $0.6750$ for the spherical
cap; on Itokawa the resulting error reversed the sign of the recovered
contrast. The volume-closure test then in use compared the fan mesh
against the same tetrahedra and could not detect the discrepancy. Only a
\emph{predicate} test---requiring every vertex of the region to satisfy
its defining inequality---distinguishes the two constructions.

The present implementation builds regions by exact plane clipping and is
verified in four independent geometric layers---against closed-form solids,
by complement closure, by the vertex predicate and by a manifold
criterion---with a fifth containment test where a shape model has more than
one component, an arithmetic audit of the mass balance on every solution,
and externally against the published lobe volume of 67P, where our head region holds $26.6\%$ of
the body against the $27\%$ given by \citet{jorda2016}. We record in
Section~\ref{sec:methods} that this external agreement is weaker than it
appears, because a single plane assigns the entire neck to one side
whereas the lobe fit of \citet{jorda2016} assigns $66\%$ and $27\%$ of
the volume to the two lobes and leaves the remaining $7\%$, the neck, in
neither; moving the plane across the neck changes our percentage by
$10.7$ points, and we report that sensitivity so that the reader may
judge it.

%-----------------------------------------------------------------------
\subsection{What is new here}
\label{sec:intro:new}

To be plain about the boundary: the relaxation argument and the
potential-variance technique are \citet{richardson2014}; the extension
to a two-region interior is \citet{kanamaru2019}; the polyhedron gravity
is \citet{werner1996}; and the finding that Itokawa's head is the denser
lobe is \citet{lowry2014} and \citet{kanamaru2019}, reproduced here but
not discovered here. \citet{scheeres2020} describe that finding as
tentative and not directly tested, which is part of what motivates an
independent implementation.

The contributions of this paper are the following.

\begin{enumerate}
\item An exact plane-clipping region construction with four verification
layers, of which the predicate and manifold tests can each fail in ways
the other three cannot, together with external validation against the
published lobe volume of 67P and an explicit account of how the neck is
apportioned (Section~\ref{sec:methods}).

\item The demonstration that the density recovered by the inversion is
resolution-stable over the meshes tested---a property distinct from, and
not implied by, the stability of the dispersion established in
Paper~IV---with no change across the three Itokawa meshes that share a
dividing plane and none across the three Arrokoth meshes at a fixed plane.
``No change'' means below the resolution of the scan, which bounds it at
$\SI{12.5}{\kilo\gram\per\cubic\metre}$, or $0.45\%$, on Itokawa; on
Arrokoth the recovered density is identical on all three meshes and the
complementary density moves by $\SI{0.4}{\kilo\gram\per\cubic\metre}$, or
$0.1\%$ (Sections~\ref{sec:itokawa} and~\ref{sec:nulls}).

\item The decomposition of the apparent stability of the inversion into
its algebraic and physical parts. Equation~\eqref{eq:excessidentity}
fixes a factor of $11.8$ between the coefficients of variation of $\Dmass$
and of $\densrest$, and the definition of $\Ddens$ as a difference inflates
its coefficient of variation over that of $\denshead$ by a further factor of
$2.8$; only the factor of $7.5$ by which $\Dmass$ is steadier than $\Ddens$
reflects compensation between contrast and region volume and is a property
of the objective function. The three combine exactly to the net factor of
$32$ between $\denshead$ and $\densrest$ over the eleven planes below the
grain-density ceiling (Section~\ref{sec:constrains}).

\item The identification of the background density, and with it the mass
anomaly to which \eqref{eq:excessidentity} ties it exactly, as the
quantity the method constrains in the strong-signal regime, together
with the demonstration by controlled synthetic experiment that this
holds in that regime rather than of the method as such
(Section~\ref{sec:constrains}).

\item The exclusion of the competing hypothesis that the conserved
quantity is the mass dipole $D=\Dmass\,(x_{R}-x_{cf})$. Dipole
conservation requires $\mathrm{d}\log\Dmass/\mathrm{d}\log\VR>0$ and
hence a scaling slope $s>-1$; quantitatively it requires $s=-0.756$ over
the eleven Itokawa planes below the grain-density ceiling, and the measured
$s=-1.141\pm0.018$ lies on the opposite side of $-1$, twenty-one standard
errors away (Sections~\ref{sec:itokawa} and~\ref{sec:constrains}).

\item The result that dispersion minimization does not constrain its own
dividing plane---at the outermost plane of the sweep the head density it
prefers stands $11.2$ standard deviations above the LL grain density---and
that an external physical bound is therefore required
(Section~\ref{sec:plane}).

\item Four weak-signal cases, establishing that the method does not
manufacture structure when the plane is correctly placed, and one worked
case showing what a misplaced plane costs (Section~\ref{sec:nulls}).

\item The finding that on a body without a measured mass the assumed
bulk density places the inferred contrast at the threshold of a sign
reversal, and the suggestion that sensitivity to that assumption may
serve as a secondary diagnostic (Sections~\ref{sec:nulls}
and~\ref{sec:constrains}).

\item An independent test of the Paper~IV resolution result against a
homogeneous ellipsoid, whose surface gravity has a closed form, removing
the need to treat the finest mesh as the truth
(Section~\ref{sec:discussion:ellipsoid}).
\end{enumerate}

Section~\ref{sec:methods} sets out the density model, the region
construction and its verification. Section~\ref{sec:itokawa} presents
the single strong signal and its stability properties.
Section~\ref{sec:nulls} presents the four weak-signal cases.
Section~\ref{sec:constrains} is the core of the paper: what the method
constrains, and the synthetic experiment that establishes in which
regime it does so. Section~\ref{sec:plane} shows that the dividing plane
is not self-determined. Section~\ref{sec:failures} catalogues the
failure modes. Section~\ref{sec:discussion} discusses the standing of
the relaxation premise on these bodies, Section~\ref{sec:guidelines}
sets out what the results imply for anyone applying the method, and
Section~\ref{sec:conclusions} states the conclusions.

%\end{document}

%=======================================================================
%  Paper V — Section 2 (Data and methods)   FINAL
%
%  Required in the master preamble (once only):
%
%  \usepackage{amsmath,amssymb,booktabs,graphicx,natbib,siunitx}
%  \sisetup{per-mode=symbol, range-phrase={--}, range-units=single,
%           separate-uncertainty=true}
%
%  \newcommand{\potvarn}{\tilde{\sigma}_{U}}
%  \newcommand{\Upot}{U}
%  \newcommand{\Ubar}{\bar{U}}
%  \newcommand{\Ufull}{U^{\mathrm{grav}}_{\mathrm{full}}}
%  \newcommand{\UR}{U^{\mathrm{grav}}_{R}}
%  \newcommand{\Urot}{U^{\mathrm{rot}}}
%  \newcommand{\denshead}{\varrho_{\mathrm{head}}}
%  \newcommand{\densrest}{\varrho_{\mathrm{rest}}}
%  \newcommand{\densR}{\varrho_{R}}
%  \newcommand{\densbulk}{\varrho_{\mathrm{bulk}}}
%  \newcommand{\densneck}{\varrho_{\mathrm{neck}}}
%  \newcommand{\densbody}{\varrho_{\mathrm{body}}}
%  \newcommand{\Ddens}{\Delta\varrho}
%  \newcommand{\Dmass}{\Delta M}
%  \newcommand{\VR}{V_{R}}
%  \newcommand{\Vtot}{V_{\mathrm{tot}}}
%  \newcommand{\Mtot}{M_{\mathrm{tot}}}
%  \newcommand{\ampfac}{\mathcal{A}}
%  \newcommand{\improve}{I}
%  \newcommand{\dCOM}{\Delta x_{\mathrm{COM}}}
%  \newcommand{\xsplit}{x_{\mathrm{split}}}
%  \newcommand{\CV}{\mathrm{CV}}
%  \newcommand{\Sstat}{\mathcal{S}}
%
%  Values superseded from earlier drafts of this section:
%      scan interval        [alpha,beta] -> 300-4000, widened to 300-8000
%      "parabola refinement"             -> removed; none is applied
%      curvature -> rho_rest  1.15 %     -> 1.38 %  (multiplier f/(1-f))
%      rho_rest interval     +/- 23      -> +/- 27 kg/m3
%      Itokawa plane in tab:planes 150 m -> 160.8 m (detector)
%      "Four checks"                     -> Five checks
%
%  All values confirmed. Spin-period provenance and the reproducibility
%  environment are now stated in the text; the only placeholder left in this
%  section is the Zenodo DOI, which set_doi.py fills.
%=======================================================================

\section{Data and methods}
\label{sec:methods}

\subsection{The density model, the mass constraint and the geopotential}
\label{sec:methods:model}

We divide a body into a region $R$ and its complement, assign a uniform
density to each, and hold the total mass at its measured value, as in
Equation~\eqref{eq:massbalance}. One free parameter remains. Writing
$\Ufull$ for the gravitational potential of the whole body at unit
density and $\UR$ for that of the region alone, linearity of the
Newtonian potential in density gives
%
\begin{equation}
U^{\mathrm{grav}}(\mathbf{r})
  = \densrest\,\Ufull(\mathbf{r})
  + \left(\densR - \densrest\right)\UR(\mathbf{r}) ,
\label{eq:superposition}
\end{equation}
%
so that a scan over $\densR$ requires only two evaluations of the
polyhedral gravity field, not one per trial density.

The quantity that relaxation acts on is not $U^{\mathrm{grav}}$ but the
geopotential in the co-rotating frame, which adds a centrifugal term,
%
\begin{equation}
\Upot(\mathbf{r})
  = \densrest\,\Ufull(\mathbf{r})
  + \left(\densR - \densrest\right)\UR(\mathbf{r})
  + \Urot(\mathbf{r}) ,
\qquad
\Urot(\mathbf{r}) = -\tfrac{1}{2}\,\omega^{2}\,d_{\perp}^{2}(\mathbf{r}) ,
\label{eq:geopotential}
\end{equation}
%
with $\omega = 2\pi/P$ and $d_{\perp}$ the perpendicular distance from
the spin axis. We adopt the convention in which the gravitational
potential is positive, so that $\Upot > 0$ over the surfaces considered
here and the normalization by $\lvert\Ubar\rvert$ in
Equation~\eqref{eq:dispersion} is unambiguous.

Equation~\eqref{eq:geopotential} carries the point on which the whole
construction turns, and we state it explicitly because it is easy to
lose. The first two terms scale linearly with density; the third does
not. Were $\Urot$ absent, multiplying every density by a constant would
multiply $\Upot$ by that constant and leave $\potvarn$, a normalized
quantity, unchanged---the objective would be flat in the overall density
scale and the minimization would be empty. It is the
density-independence of $\Urot$ that breaks the scaling and gives the
objective an interior minimum. The same argument applied at fixed total
mass is what gives the two-region problem a determinate solution.

Spin axes are taken from the published pole solutions cited in
Table~\ref{tab:shapes}, and the shape models are used in the body-fixed frame
in which they are supplied, with $z$ along the observed pole. That frame is not
exactly the principal frame even for a uniform body: on Itokawa the axis of
maximum moment of the uniform shape lies $\ang{0.77}$ from $z$. We note a consequence of
the two-region model that is not present in the uniform case: a non-zero
$\Ddens$ displaces the centre of mass, and in general rotates the
principal axes, so that the spin axis of the trial body is not exactly
the $z$-axis of the supplied mesh. On Itokawa the displacement is
$\SI{16.86}{\metre}$ along $x$ at the plane returned by the detector, and
ranges from $14.89$ to $\SI{20.52}{\metre}$ across the sweep of
Table~\ref{tab:itokawasweep}; the axis of maximum moment lies $\ang{0.77}$ from the mesh $z$-axis for the
uniform body and $\ang{1.87}$ for the recovered interior, an induced rotation
of $\ang{1.10}$. The recovered centre of mass computed from the same inertia
calculation lies at $x = \SI{17.22}{\metre}$, which is the $\SI{16.86}{\metre}$
displacement plus the $\SI{0.36}{\metre}$ offset of the centre of figure, so
the two calculations agree.
We hold the spin axis fixed at the shape-model $z$-axis, through the
origin, throughout. To measure what that choice costs we repeated the Itokawa
reference inversion with the rotational potential taken instead about the
recovered axis of maximum moment, through the recovered centre of mass, which
changes both the direction and the position of the axis. The recovered head
density moved from $2800$ to $\SI{2875}{\kilo\gram\per\cubic\metre}$, three
scan steps; the contrast from $1.507$ to $1.560$, an increase of $3.5\%$; and
the improvement from $16.9\%$ to $19.2\%$. The shift in $\denshead$ lies
inside the $\pm\SI{129}{\kilo\gram\per\cubic\metre}$ curvature interval of
Section~\ref{sec:itokawa}, and the change in contrast inside the $6.0\%$
error budget given there, so the choice does not alter the result; but it is
not negligible, and it enters that budget as a further term of comparable
size. We do not adopt the recomputed axis as the reference, because the
observed pole is what a real rotation defines and an interior whose principal
axis departs from it is, by that measure, less rather than more consistent
with the data.

Potentials are computed with the Werner--Scheeres formulation
\citep{werner1996,tsoulis2012} as implemented in
\texttt{polyhedral-gravity} version~\texttt{3.3.1}
\citep{gravitylib},
at field points placed at each facet centroid and displaced inward along
the outward normal by $10^{-4}$ of the bounding-box diagonal. The
displacement avoids the removable singularity of the Werner--Scheeres
expressions at a field point lying on a facet, but it is not physically
neutral: the same fractional offset is $\SI{6.5}{\centi\metre}$ on
Itokawa and $\SI{3.8}{\metre}$ on Eros, and it shifts the potential by
$\delta U \simeq g\,d$, of order $0.04\%$ of $\lvert\Ubar\rvert$ on both.
That is comparable to the dispersion improvements reported for the
weak-signal bodies in Section~\ref{sec:nulls}, so we repeated the Itokawa
reference run and the Eros run at Himeros---the strong case and the
best-resolved weak one---with the offset scaled by $10$ and by $0.1$. The
recovered contrasts did not move at the resolution of the density scan. The
improvements moved by at most $0.29$ percentage points on Itokawa and
$0.007$ on Eros, in both cases about $1.7\%$ of the improvement itself, and
in the same direction, a larger offset giving a larger improvement.

The normalized geopotential dispersion follows Paper~IV:
%
\begin{equation}
\potvarn = \frac{\sqrt{\sum_i A_i (\Upot_i - \Ubar)^2 \big/ \sum_i A_i}}
                {\lvert \Ubar \rvert},
\qquad
\Ubar = \frac{\sum_i A_i \Upot_i}{\sum_i A_i},
\label{eq:dispersion}
\end{equation}
%
with $A_i$ the facet areas. The sums run over the facets of the
\emph{whole body}, not of the region: the region enters only through
$\UR$ in Equation~\eqref{eq:geopotential}. Where facets are excluded from
the surface---on Arrokoth, those buried inside the opposite shell
(Section~\ref{sec:methods:shapes})---both $\Ubar$ and $\sum_i A_i$ are
recomputed over the retained set. Equation~\eqref{eq:dispersion} is the
same quantity that \citet{richardson2014} plot as the normalized
potential standard deviation. We report the improvement over the uniform
case,
%
\begin{equation}
\improve = \frac{\potvarn(\text{uniform}) - \potvarn(\densR^{\ast})}
                {\potvarn(\text{uniform})},
\label{eq:improvement}
\end{equation}
%
and the two derived quantities that will matter in
Section~\ref{sec:constrains}: the excess mass and the resulting
displacement of the centre of mass from the centre of figure,
%
\begin{equation}
\Dmass = (\densR - \densrest)\,\VR ,
\qquad
\dCOM = \frac{\Dmass\,(x_R - x_{\mathrm{cf}})}{\Mtot},
\label{eq:excessmass}
\end{equation}
%
where $x_R$ and $x_{\mathrm{cf}}$ are the $x$-coordinates of the
centroids of the region and of the whole body.
Equation~\eqref{eq:excessmass} is exact, not a small-anomaly
approximation: it follows from the definition of the centre of mass of a
two-density body with the complement's centroid eliminated through
Equation~\eqref{eq:massbalance}.

\subsection{Two consequences of the mass constraint}
\label{sec:methods:identity}

Two consequences of Equation~\eqref{eq:massbalance} should be stated
before any result rests on them. The first is the
identity~\eqref{eq:excessidentity}, an exact affine relation between
$\Dmass$ and $\densrest$ at fixed total mass and volume. The two are not
independent quantities: a statement that one is invariant under a change
of dividing plane is the same statement about the other. Differentiating,
%
\begin{equation}
\frac{\delta \Dmass}{\Dmass}
= -\,\ampfac\,\frac{\delta \densrest}{\densrest} ,
\label{eq:excessamplify}
\end{equation}
%
with $\ampfac$ as defined in Equation~\eqref{eq:amplification}. For the
Itokawa reference solution
$\densbulk = \SI{2012.87}{\kilo\gram\per\cubic\metre}$ and
$\densrest = \SI{1858.1}{\kilo\gram\per\cubic\metre}$ give
$\ampfac = 12.0$; across the thirteen planes of
Table~\ref{tab:itokawasweep} it ranges from $11.34$ to $12.70$ with a
mean of $11.77$, and over the nine admissible planes its mean is $11.90$.
Fractional ranges in $\Dmass$ and in $\densR$ are therefore not directly
comparable, and Section~\ref{sec:constrains} reports $\densrest$
alongside both. The identity can be seen operating in the output as well
as in the derivation: the ratio of the measured coefficients of
variation, $\CV(\Dmass)/\CV(\densrest) = 2.946/0.2506 = 11.76$ over all
thirteen planes and $3.00/0.252 = 11.90$ over the admissible nine,
reproduces the corresponding mean $\ampfac$ in both cases.

The second consequence is that $\ampfac$ is itself a poorly conditioned
quantity, because $\densbulk - \densrest$ is a difference of two
comparable numbers. Differentiating
Equation~\eqref{eq:amplification} at fixed $\densbulk$,
%
\begin{equation}
\frac{\delta \ampfac}{\ampfac}
 = \frac{\densbulk}{\densbulk - \densrest}\,
   \frac{\delta \densrest}{\densrest}
 \;\simeq\; 13\,\frac{\delta \densrest}{\densrest}
\quad\text{on Itokawa,}
\label{eq:ampsensitivity}
\end{equation}
%
so that a shift of $0.13\%$ in the background density moves $\ampfac$ by
$1.7\%$. This is why $\ampfac$ varies by $\pm5.6\%$ across a sweep over
which $\densrest$ varies by $\pm0.434\%$, and why any quoted value must
carry the plane at which it was evaluated.

Two region shapes are used. The \emph{head} model takes $R$ to be the
material beyond a plane, $x > \xsplit$. The \emph{neck} model takes $R$
to be a slab, $\lvert x - x_{\mathrm{neck}}\rvert < w/2$. The two are
alternative hypotheses about where the anomaly sits, not independent
constraints, and they overlap whenever $\xsplit$ falls inside the slab.
Section~\ref{sec:itokawa:ambiguity} shows that on Itokawa the two return
background densities differing by more than three times the combined
uncertainty on either, so the choice between them is a modelling decision
with consequences, not a matter of presentation.

\subsection{Region construction by exact plane clipping}
\label{sec:methods:clip}

Each region is built as a closed, consistently oriented triangle mesh.
Every facet crossing the plane is clipped against the half-space by the
reentrant-polygon algorithm of \citet{sutherland1974}; since the subject
polygon is a triangle and the clip region is a single half-space, the
clipped facet is a triangle or a quadrilateral and the general reentrant
machinery reduces to one pass. Intersection points are snapped exactly
onto the plane, and the opening is closed by fanning each boundary edge,
traversed in reverse, to an apex lying in the plane. Fan triangulation of
a planar loop gives the correct signed area from any coplanar apex, so
non-convex cross-sections need no special treatment. The slab of the neck
model is obtained by clipping twice.

Two details of the implementation are load-bearing. An intersection point
is created once per \emph{edge}, not once per facet, so that the two
facets sharing a cut edge reference the same vertex index; created per
facet, they would reference two distinct indices at the same position and
the clipped surface would be torn along the cut---a mesh that still
encloses the correct volume, since that is a sum of signed tetrahedra,
but is not a closed surface. And the cut boundary is traced into separate
loops, each capped to its own apex. A cross-section can be multiply
connected, or the region can fall into more than one piece, in which case
fanning every boundary edge to a single apex would stitch unrelated loops
together. Both conditions arise in this study: the Arrokoth shape model
is supplied as two unjoined shells, so a plane beyond its neck leaves the
region in two pieces (Section~\ref{sec:methods:shapes}).

\subsection{Why the obvious construction fails}
\label{sec:methods:cone}

An earlier version of this work closed each region differently, by
fanning every boundary edge to a single interior point $O$. That mesh is
watertight, and its volume agrees with the sum of the underlying
tetrahedra to machine precision. It nevertheless encloses the wrong
solid: the radial cone over the retained surface patch, with apex at
$O$---for a sphere, a spherical sector---rather than the plane-cut lobe
intended.

\begin{figure}[htbp]
\centering
\includegraphics[width=0.86\textwidth]{fig01_cone_vs_lobe.pdf}
\caption{The two solids, on a unit sphere cut at $x=0.5$. Top: the
plane-cut lobe intended by the density model, of volume $0.654498$ with
centroid at $0.675000$. Bottom: the radial cone---here a spherical
sector---produced by closing the region to an interior point $O$, of volume
$1.047198$ with centroid at $0.562500$. Both are watertight and both have
exact volumes; the sector has $1.600$ times the volume and its centroid lies
$0.1125R$ further in. Of the five checks of
Section~\ref{sec:methods:verify} only the vertex predicate distinguishes
them.}
\label{fig:cone}
\end{figure}

A unit sphere cut at $x=0.5$ makes the difference exact
(Figure~\ref{fig:cone}). The intended spherical cap has volume
$\pi h^{2}(3R-h)/3 = 0.654498$ and centroid
$3(2R-h)^{2}/[4(3R-h)] = 0.675000$; the spherical sector has volume
$\tfrac{1}{3}R\,A_{\rm cap} = 2\pi R^{2}h/3 = 1.047198$ and centroid
$\tfrac{3}{8}(2R-h) = 0.562500$. The constructions differ by a factor
$1.600$ in volume and by $0.1125\,R$ in centroid. On Itokawa the
consequence was not a small bias but a reversal: raising the density of a
sector adds mass along a wedge reaching back to the centre of the body
rather than in the head, and the recovered contrast came out inverted.

\subsection{Degenerate faces and numerical tolerance}
\label{sec:methods:tol}

Vertices lying within a tolerance of the cutting plane are classified as
\emph{on} it, snapped exactly onto it, and never duplicated as
intersection points; only strict sign changes generate an intersection.
Without this, a vertex sitting on the plane makes both of its edges
register as crossings, producing two coincident points and hence
zero-area triangles on both the surface and the cap. A single null-normal
facet makes the polyhedral evaluation return NaN at \emph{every} field
point, not merely near the plane, so the failure is total rather than
local. The tolerance is $10^{-10}$ of the bounding-box diagonal.

At that tolerance a vertex a few millimetres from the plane on a
kilometre-scale body still produces a sliver, which the gravity library
flags as numerically unstable. We tested whether such slivers matter by
repeating runs with the tolerance raised by four orders of magnitude,
which removes them. On 67P the minimum facet area rose from
$2.381\times10^{-12}$ to $3.660\times10^{-8}$~km$^{2}$ while the
recovered contrast, the improvement and the excess mass changed in the
seventh significant figure or beyond; on Itokawa, the body from which the
principal result of this paper comes, the recovered $\denshead$ was
unchanged at every plane of Table~\ref{tab:itokawasweep} and $\densrest$
was unchanged to four decimal places;
on a synthetic body every printed digit was unchanged. The slivers are
harmless, as the physics requires---their contribution to
Equation~\eqref{eq:dispersion} scales with their area---but this was
established by measurement rather than assumed.

\subsection{Verification}
\label{sec:methods:verify}

Five checks guard the calculation. Three of them can fail in ways the
others cannot, and that is the reason for having five.

\emph{Analytic.} The reference solid is an icosphere of $20{,}480$ facets
inscribed in the unit sphere.
The clipped cap has volume $0.653931$ against the closed form $0.654498$
($-0.087\%$) and centroid $0.674936$ against $0.675000$ ($-0.009\%$); a
slab between $x=\pm0.2$ has volume $1.239437$ against $1.239882$
($-0.036\%$). These residuals are the polyhedral approximation of the
sphere itself---the whole-body volume of the same mesh falls short of
$4\pi/3$ by $0.05\%$, and an inscribed polyhedron necessarily
underestimates---not error in the clipping, which is exact for the given
polyhedron. That the cap and slab deficits bracket the whole-body
deficit, and that all three carry the same sign, is what identifies them
as discretization rather than clipping error.

\emph{Complement.} $V(x>c) + V(x<c)$ must equal the body volume. The
residual on Itokawa is $2.0\times10^{-14}$ per cent, that is
$2.0\times10^{-16}$ in relative terms. That is below the double-precision
epsilon of $2.2\times10^{-16}$, and below the $\sqrt{N}\,\epsilon \approx
7\times10^{-14}$ that an uncorrelated sum of $10^{5}$ signed tetrahedra would
give. The explanation is that the two half-space meshes share most of their
facets with the whole-body mesh, so the round-off is common to both sides of
the comparison and cancels rather than accumulating. The test is therefore
weaker than its magnitude suggests: it confirms that the clipping partitions
the tetrahedral sum, which is why Section~\ref{sec:methods:cone} adds the
predicate test.

\emph{Predicate.} Every vertex of the region must satisfy its defining
inequality. The plane-cut lobe passes with zero violation. The spherical
sector of Section~\ref{sec:methods:cone} fails outright, its apex lying
at the centre of the body, with a violation of $0.5$ in units where the
plane sits at $x=0.5$.

\emph{Manifold.} Every directed edge must appear exactly once and its
reverse must be present, and the Euler characteristic
$\chi = n_V - n_E + n_F$ must equal $2$ for each connected component,
where $n_V$, $n_E$ and $n_F$ count vertices, edges and facets. This is
independent of the other checks, and necessarily so: the
divergence-theorem volume is a sum of signed tetrahedra and requires only
that the oriented facets tile the boundary with correct signs, not that
they are stitched together. A surface torn along the cut therefore passes
the analytic, complement and predicate tests while not being a closed
surface at all. The criterion is stated per component rather than as a
single $\chi = 2$, because a region cut from a body supplied as several
unjoined shells is legitimately in several pieces.

\emph{Mass balance.} Every solution must satisfy
$\Dmass = (\densbulk - \densrest)\Vtot$ and, independently,
$\Dmass = \Ddens\,\VR$. Both hold to the printed precision on all
thirteen rows of Table~\ref{tab:itokawasweep}. The check is enforced by
construction and cannot fail unless the region volume passed to the mass
constraint differs from the one passed to the gravity evaluation---which
is exactly the failure mode of Section~\ref{sec:methods:cone}, where the
two differed by a factor $1.600$. It is therefore the layer that would
have caught that error at the level of the output rather than of the
geometry, and it is the reason we tabulate $\densrest$ and $\Dmass$ side
by side rather than quoting either alone.

The first three checks run on every region and abort the model if they
fail. The manifold check reports rather than aborts, and distinguishes a
defect inherited from the source mesh from one introduced by the
clipping; the distinction matters because a volume and a potential are
unaffected by tearing, so a defect that the source already carries cannot
be repaired here and should not stop the calculation. The mass-balance
audit runs on every solution and is reported with the results.
Section~\ref{sec:methods:shapes} records which shape models carry
topological defects.

\subsection{External validation of the geometry}
\label{sec:methods:external}

The preceding subsections test the construction against solids of our own
choosing. A stronger check is available for 67P, whose small lobe has
been measured independently. Our clipped region holds $26.6\%$ of the
body volume against the $27\%$ of \citet{jorda2016}, at the plane the
detector of Section~\ref{sec:methods:neck} returns with nothing tuned.
The 67P model is a single connected surface, so producing that figure
requires the clipping to do real geometric work.

Three qualifications belong with that number, and we state them rather
than let the agreement stand unexamined. First, \citet{jorda2016}
decompose the nucleus into two lobes \emph{and} a neck, the neck
accounting for about $7\%$ of the total volume and assigned to neither
lobe, whereas a single plane necessarily assigns the entire neck to one
side. Our figure therefore agrees with theirs only because the detected
plane falls at the far edge of the neck; displacing it across the neck
moves the fraction by $10.7$ percentage points, from $31.80\%$ at
$\xsplit = \SI{0.70}{\kilo\metre}$ to $21.14\%$ at $\SI{1.20}{\kilo\metre}$,
a gradient of $-21.2$ points per kilometre. Sweeping the $7\%$ that
\citet{jorda2016} assign to the neck therefore corresponds to about
$\SI{330}{\metre}$ of plane travel. We report that sensitivity in
Table~\ref{tab:planes} so that the reader may judge how much of the
agreement is geometric fidelity and how much is plane placement.
Second, the comparison is between a volume fraction and a volume
fraction; the mass fractions quoted by \citet{jorda2016} coincide with
them only under the assumption of uniform density, which is the
assumption this paper exists to test. Third, our $\Vtot$ of
$\SI{18.5912}{\cubic\kilo\metre}$ comes from an SPG model, whereas the
lobe decomposition of \citet{jorda2016} is based on the SPC SHAP5 model
with $\Vtot = \SI{18.8+-0.3}{\cubic\kilo\metre}$; the $1.1\%$ difference
in total volume is smaller than the effect just described but is not
zero.

The corresponding comparison for Arrokoth is weaker, and we report it as
such. Our region holds $66.33\%$ against $66.30\%$ from
\citet{porter2024}, but the model is supplied as two unjoined shells and
the larger of them is $66.30\%$ of the summed component volumes by itself
(Section~\ref{sec:methods:shapes}). The agreement is therefore largely a
property of the model as supplied rather than an independent test of the
clipping, and only the 67P comparison should be read as validation.

\subsection{Locating the dividing plane}
\label{sec:methods:neck}

The cross-sectional half-width is binned along $x$ and a saddle sought
within it. A neck is a \emph{local} minimum flanked by a hump on either
side, not merely the smallest half-width in the interior of the profile.
The distinction is not academic: on a body that tapers at both ends the
smallest interior half-width lies on a flank. On Eros the
smallest-interior-minimum rule would select
$x = \SI{-10.76}{\kilo\metre}$, on a shoulder, whereas the saddle rule
adopted here selects Himeros at $x = \SI{+4.47}{\kilo\metre}$.
% Confirmed against the run records: -10.7619 km is returned by the rejected
% smallest-interior-minimum rule, +4.4698 km by the saddle rule used here.
Section~\ref{sec:nulls} shows what the substitution costs.

Candidates are therefore detected on runs of equal profile value---a
strict point test finds nothing on a symmetric profile, where two samples
straddling the minimum carry equal values---and each is scored by
%
\begin{equation}
p = \frac{\min(w_{\mathrm{left}}, w_{\mathrm{right}}) - w_{\mathrm{neck}}}
         {\max_x w(x)},
\label{eq:prominence}
\end{equation}
%
with $w_{\mathrm{left}}$ and $w_{\mathrm{right}}$ the flanking maxima.
The most prominent candidate is selected provided $p \ge 0.02$;
otherwise the routine returns no neck and the plane must be supplied
explicitly. Where the selected saddle differs from the absolute interior
minimum, the rejected position is reported.

Measured prominences separate the bodies cleanly: Arrokoth $0.490$, Eros
$0.216$, Itokawa $0.170$, 67P $0.091$, and Bennu $0.0004$, which is
refused. The gap between the lowest accepted value and the highest
rejected one spans more than two orders of magnitude, so the threshold is
not a fine judgement. Figure~\ref{fig:profiles} shows the three cases the
detector must distinguish.

\begin{figure}[htbp]
\centering
\includegraphics[width=\textwidth]{fig02_profiles.pdf}
\caption{Cross-sectional half-width profiles for the three cases the saddle
detector must distinguish. Itokawa (top): a clear saddle that is also the
smallest interior half-width, accepted at a prominence of $0.170$, one of
three local minima. Eros (middle): the true saddle at Himeros, prominence
$0.216$, while the smallest interior half-width lies on a flank at
$x=-10.76$~km and is rejected---taking it instead costs what
Section~\ref{sec:nulls} reports. Bennu (bottom): no saddle at all, the
largest of three local minima having a prominence of $0.0004$, and the
detector declines to return a plane.}
\label{fig:profiles}
\end{figure}

Bennu is retained in the study precisely because
the detector refuses it: it supplies a case in which the two-region model
is applied to a body with no neck at all, and Section~\ref{sec:nulls}
reports the inversion at a plane through the centre of figure, chosen by
us rather than by the detector, to show what the objective returns when
the model has no geometric warrant.
% Confirmed: the Bennu plane is x = 0 exactly, imposed by us, enclosing
% 50.24 per cent of the volume; the detector returns no plane on this body.

Table~\ref{tab:planes} collects the plane used for every body, together
with its provenance, so that no reported result depends on a plane whose
origin is stated only in passing. The two Itokawa entries are listed
separately and should not be conflated: the detector returns
$\xsplit = \SI{160.8}{\metre}$, enclosing $16.43\%$ of the volume, while
the plane of \citet{kanamaru2019} lies at $\SI{150}{\metre}$ and encloses
$17.89\%$. The first is the reference solution of this paper; the second
is used only where a direct comparison with published values is intended,
and every quoted number states which of the two it belongs to.

\begin{table}[htbp]
\centering
\small
\caption{Dividing planes used in this work. ``Detector'' means the saddle
returned by Equation~\eqref{eq:prominence} with no tuning; ``imposed''
means supplied by us, with the reason given. $f = \VR/\Vtot$ is the
volume fraction of the region. The two Itokawa rows are distinct planes
separated by $\SI{10.8}{\metre}$ and are kept apart throughout the paper.}
\label{tab:planes}
\setlength{\tabcolsep}{4pt}
\begin{tabular}{llrrrl}
\toprule
Body & Model & $\xsplit$ & $p$ & $f$ (\%) & Provenance \\
\midrule
Itokawa  & head & \SI{160.8}{\metre}      & $0.170$  & $16.43$ & detector; reference solution \\
Itokawa  & head & \SI{150}{\metre}        & ---      & $17.89$ & imposed; \citet{kanamaru2019}, comparison only \\
Itokawa  & neck & \SI{118.7}{}--\SI{202.9}{\metre} & ---   & $12.05$ & imposed; slab, Section~\ref{sec:itokawa:ambiguity} \\
Eros     & head & \SI{+4.47}{\kilo\metre} & $0.216$  & $34.6$  & detector (Himeros) \\
Bennu    & head & \SI{0}{\kilo\metre}    & $0.0004$ & $50.2$  & imposed; detector refuses \\
67P      & head & \SI{+0.952}{\kilo\metre} & $0.091$  & $26.6^{a}$  & detector \\
Arrokoth & head & \SI{-4.786}{\kilo\metre} & $0.490$  & $66.33$ & detector \\
\bottomrule
\end{tabular}
\end{table}
\par\smallskip{\footnotesize
\textsuperscript{a}\,Across the 67P neck $\mathrm{d}f/\mathrm{d}\xsplit =
-21.2$ percentage points per kilometre; $f$ runs from $31.80\%$ at
$\SI{0.70}{\kilo\metre}$ to $21.14\%$ at $\SI{1.20}{\kilo\metre}$. The
Itokawa neck slab has width $w = \SI{84.2}{\metre}$.}

\subsection{The minimization and its sensitivity interval}
\label{sec:methods:minimisation}

The single free parameter of Equation~\eqref{eq:massbalance} is scanned
rather than solved for, because the objective is cheap once $\Ufull$ and
$\UR$ are in hand: by Equation~\eqref{eq:geopotential} each trial density
costs one vector addition. We scan $\densR$ on a uniform grid of step
$\SI{25}{\kilo\gram\per\cubic\metre}$ over an interval set for each body
relative to its own bulk density, since a range appropriate to a rocky
asteroid would miss the minimum entirely on a cometary nucleus or a
Kuiper Belt object, and report the grid point at which $\potvarn$ is
least. On Itokawa the interval is $300$--$4000$~kg\,m$^{-3}$, widened to
$300$--$8000$~kg\,m$^{-3}$ for the outer planes of the sweep of
Section~\ref{sec:itokawa:sweep}; because the two intervals share the same
grid, widening adds candidate points without displacing existing ones.
The intervals used elsewhere are
$500$--$5000$ on Eros, $200$--$3000$ on Bennu, $100$--$4000$ on both 67P
models, and $25$--$3000$~kg\,m$^{-3}$ on Arrokoth at each of the three assumed
bulk densities, with grid steps of $25$, $10$, $5$ and $5$~kg\,m$^{-3}$
respectively.

A minimum returned at either end of the scan interval is flagged as a
boundary minimum and the interval widened until an interior minimum is
obtained; no boundary minimum is reported as a result. Every solution in
this paper satisfies that condition, and Section~\ref{sec:failures}
describes the one case in which the diagnostic fired during development.

The upper limit of the scan also bounds the range of dividing planes the
sweep can probe. Holding $\Dmass$ at its Itokawa sweep value of
$2.80\times10^{9}$~kg, a region small enough to require $\densR$ at the cap
of $\SI{8000}{\kilo\gram\per\cubic\metre}$ has a volume fraction of
$\Dmass/[(8000-\densrest)\Vtot] = 2.56\%$, which the sweep reaches at
$\xsplit \approx \SI{0.257}{\kilo\metre}$. No plane in the sweep comes
within $\SI{17}{\metre}$ of that, but the figure matters in
Section~\ref{sec:plane}, where an extrapolation beyond the sweep would land
close to it.

No sub-grid refinement is applied, and the reported densities are
therefore quantized at that step. We state this rather than leave it to
be inferred, because it is visible in the results: in
Table~\ref{tab:itokawasweep} the planes at $\SI{160.0}{\metre}$ and
$\SI{160.8}{\metre}$, separated by $\SI{0.8}{\metre}$, return the same
$\denshead$ but different $\densrest$, and at two places $\dCOM$
decreases as the plane advances even though $x_R - x_{\mathrm{cf}}$
increases monotonically throughout. The size of the effect follows from
the mass constraint. A half-step error of
$\SI{12.5}{\kilo\gram\per\cubic\metre}$ in $\denshead$ propagates as
$[f/(1-f)]\times\SI{12.5}{\kilo\gram\per\cubic\metre}
= \SI{2.4}{\kilo\gram\per\cubic\metre}$, or $0.13\%$, against an observed
half-range in $\densrest$ of $\SI{8.05}{\kilo\gram\per\cubic\metre}$
($0.434\%$); in standard-deviation terms, a uniformly distributed
quantization error contributes
$\SI{25}{\kilo\gram\per\cubic\metre}/\!\sqrt{12}$ in $\denshead$ and
hence $\SI{1.4}{\kilo\gram\per\cubic\metre}$ in $\densrest$, against an
observed $\SI{4.65}{\kilo\gram\per\cubic\metre}$. Roughly one tenth of
the observed variance in $\densrest$ is therefore numerical rather than
physical, and the quoted plane-placement scatter is an upper bound on the
physical effect.
% NOTE: Equation (3) is affine in rho_R once the mass constraint is
% substituted, so sigma_U^2 is a ratio of quadratics in rho_R and the
% stationarity condition is linear, with a single root. Solving it in
% closed form would remove the grid step from the budget entirely and
% would also give the curvature interval as the two roots of a quadratic
% rather than as a fitted quantity. This is left for a future revision.

The objective is deterministic: there is no noise model, and repeated
evaluation returns identical values. The interval we quote alongside
$\densR^{\ast}$ is therefore \emph{not} a confidence interval, and we do
not call it one. It is a \emph{curvature sensitivity interval}: the range
of $\densR$ over which $\potvarn$ rises by no more than $0.60\%$ above
its minimum, the threshold being the plateau width between the two finest
meshes as reported in Paper~IV and reproduced in
Appendix~\ref{app:envelope}. Recomputing that plateau for the present study
gives a tighter figure: between the two finest meshes the dispersion moves by
$0.008\%$ on Itokawa and at most $0.098\%$ on any of the four resolution
bodies, with either decimator. We retain $0.60\%$ because it is the published
value and because a wider threshold gives a wider interval, so the sensitivity
intervals quoted throughout are conservative by roughly a factor of six.
It answers the question ``how far can the density move before the
objective would have preferred a different mesh?'' rather than ``what is
the standard error of the estimate?'', and on Itokawa it spans
$\pm\SI{130}{\kilo\gram\per\cubic\metre}$.

Because the mass constraint ties the two densities, that interval
propagates to the background density as
$\delta\densrest = [f/(1-f)]\,\delta\denshead$, which at $f = 0.1643$
gives $\SI{25.6}{\kilo\gram\per\cubic\metre}$, or $1.38\%$. This is three
times the plane-placement half-range of $0.434\%$ and is the dominant
term in the budget; any statement of the form ``the background density is
stable to four parts in a thousand'' is a statement about plane placement
alone and must not be read as an uncertainty. Combining the two in
quadrature at the reference plane,
%
\begin{multline}
\densrest = 1858 \pm \SI{27}{\kilo\gram\per\cubic\metre},
\qquad
\Dmass = (2.75 \pm 0.48)\times10^{9}~\mathrm{kg}
       = 7.7 \pm 1.3\%~\text{of}~\Mtot ,
\qquad \\
\dCOM = 16.9 \pm \SI{2.9}{\metre} .
\label{eq:budget}
\end{multline}
%
The background density differs from the bulk density by
%
\begin{equation}
\Sstat \;\equiv\;
\frac{\lvert \densbulk - \densrest \rvert}{\sigma(\densrest)}
\;=\; 5.8 ,
\label{eq:sstat}
\end{equation}
%
times its own combined uncertainty. That ratio, and not the scatter of
$\densrest$ alone, is what distinguishes a detection from a uniform body:
for a uniform body $\densrest \equiv \densbulk$ holds identically and the
scatter would be zero, so stability alone is evidence of nothing.
$\Sstat$ is used in Section~\ref{sec:itokawa:ambiguity} to compare region
shapes and in Section~\ref{sec:nulls} to compare bodies.

\subsection{Shape models, masses and spin states}
\label{sec:methods:shapes}

\begin{table}[htbp]
\centering
\scriptsize
\setlength{\tabcolsep}{4pt}
\caption{Shape models and parameters. Facets are those used after
decimation, where any was applied. Arrokoth has no measured mass; three
assumed bulk densities were run. The $\chi$ column is the Euler
characteristic of the mesh as used. $\densbulk = \Mtot/\Vtot$ is listed
beside the published value because it enters
Equation~\eqref{eq:amplification} with the amplification of
Equation~\eqref{eq:ampsensitivity}.}
\label{tab:shapes}
\small
\setlength{\tabcolsep}{4.5pt}
\begin{tabular}{lrrrrrrrr}
\toprule
Body & Native & Used & $\chi$ & $\Vtot$ & $\Mtot$ (kg)
     & $\densbulk$ & Published $\densbulk$ & $P$ (h) \\
\midrule
Itokawa  & $49{,}152$      & $49{,}152$ & $2$ & $0.0177856$~km$^3$ & $3.58\times10^{10}$   & $2012.9$ & $1950\pm140$ & $12.1324$ \\
Eros     & $49{,}152$      & $49{,}152$ & $2$ & $2506.39$~km$^3$   & $6.687\times10^{15}$  & $2668$   & $2670\pm30$  & $5.27$ \\
Bennu    & $196{,}608$     & $50{,}000$ & $2$ & $0.061553$~km$^3$  & $7.330\times10^{10}$  & $1191$   & $1190\pm13$  & $4.296061$ \\
67P      & $1{,}309{,}996$ & $50{,}000$ & $2$ & $18.5912$~km$^3$   & $9.982\times10^{12}$  & $537$    & $537.8\pm0.7$ & $12.4043$ \\
Arrokoth & $40{,}960$      & $40{,}960$ & $4$ & $4123.81$~km$^3$   & assumed               & ---      & $155$--$600$ & $15.92$ \\
\bottomrule
\end{tabular}
\end{table}
% Published densities: Itokawa 1950+/-140 (Abe et al. 2006); Eros 2670+/-30
% (Yeomans et al. 2000); Bennu 1190+/-13 (Scheeres et al. 2019); 67P
% 537.8+/-0.7, 3-sigma, from the SHAP7 volume of 18.56 km^3 (Preusker et al.
% 2017), against 532+/-7 from SHAP5 (Jorda et al. 2016); Arrokoth 235, 1-sigma
% range 155-600 (Keane et al. 2022). Arrokoth period 15.92+/-0.02 h (Spencer
% et al. 2020); 67P initial period 12.4041+/-0.0001 h (Jorda et al. 2016).
Spin periods and poles are taken from the following. Itokawa,
$\SI{12.1324(1)}{\hour}$ from the ground-based lightcurve campaign of
\citet{nishihara2019}, which refines the $\SI{12.132(5)}{\hour}$ of
\citet{kaasalainen2003}; Eros, $\SI{5.27025547}{\hour}$ with the pole of
\citet{miller2002}, of which the $\SI{5.27}{\hour}$ used here is a rounding
that changes the rotational term by $8\times10^{-5}$ in relative terms;
Bennu, $\SI{4.296061}{\hour}$, the value determined in flight
\citep{hergenrother2019}, against $\SI{4.297461}{\hour}$ from the
pre-encounter radar and lightcurve solution \citep{nolan2013}---the nucleus
is spinning up, and the two differ by five seconds, which changes the
rotational term by $7\times10^{-4}$; 67P and Arrokoth as given above. The
shape models themselves are identified by their published references rather
than by archive identifiers, which the deposit records
(Section~\ref{sec:conclusions}, data availability): \citet{gaskell2008} for
Itokawa and Eros, the OLA v20 product of \citet{barnouin2019} for Bennu,
SHAP2 and SPC~2017 for 67P \citep{preusker2017,jorda2016}, and
\citet{porter2024} for Arrokoth.
% Confirmed: the Arrokoth V_tot of 4123.81 km^3 is the sum of the two
% components, 1389.53 and 2734.28 km^3; they do not overlap, the intersection
% being 0.044 per cent of the total (Section 2.7).

Masses are taken from \citet{abe2006} for Itokawa, \citet{yeomans2000}
for Eros, \citet{scheeres2019} for Bennu, whose $7.329\pm0.009\times10^{10}$~kg
we round to $7.330$, and \citet{patzold2016} for 67P. The published 67P
density of $537.8\pm0.7$~kg\,m$^{-3}$ is the SHAP7 value of
\citet{preusker2017}; the SHAP5 volume of \citet{jorda2016} gives
$532\pm7$, and our $537$ from a decimated model of $18.5912$~km$^{3}$ lies
between them. The 67P rotation period is the initial value, which held from
arrival until October 2014; \citet{jorda2016} give $12.4041\pm0.0001$~h for
it, $0.7$~s from the value used here. The period then lengthened to
$12.4304$~h by May 2015 and fell to $12.305$~h just before the August 2015
perihelion, so the epoch matters: the change is several hundred times the
quoted uncertainty.
Arrokoth has no mass determination \citep{spencer2020}, and we run three
assumed bulk densities, $250$, $400$ and
$\SI{500}{\kilo\gram\per\cubic\metre}$, with a fourth at $235$, spanning the range discussed by
\citet{mckinnon2020} and \citet{spencer2020}. The published constraints do not
agree with one another. \citet{keane2022} infer $235$~kg\,m$^{-3}$ with a
$1\sigma$ range of $155$--$600$; \citet{spencer2020} require more than
$290$ if the neck is not in tension; and \citet{mckinnon2020} find
$250$--$500$ for a comet-like strength. Our three principal values span the last
of these, and we add a fourth run at the central value of \citet{keane2022}.
At $235$~kg\,m$^{-3}$ the recovered contrast is $1.066$, against $0.9999$ at
$250$: the sense of the answer reverses between the two, inside the published
$1\sigma$ range and just below the bound of \citet{spencer2020}. A linear
extrapolation from the three higher runs would have placed the contrast at
$235$ near $1.006$; the measured value is ten times further from unity, so the
dependence is not linear in this range and the crossing is not predictable
from runs on one side of it. On this body the assumed density, not the data,
sets the direction of the answer (Section~\ref{sec:nulls}).

Four points on the models deserve statement.

First, our Itokawa model encloses $0.0177856$~km$^3$ against the
$(1.84\pm0.092)\times10^{7}$~m$^3$ adopted by \citet{abe2006}, a
difference of $-0.67\sigma$; the bulk density that follows,
$\SI{2012.9}{\kilo\gram\per\cubic\metre}$, is $+0.45\sigma$ from their
$1950\pm140$. The difference is ordinary and within the published
uncertainty, but it is not negligible for our purposes, because
Equation~\eqref{eq:ampsensitivity} amplifies it: a $3.2\%$ offset in
$\densbulk$ moves $\ampfac$ by a factor of order unity, and $\ampfac$ is
the number by which the headline stability ratio must be divided before
any part of it can be called physical. Section~\ref{sec:itokawa:scale}
shows that the contrast and the excess mass move by under one per cent
when the model is rescaled to the published volume, and reports $\ampfac$
for both.

Second, the Arrokoth model of \citet{porter2024} is supplied as two
unjoined shells: $\chi = 4$ over two connected components, with no
repeated or unmatched edges. Geometrically the two lobes touch;
topologically they are separate surfaces. This is why the manifold
criterion of Section~\ref{sec:methods:verify} is stated per component,
and why a plane placed beyond the neck yields a region legitimately in
two pieces.

Two consequences of that arrangement were measured rather than assumed.
The shells touch but barely interpenetrate: a containment test on
$2\times10^{5}$ uniformly distributed samples places $0.044\%$ of the
volume inside both shells. With $88$ samples falling in the intersection
the Poisson standard error is $0.005\%$, so the overlap is resolved
rather than consistent with zero; it is nevertheless small enough that
the double-counted mass is $4\times10^{-4}$ of the total and affects no
reported digit.
% NOTE: an exact boolean mesh intersection would remove the need for this
% qualification and is preferable if the geometry kernel allows it.
Separately, $455$ facets, carrying $0.943\%$ of the surface area, have
their centroids inside the other shell. Those facets are not on the true
surface and have no business in an average taken over it. All Arrokoth
results reported here exclude them, with $\Ubar$ and $\sum_i A_i$
recomputed over the retained set; Section~\ref{sec:nulls:arrokoth} gives
the size of the effect, which is negligible at most planes and decisive
at one.

Third, three shape models of 67P were available. The two that reproduce
their own published volumes to better than $1\%$ agree closely with each
other, while a third, a decimated product, encloses $3.9\%$ less volume
and is not used. The SPC~2017 model is not a closed surface as supplied,
carrying one repeated and two unmatched directed edges; the defect is
inherited by any region cut from it, and is reported where that model is
used.

Fourth, meshes above $50{,}000$ facets were reduced by quadric edge
collapse with the link condition enforced, which preserves the manifold
property and returns the requested facet count exactly. Papers~III
and~IV used vertex clustering, which is equally accurate in the fields it
produces---tested against a homogeneous ellipsoid in
Section~\ref{sec:discussion:ellipsoid}---but does not preserve topology.
The change of method alters the numbers where the reduction is severe: at
a $16\%$ reduction on Itokawa the recovered contrast moves by $0.005\%$,
at $68\%$ on Bennu by $0.054\%$, and at $96\%$ on 67P by $-3.9\%$. The
last figure is larger than the dispersion improvement reported for 67P in
Section~\ref{sec:nulls}, and we carry it into the error budget rather
than record it in passing: on that body the decimation, not the
inversion, dominates.
% Confirmed independent: the -3.875 % is the change in recovered contrast
% between vertex clustering and quadric collapse at 96 per cent reduction on
% the model actually used; the 3.9 % is the volume deficit of a different,
% discarded model. The agreement to two figures is coincidental.

\subsection{Reproducibility}
\label{sec:methods:repro}

The self-test of Section~\ref{sec:methods:verify}, including the
mass-balance audit over all thirteen Itokawa planes, is released as a
runnable script, and the identical values were reproduced on a second
machine: the results reported here were computed under Windows~10
(build~19045) with Python~3.13 and \texttt{polyhedral-gravity}~3.3.1, and
the geometry self-test reproduces its printed values digit for digit under
Ubuntu~24.04 with Python~3.12.3. All code, the driver
scripts and the derived tables are archived at
\href{https://doi.org/10.5281/zenodo.23003040}{doi:10.5281/zenodo.23003040};
the shape
models are not redistributed and are identified by their published
references in Section~\ref{sec:methods:shapes}, with the archive holding
each one recorded in the deposit. The scan intervals of
Section~\ref{sec:methods:minimisation} are recorded per plane in the
released configuration, so that the widening described there can be
audited rather than taken on trust.
% DOI reserved 2026-09-28: 10.5281/zenodo.23003040. Confirm before
% submission that the published record is the one that produced these numbers.
%=======================================================================
%  Paper V — Section 3 (Itokawa: the one strong inversion signal) FINAL
%
%  Requires the macros listed at the head of Section 2.
%
%  Values superseded from earlier drafts of this section:
%      scan interval 300-3900        -> 300-4000, widened to 300-8000
%      I = 16.888 %                  -> 16.9 %
%      contrast 1.507                -> 1.51 +/- 0.09
%      Dx_COM 16.86 m                -> 16.9 +/- 2.9 m
%      "identical across meshes"     -> bounded at half a scan step
%      compositional bounds          -> now carry their own uncertainties
%      neck model reported as a fit  -> reported as a non-detection
%
%  All values confirmed against paperV_final.json and paper4_recheck.json.
%  Further corrections in this pass:
%      curvature bounds 130 -> 129, neck 261 -> 259 kg/m3
%      neck slab fraction 12.01 -> 12.05 %
%      dipole test moved to the eleven planes below the grain-density
%      ceiling, to match Section 5: g = -0.244, s = -1.141 +/- 0.018, 21 sigma;
%      eleven-plane block added to the sweep table
%=======================================================================

\section{Itokawa: the one strong inversion signal}
\label{sec:itokawa}

Of the five bodies examined, Itokawa is the only one on which the
two-region inversion produces a dispersion improvement that cannot be
accounted for by mesh resolution alone. This section reports that result
and then tries to break it: by changing the mesh, by moving the dividing
plane, by changing the region shape, by changing the assumed mass, and by
rescaling the shape model. The result survives in the sense that the
recovered contrast is always positive and of the same order; it does not
survive in the sense of being independent of the choices made, and
Section~\ref{sec:constrains} takes up what that means.

\subsection{Result at the detected neck}
\label{sec:itokawa:reference}

The neck detector of Section~\ref{sec:methods:neck} returns
$\xsplit = \SI{0.1608}{\kilo\metre}$ with a prominence of $0.170$,
placing $16.43\%$ of the volume in the head, whose centroid lies at
$x_R = \SI{+0.2197}{\kilo\metre}$. This is the reference solution of the
paper. It is not the plane used by \citet{kanamaru2019}, which lies at
$\SI{0.150}{\kilo\metre}$ and encloses $17.89\%$; the two are kept apart
throughout, and Section~\ref{sec:itokawa:comparison} quotes the second
where a comparison with published values is intended.

Scanning the head density over $4000$~kg\,m$^{-3}$ in steps of
$\SI{25}{\kilo\gram\per\cubic\metre}$ and imposing the mass constraint at
each step gives
%
\begin{equation}
\denshead = 2800 \pm \SI{129}{\kilo\gram\per\cubic\metre},
\qquad
\densrest = 1858 \pm \SI{27}{\kilo\gram\per\cubic\metre},
\qquad
C = \denshead/\densrest = 1.51 \pm 0.09 ,
\label{eq:itokawaref}
\end{equation}
%
with a dispersion improvement of $\improve = 16.9\%$ over the uniform
body. The interval on $\denshead$ is the curvature sensitivity interval
defined in Section~\ref{sec:methods:minimisation}, not a confidence
interval; that on $\densrest$ combines it with the plane-placement
scatter of Table~\ref{tab:itokawasweep}; and that on $C$ follows from the
first two, which move in opposite directions because the mass constraint
ties them, so that their fractional contributions add rather than combine
in quadrature: $129/2800 + 25.4/1858 = 4.6\% + 1.4\% = 6.0\%$, the
$\SI{25.4}{\kilo\gram\per\cubic\metre}$ being the curvature interval
propagated through the mass constraint, $[f/(1-f)]\times129$ at
$f = 0.1643$. The $\pm\SI{27}{\kilo\gram\per\cubic\metre}$ quoted for
$\densrest$ above is larger because it adds, in quadrature, the
$\SI{8}{\kilo\gram\per\cubic\metre}$ by which $\densrest$ moves across the
admissible planes; that term belongs to the interval on $\densrest$ but not
to the budget on $C$, which is evaluated at a fixed plane. A further
term comes from the choice of spin axis: taking the rotational potential about
the recovered principal axis instead of the observed pole moves the contrast
by $3.5\%$ (Section~\ref{sec:methods}). That term is a modelling choice
independent of the other two, so it combines with them in quadrature, giving
$6.9\%$; added linearly, as the most conservative reading would have it, the
total is $9.5\%$. We note
that the agreement of the central value with published contrasts,
discussed in Section~\ref{sec:itokawa:comparison}, is well inside this
interval and should not be quoted to a precision it does not support.

The derived quantities are
%
\begin{equation}
\Dmass = (2.75 \pm 0.48)\times10^{9}~\mathrm{kg}
       = 7.7 \pm 1.3\%~\text{of}~\Mtot ,
\qquad
\dCOM = 16.9 \pm \SI{2.9}{\metre} .
\label{eq:itokawaderived}
\end{equation}
%
Both follow from $\densrest$ through
Equations~\eqref{eq:excessidentity} and~\eqref{eq:excessmass} and are not
independent measurements; they are reported because
Section~\ref{sec:itokawa:comparison} compares the second against an
external constraint, and because Section~\ref{sec:constrains} uses the
first. Computed directly from the shape model, the centre of figure lies at
$x_{\mathrm{cf}} = \SI{0.36}{\metre}$, within half a metre of the origin of
the supplied model; inverting Equation~\eqref{eq:excessmass} at this plane
gives $x_R - x_{\mathrm{cf}} = \SI{219.3}{\metre}$ against a region centroid
at $\SI{219.7}{\metre}$, consistent to the rounding.

Two figures should be recorded alongside these. The amplification factor
of Equation~\eqref{eq:amplification} is $\ampfac = 12.0$ at this plane,
so a fractional scatter quoted for $\densrest$ corresponds to twelve
times that in $\Dmass$. And the displacement statistic of
Equation~\eqref{eq:sstat} is $\Sstat = 5.8$, which is the quantity that
separates a detection from a uniform body: for a uniform body
$\densrest \equiv \densbulk$ holds identically, $\Sstat$ is zero however
small the scatter, and no amount of stability in $\densrest$ is evidence
of anything.

\begin{figure}[htbp]
\centering
\includegraphics[width=0.72\textwidth]{fig03_dispersion_curves.pdf}
\caption{Normalized geopotential dispersion against region density, each
body at its own reference plane, normalized by its own bulk density and by
its own uniform-density dispersion so that the five are comparable. Each
body has its own line style and its own marker at the minimum, so the five
remain distinguishable in monochrome. Itokawa has a deep, broad well,
reaching $0.831$ of
the uniform value at a region density $1.39$ times the bulk; the other four
have shallow wells pinned near it, the deepest of them Arrokoth at $0.987$.
Every curve is the run its section reports, the Arrokoth one with the buried
facets of Section~\ref{sec:methods:shapes} excluded, so the improvements in
the legend are those of Table~\ref{tab:nulls}.}
\label{fig:curves}
\end{figure}

\subsection{Resolution stability of the inferred density}
\label{sec:itokawa:mesh}

Paper~IV established that the objective function is stable under
decimation. That does not establish that the density recovered by
minimizing it is stable, which is a separate question and the one that
matters here. We repeated the inversion on meshes of $20{,}000$,
$30{,}000$, $40{,}000$ and $49{,}152$ facets (Figure~\ref{fig:meshconv}).

\begin{figure}[htbp]
\centering
\includegraphics[width=0.92\textwidth]{fig04_mesh_convergence.pdf}
\caption{Recovered density against facet count with the dividing plane held
fixed. Top: Itokawa. The three finest meshes share the plane
$\xsplit = \SI{0.1608}{\kilo\metre}$ and are flat to the resolution of the
scan; the coarsest, which the detector placed at
$\SI{0.1702}{\kilo\metre}$, is shown as a single point and differs by the
plane and not by the resolution. Bottom: Arrokoth at its neck, unchanged
across all three meshes. The shaded band is $\pm$ one scan step, which is
the resolution of the answer; a change smaller than the band is not a
change.}
\label{fig:meshconv}
\end{figure}

The three finest meshes return the same dividing plane,
$\SI{0.1608}{\kilo\metre}$, and on those three the recovered $\denshead$
is unchanged. We state what that does and does not mean. Because no
sub-grid refinement is applied, ``unchanged'' bounds the variation at
half a scan step, $\SI{12.5}{\kilo\gram\per\cubic\metre}$ or $0.45\%$;
it does not resolve it. The accompanying figure, that $\densrest$ moves
by $\SI{0.1}{\kilo\gram\per\cubic\metre}$ or $0.005\%$, is a weaker
statement still: with $\denshead$ pinned to the same grid point and
$\Mtot$ fixed, Equation~\eqref{eq:massbalance} makes $\densrest$ a
function of $\VR$ and $\Vtot$ alone, so the figure measures how little
the enclosed volumes change under decimation rather than how little the
inversion changes. Both are worth reporting; neither is a measurement of
the objective's resolution sensitivity.

The coarsest mesh is the informative case. At $20{,}000$ facets the
detector returns $\SI{0.1702}{\kilo\metre}$ instead of
$\SI{0.1608}{\kilo\metre}$, a displacement of $\SI{9.4}{\metre}$.
Table~\ref{tab:itokawasweep} gives
$\mathrm{d}\denshead/\mathrm{d}\xsplit \simeq
\SI{10.9}{\kilo\gram\per\cubic\metre\per\metre}$ over that interval, so
the displacement alone would move the recovered head density by about
$\SI{100}{\kilo\gram\per\cubic\metre}$, or $3.6\%$---eight times the
bound obtained on the three finer meshes. The value actually returned on
that mesh is $\SI{2900}{\kilo\gram\per\cubic\metre}$, at the plane
$\xsplit = \SI{0.1702}{\kilo\metre}$ that the detector returns there. That is
exactly the value the unreduced mesh returns at
$\SI{0.1700}{\kilo\metre}$ (Table~\ref{tab:itokawasweep}), so the
disagreement between the two meshes is attributable to the plane in full and
to the objective function not at all: each returns the same density at its
own plane, and they differ only because the detector places the plane
$\SI{9.4}{\metre}$ apart on them.
Resolution sensitivity therefore enters this calculation through the
plane detector rather than through the objective function, and a study
that fixed the plane externally would not see it at all. We regard this
as the more useful statement of the two, and it is the reason
Section~\ref{sec:guidelines} recommends reporting the detected plane on
every mesh rather than the density alone.

At the reference plane the improvement of $16.9\%$ is $12.5$ times the
value the envelope of Paper~IV takes on Itokawa specifically, $1.35\%$ over
the ladder from $2{,}000$ to $42{,}000$ facets with quadric decimation. It is
$2.9$ times the upper edge of the cross-body range, $5.9\%$, but that edge is
attained on Eros, so the ratio to it is the weaker of the two comparisons.
The envelope is a scale, not a threshold
(Section~\ref{sec:intro:relaxation}), and the comparison establishes only
that the improvement is not of the size that a change of discretization
would produce.

\subsection{Sensitivity to the dividing plane}
\label{sec:itokawa:sweep}

The plane is chosen, not measured. We therefore swept it from
$\SI{0.130}{\kilo\metre}$ to $\SI{0.240}{\kilo\metre}$ and repeated the
inversion at each position. Table~\ref{tab:itokawasweep} gives the
result. It is the central piece of evidence in this paper, and
Section~\ref{sec:constrains} is built on it.

A statement about the scan is needed before the outer planes are read.
The interval searched was $4000$~kg\,m$^{-3}$ for
$\xsplit \le \SI{0.1800}{\kilo\metre}$ and $8000$~kg\,m$^{-3}$ for
$\xsplit \ge \SI{0.1850}{\kilo\metre}$, the widening being made in
anticipation of the steep rise in recovered density at the outer planes
rather than in response to a boundary minimum: at
$\SI{0.1850}{\kilo\metre}$ the recovered value of $3075$ lies at $38\%$
of its scan limit, whereas at $\SI{0.1800}{\kilo\metre}$, where the
interval was not widened, $3025$ lies at $76\%$ of its own. The two
intervals share the same $\SI{25}{\kilo\gram\per\cubic\metre}$ grid, so
widening adds candidate points without displacing existing ones and no
discontinuity is introduced at the changeover. Every plane in
Table~\ref{tab:itokawasweep} returns an interior minimum under the rule
of Section~\ref{sec:methods:minimisation}; the closest approach to any
scan boundary is $3025$ against a limit of $4000$, a margin of
$\SI{975}{\kilo\gram\per\cubic\metre}$. No row is a boundary minimum and
none is discarded.

Three features of the table should be read before its statistics. First,
the recovered head density rises monotonically and at an accelerating
rate: the gradient $\mathrm{d}\denshead/\mathrm{d}\xsplit$ is
$\SI{7.5}{\kilo\gram\per\cubic\metre\per\metre}$ over the first four
planes and reaches
$\SI{55.0}{\kilo\gram\per\cubic\metre\per\metre}$ over the last two,
while $\densrest$ changes by only
$-\SI{0.22}{\kilo\gram\per\cubic\metre\per\metre}$ over the admissible
range---smaller by a factor of $34$ than the gentlest part of the
$\denshead$ gradient and $250$ than the steepest. Second, the improvement
also rises monotonically, from $12.9\%$ to $26.4\%$, with no interior
turning point: the objective prefers the planes furthest along the body,
which are the planes the composition forbids
(Section~\ref{sec:itokawa:ceiling}). The reported $\improve = 16.9\%$ is
therefore the improvement \emph{at the detected neck}, not the best the
objective can do, and we describe it that way throughout. Interpolating
between the planes at $0.190$ and $\SI{0.195}{\kilo\metre}$, the improvement
at the block-density bound is $21.5\%$: the objective prefers the edge of the
admissible range to the neck by $4.6$ percentage points, a difference that
itself lies inside the mesh-resolution sensitivity envelope. Third,
$\densrest$ is not flat but traces a shallow minimum near
$\SI{0.200}{\kilo\metre}$, falling from $1865.9$ to $1849.8$ and rising
again to $\SI{1855.5}{\kilo\gram\per\cubic\metre}$; the variation is
systematic rather than noise, and the minimum lies outside the admissible
range, so it should not be read as an optimum.

\begin{table}[htbp]
\centering
\small
\setlength{\tabcolsep}{4pt}
\caption{The Itokawa split-plane sweep, thirteen planes on the unreduced
$49{,}152$-facet mesh at $\Mtot = 3.58\times10^{10}$~kg. The planes at
$\SI{0.1600}{\kilo\metre}$ and $\SI{0.1608}{\kilo\metre}$ are distinct,
separated by $\SI{0.8}{\metre}$: the first belongs to the sweep grid, the
second is the saddle returned by the detector and is the reference
solution (bold). The scan column gives the interval over which $\densR$
was searched; it was widened at $\xsplit \ge \SI{0.1850}{\kilo\metre}$ in
anticipation of the steep rise in recovered density, not in response to a
boundary minimum, and the two intervals share the same
$\SI{25}{\kilo\gram\per\cubic\metre}$ grid. Every plane returns an
interior minimum, the closest approach to a boundary being $3025$ against
a limit of $4000$. Planes below the first rule exceed the empirical
block-density bound and those below the second the LL-chondrite
grain-density ceiling of Section~\ref{sec:itokawa:ceiling}; they are
listed because the monotone rise of $\improve$ through them is the
evidence for Section~\ref{sec:plane}, not because they describe a
possible interior. The four quantity columns are ordered by how nearly
invariant they are, and the ordering is the result. $\denshead$ is
recovered on the scan grid with no sub-grid refinement, and the two
places where $\dCOM$ decreases as the plane advances are that grid rather
than a geometric effect, since $x_R - x_{\mathrm{cf}}$ increases
monotonically from $203.9$ to $\SI{262.6}{\metre}$ throughout.}

\label{tab:itokawasweep}
\setlength{\tabcolsep}{10pt}
\begin{tabular}{lrrrrrrr}
\toprule
$\xsplit$ & $\VR/\Vtot$ & $\densrest$ & $\Dmass$ & $\dCOM$
 & $\denshead$ & scan & $\improve$ \\
(km) & (\%) & (kg\,m$^{-3}$) & ($10^{9}$~kg) & (m)
 & (kg\,m$^{-3}$) & (kg\,m$^{-3}$) & (\%) \\
\midrule
$0.1300$ & $20.73$ & $1865.9$ & $2.614$ & $14.89$ & $2575$ & $300$--$4000$ & $12.9$ \\
$0.1400$ & $19.28$ & $1860.7$ & $2.707$ & $15.81$ & $2650$ & $300$--$4000$ & $14.1$ \\
$0.1500$ & $17.89$ & $1857.7$ & $2.761$ & $16.51$ & $2725$ & $300$--$4000$ & $15.4$ \\
$0.1600$ & $16.54$ & $1856.9$ & $2.774$ & $16.97$ & $2800$ & $300$--$4000$ & $16.8$ \\
$\mathbf{0.1608}$ & $\mathbf{16.43}$ & $\mathbf{1858.1}$ & $\mathbf{2.753}$
 & $\mathbf{16.86}$ & $\mathbf{2800}$ & $\mathbf{300}$--$\mathbf{4000}$ & $\mathbf{16.9}$ \\
$0.1700$ & $15.18$ & $1854.1$ & $2.825$ & $17.66$ & $2900$ & $300$--$4000$ & $18.2$ \\
$0.1800$ & $13.80$ & $1850.9$ & $2.881$ & $18.40$ & $3025$ & $300$--$4000$ & $19.6$ \\
$0.1850$ & $13.08$ & $1853.0$ & $2.844$ & $18.37$ & $3075$ & $300$--$8000$ & $20.3$ \\
$0.1900$ & $12.36$ & $1852.5$ & $2.852$ & $18.63$ & $3150$ & $300$--$8000$ & $21.0$ \\
\cmidrule(l){1-8}
\multicolumn{8}{l}{\footnotesize\itshape
 below: $\denshead$ exceeds the empirical block-density bound,
 $3220 \pm \SI{220}{\kilo\gram\per\cubic\metre}$} \\
$0.1950$ & $11.62$ & $1850.2$ & $2.893$ & $19.11$ & $3250$ & $300$--$8000$ & $21.7$ \\
$0.2000$ & $10.87$ & $1849.8$ & $2.901$ & $19.38$ & $3350$ & $300$--$8000$ & $22.4$ \\
\cmidrule(l){1-8}
\multicolumn{8}{l}{\footnotesize\itshape
 below: $\denshead$ exceeds the LL-chondrite grain density,
 $3540 \pm \SI{130}{\kilo\gram\per\cubic\metre}$} \\
$0.2200$ & $7.86$ & $1851.9$ & $2.863$ & $20.04$ & $3900$ & $300$--$8000$ & $24.8$ \\
$0.2400$ & $5.00$ & $1855.5$ & $2.798$ & $20.52$ & $5000$ & $300$--$8000$ & $26.4$ \\
\midrule
\multicolumn{8}{l}{\bfseries All thirteen planes\quad
 ($\VR$ spanning a factor $4.15$)} \\
mean       & $13.90$ & $1855.2$ & $2.805$ & $17.93$ & $3169.2$ & & \\
half-range &         & $0.434\%$ & $5.12\%$ & $15.70\%$ & $38.26\%$ & & \\
CV         &         & $\mathbf{0.251\%}$ & $\mathbf{2.95\%}$
           & $\mathbf{9.31\%}$ & $20.61\%$ & & \\
$\ampfac$  &         & \multicolumn{6}{l}{mean $11.77$, range $11.34$--$12.70$} \\
\midrule
\multicolumn{8}{l}{\bfseries Below the grain-density ceiling, $n=11$\quad 

 ($\VR$ spanning a factor $1.907$) %; used for scaling in Section~\ref{sec:constrains})
 } \\
mean       & $15.25$ & $1855.4$ & $2.800$ & $17.51$ & $2936$ & & \\
half-range &         & $0.434\%$ & $5.12\%$ & $12.82\%$ & $13.20\%$ & & \\
CV         &         & $\mathbf{0.268\%}$ & $\mathbf{3.16\%}$
           & $\mathbf{8.15\%}$ & $8.60\%$ & & \\
$\ampfac$  &         & \multicolumn{6}{l}{mean $11.80$} \\
\midrule
\multicolumn{8}{l}{\bfseries Admissible planes only, $n=9$\quad
 ($\VR$ spanning a factor $1.68$)} \\
mean       & $16.14$ & $1856.6$ & $2.779$ & $17.12$ & $2855.6$ & & \\
half-range &         & $0.404\%$ & $4.80\%$ & $10.92\%$ & $10.07\%$ & & \\
CV         &         & $\mathbf{0.252\%}$ & $\mathbf{3.00\%}$
           & $\mathbf{7.42\%}$ & $6.89\%$ & & \\
$\ampfac$  &         & \multicolumn{6}{l}{mean $11.90$} \\
\bottomrule
\end{tabular}


\vspace{4pt}
\begin{minipage}{\textwidth}
\footnotesize
\textbf{Mass-balance audit.} Every row satisfies
$\Dmass = (\densbulk-\densrest)\Vtot$ with
$\densbulk = \SI{2012.87}{\kilo\gram\per\cubic\metre}$, and independently
$\Dmass = \Ddens\,\VR$, to the printed precision
(Section~\ref{sec:methods:verify}). The $\dCOM$ column is likewise not
independent: by Equation~\eqref{eq:excessmass} it is $\Dmass$ multiplied
by the purely geometric factor $(x_R-x_{\mathrm{cf}})/\Mtot$. Its value
lies in testing a rival hypothesis, not in adding data.
\\[3pt]
\textbf{Exclusion of the dipole hypothesis.} If the conserved quantity
were the mass dipole $D = \Dmass\,(x_R-x_{\mathrm{cf}})$ rather than the
excess mass, $\dCOM$ would be the flatter column. It is not: its
coefficient of variation exceeds that of $\Dmass$ by $2.6$ over the
eleven planes below the grain-density ceiling and $3.2$ over all thirteen.
Quantitatively, $\mathrm{d}\log(x_R-x_{\mathrm{cf}})/\mathrm{d}\log\VR =
-0.244$ over the eleven planes, so dipole conservation requires a scaling
slope $s = -0.756$; the measured value is $s = -1.141 \pm 0.018$,
twenty-one standard errors away (Section~\ref{sec:constrains:slope}).
\end{minipage}
\end{table}

\subsection{A physical bound on the plane}
\label{sec:itokawa:ceiling}

Because the objective rises monotonically with plane position, it cannot
select the plane, and an external bound is needed. Two are available, and
both are bounds on the recovered head density rather than on the plane
directly; each is converted into a plane position through the measured
gradient of Table~\ref{tab:itokawasweep}, and each therefore inherits an
uncertainty that we state rather than suppress.

The hard bound is compositional. Itokawa is an LL~chondrite parent body
\citep{nakamura2011}, and no assemblage of LL material can exceed its
grain density, $3540 \pm \SI{130}{\kilo\gram\per\cubic\metre}$
\citep{consolmagno2008}. A head denser than this would require a
composition the returned samples exclude. Interpolating
Table~\ref{tab:itokawasweep} between $\SI{0.2000}{\kilo\metre}$ and
$\SI{0.2200}{\kilo\metre}$, where the local gradient is
$\SI{27.5}{\kilo\gram\per\cubic\metre\per\metre}$, places the crossing at
$\xsplit = \SI{0.2069}{\kilo\metre}$, and the $\pm130$ on the grain
density makes that $\SI{0.2069+-0.0047}{\kilo\metre}$. Two caveats
attach. The interpolation spans an interval over which
$\denshead(\xsplit)$ is convex---the gradient rises from $20$ to
$\SI{27.5}{\kilo\gram\per\cubic\metre\per\metre}$ across it---so the
secant lies above the curve and the true crossing is at a slightly larger
$\xsplit$: the figure is a conservative lower bound on the admissible
range, not a best estimate. And the grain density is a zero-porosity
limit, so a head at exactly that value would be a solid monolith, which
the surface morphology does not support.

The softer bound is empirical. The bulk density of LL~chondrite falls,
which includes their internal microporosity, is
$3220 \pm \SI{220}{\kilo\gram\per\cubic\metre}$ \citep{consolmagno2008}.
This is a meteorite measurement, not a measurement of boulders on Itokawa,
and it bounds the density of the blocks a rubble pile is built from only on
the assumption that those blocks resemble the falls.
A head region denser than the densest individual block it could be made
of would require the head to be less porous than its own constituents.
The crossing falls between $\SI{0.1900}{\kilo\metre}$ and
$\SI{0.1950}{\kilo\metre}$, where the local gradient is
$\SI{20}{\kilo\gram\per\cubic\metre\per\metre}$, at
$\xsplit = \SI{0.1935}{\kilo\metre}$; the $\pm220$ on the bound makes
that $\SI{194+-11}{\metre}$, which should not be quoted to the tenth of a
metre.

We adopt the softer bound as the edge of the admissible range, giving the
nine planes from $\SI{0.1300}{\kilo\metre}$ to
$\SI{0.1900}{\kilo\metre}$ over which Table~\ref{tab:itokawasweep}
reports its second set of statistics. The detected neck at
$\SI{160.8}{\metre}$ sits comfortably inside it, $\SI{33}{\metre}$ from
the edge and three times the uncertainty on the edge. The essential point
is not the numerical value of either bound but that a bound from outside
the method was needed at all: the objective function, left to itself,
runs to a head density that the meteorite collection forbids.

\subsection{Comparison with published determinations}
\label{sec:itokawa:comparison}

Table~\ref{tab:itokawacompare} sets our result beside the published ones.
Two cautions govern how it should be read.

First, the comparison must be made at a common plane.
\citet{kanamaru2019} divide at $\SI{0.150}{\kilo\metre}$; at that plane
our inversion returns $\denshead = 2725$, $\densrest = 1857.7$,
$C = 1.467$ and $\dCOM = \SI{16.51}{\metre}$, and those are the values
entered in the table. The reference solution of
Section~\ref{sec:itokawa:reference} belongs to a different plane and a
different row, and the two sets should not be interchanged.

Second, three of the four entries are not independent of one another.
\citet{motooka2012}, \citet{kanamaru2019} and this work all minimize a
functional of the surface geopotential, differing in the weighting and in
the surface subset. Agreement among them tests implementations, not
physics, and the recovery of $1.467$ against $1.460$ is expected rather
than confirmatory---the more so since our own interval on $C$ is
$\pm0.09$, two orders of magnitude wider than the difference. The only
entry derived from independent physics is \citet{lowry2014}, whose YORP
spin-up measurement gives $\densbody = 1750 \pm 110$ and
$\denshead = 2850 \pm \SI{500}{\kilo\gram\per\cubic\metre}$, hence a
contrast of $1.63 \pm 0.30$ if the quoted uncertainties are combined in
quadrature, and a centre-of-mass offset of $\SI{21+-12}{\metre}$. Our
$C = 1.47 \pm 0.09$ and $\dCOM = \SI{16.5+-2.9}{\metre}$ lie inside those
intervals, but the intervals are wide, and the test is correspondingly
weak: a contrast anywhere between $1.33$ and $1.93$ would have passed it.

\begin{table}[htbp]
\centering
\footnotesize
\setlength{\tabcolsep}{2pt}
\caption{Published two-density determinations for Itokawa and this work,
all quoted at or referred to the dividing plane
$\xsplit = \SI{0.150}{\kilo\metre}$. The first three rows minimize a
functional of the surface geopotential and are not independent of one
another; only \citet{lowry2014} rests on different physics. Contrasts are
computed from the tabulated densities.}
\label{tab:itokawacompare}
\begin{tabular}{lllrrrr}
\toprule
Source & Method & Surface used & $\denshead$ & $\densbody$ & $C$ & $\dCOM$ \\
       &        &              & \multicolumn{2}{c}{(kg\,m$^{-3}$)} & & (m) \\
\midrule
\citet{lowry2014}       & YORP spin-up            & ---            & $2850\pm500$ & $1750\pm110$ & $1.63\pm0.30$ & $21\pm12$ \\
\citet{kanamaru2019}    & potential variance      & global         & $2730$       & $1870$       & $1.460$       & $16$ \\
\citet{kanamaru2019pss} & potential variance      & smooth terrain & $2450$       & $1930$       & $1.269$       & --- \\
This work               & dispersion minimization & global         & $2725$       & $1857.7$     & $1.47\pm0.09$ & $16.5\pm2.9$ \\
\bottomrule
\end{tabular}
\end{table}
% The centre-of-mass offset of Kanamaru et al. (2019, PSS) is not quoted in
% the same form as the others and is left empty rather than converted.

The spread among the three equipotential-family results is itself
informative. \citet{kanamaru2019pss}, restricting the average to the
smooth terrains, obtain a contrast of $1.269$ against $1.460$ on the
global surface: a change of surface subset moves the answer by more than
our quoted interval. The choice of which facets enter
Equation~\eqref{eq:dispersion} is therefore a modelling decision of the
same order as the choice of dividing plane, and we have not explored it.

\subsection{The head--neck ambiguity}
\label{sec:itokawa:ambiguity}

The head model is one hypothesis about where the anomaly sits. The neck
model of Section~\ref{sec:methods:identity}, in which the dense region is
a slab spanning the constriction, is another, and
\citet{kanamaru2019pss} report that on the smooth terrains it fits
marginally better. On the global surface we find the opposite, and by a
margin that can be stated quantitatively.

The neck slab, with upper edge at $\SI{0.2029}{\kilo\metre}$ and lower
edge at $\SI{0.1187}{\kilo\metre}$, a width of
$w = \SI{0.0842}{\kilo\metre}$, encloses $12.05\%$ of the volume.
The inversion returns
$\densneck = 2525 \pm \SI{259}{\kilo\gram\per\cubic\metre}$ and
$\densrest = \SI{1943}{\kilo\gram\per\cubic\metre}$, a contrast of
$1.300$, with an improvement of $\improve = 2.3\%$. Four comparisons
follow.

The improvement does not clear the bar the paper sets for itself. At
$2.3\%$ it lies inside the $0.9$--$5.9\%$ mesh-resolution envelope, which
is where the four weak-signal bodies of Section~\ref{sec:nulls} also sit.
By the criterion applied to them, the neck model is not a detection, and
we do not report it as one.

The curvature interval is wider in relative terms, $\pm259$ on $2525$
being $10.3\%$ against $4.6\%$ for the head model, so the objective is
less sharply curved about its minimum and the parameter is less well
determined.

The displacement statistic is smaller and its denominator larger.
Propagating the curvature interval through the mass constraint gives
$\delta\densrest = [f/(1-f)]\times 259 =
\SI{35.5}{\kilo\gram\per\cubic\metre}$ at $f = 0.1205$, against
$\densbulk - \densrest = \SI{69.9}{\kilo\gram\per\cubic\metre}$, so
$\Sstat = 2.0$ for the neck model against $5.8$ for the head model. On
the statistic of Equation~\eqref{eq:sstat}, which is the one that does
not reward a model merely for being stable, the head model is favoured by
a factor of nearly three.

The amplification factor differs sharply between the two models and must
be quoted with either. Because the neck model places less excess mass,
$\ampfac = \densrest/(\densbulk-\densrest) = 1943/69.9 = 27.8$, against
$12.0$ for the head model. A reader comparing the fractional stability of
$\densrest$ between the two models without dividing by $\ampfac$ would
conclude that the neck model was more than twice as well constrained,
when the opposite is the case.

One consequence deserves emphasis rather than burial. The two models
return background densities of $1858$ and
$\SI{1943}{\kilo\gram\per\cubic\metre}$, differing by
$\SI{85}{\kilo\gram\per\cubic\metre}$, or $3.2$ times the combined
uncertainty on the head-model value. The quantity that
Section~\ref{sec:constrains} identifies as the one the method constrains
is therefore not invariant under a change of region shape, only under a
change of plane position within a fixed region shape. That is a narrower
claim than the stability of Table~\ref{tab:itokawasweep} taken alone
would suggest, and we make it in that narrower form.

\subsection{Insensitivity to the assumed mass}
\label{sec:itokawa:mass}

The total mass enters as a constraint, so it is worth asking how much of
the result depends on it. We repeated the reference inversion with
$\Mtot$ varied over a factor of four, from half to twice the measured
value---far beyond the published uncertainty on $GM$, which is about five
per cent \citep{abe2006}. The rotation period was held fixed. Across that
range the recovered contrast moved by $11.7\%$, from $1.404$ to $1.569$,
and the improvement by $0.75$ percentage points, from $16.30\%$ to
$17.05\%$.

Two things should be said about this insensitivity, because it is easy to
over-read. It is not surprising: the contrast is a ratio, and to first
order a change in $\Mtot$ rescales both densities together. What does
\emph{not} rescale is the amplification factor, because
Equation~\eqref{eq:ampsensitivity} makes $\ampfac$ respond to a shift in
$\densbulk$ with a gain of about thirteen. A one per cent error in the
mass therefore moves $\ampfac$ by roughly thirteen per cent, and with it
the division between the algebraic and physical parts of the stability
ratio reported in Section~\ref{sec:constrains}. The contrast is robust to
the assumed mass; the interpretation of the stability is not, and the
error budget of Section~\ref{sec:constrains} carries the mass term for
that reason.

\subsection{Sensitivity to the shape model volume}
\label{sec:itokawa:scale}

Our mesh encloses $\Vtot = \SI{0.0177856}{\cubic\kilo\metre}$ against the
$(1.84\pm0.092)\times10^{7}~\mathrm{m^{3}}$ adopted by \citet{abe2006}, a
difference of $-3.3\%$ and $-0.67\sigma$
(Section~\ref{sec:methods:shapes}). Since $\densbulk = \Mtot/\Vtot$
enters Equation~\eqref{eq:amplification} directly, the difference matters
more than its size suggests.

Rescaling the model linearly to the published volume moves the contrast
by $-0.32\%$ and the excess mass by $-0.92\%$, both under one per cent, while
$\densbulk$ falls from $2012.9$ to
$\SI{1945.7}{\kilo\gram\per\cubic\metre}$ and $\ampfac$ moves from $12.0$
to $12.4$. The rotation period was held at its measured value, so the
rescaling alters the rotational term relative to the gravitational one: it
is a physical change to the problem, not a relabelling of it.
The contrast and the excess mass are robust; $\ampfac$, and hence the
decomposition of Section~\ref{sec:constrains}, is not. We report both
values of $\ampfac$ wherever the decomposition is quoted.

\subsection{Summary of this section}
\label{sec:itokawa:summary}

Itokawa returns a dispersion improvement of $16.9\%$ at the detected
neck, with a head density of
$2800 \pm \SI{129}{\kilo\gram\per\cubic\metre}$ against a background of
$1858 \pm \SI{27}{\kilo\gram\per\cubic\metre}$, a contrast of
$1.51 \pm 0.09$, an excess mass of $7.7 \pm 1.3\%$ of the total, and a
centre-of-mass offset of $\SI{16.9+-2.9}{\metre}$. The background density
is displaced from the bulk density by $5.8$ times its own uncertainty.

The result is stable under decimation at the level the scan grid can
resolve, and the residual resolution sensitivity enters through the plane
detector rather than the objective. It is stable under a change of
dividing plane in the sense that $\densrest$ varies by a quarter of a per
cent across thirteen planes, but not in the sense that $\denshead$ does,
which varies by twenty. It is not stable under a change of region shape:
the neck model returns a background density
$\SI{85}{\kilo\gram\per\cubic\metre}$ higher, though with an improvement
inside the resolution envelope and an $\Sstat$ of $2.0$ against $5.8$, so
the head model is preferred on the paper's own criteria. It agrees with
published determinations, but three of the four rest on the same physics
as ours, and the fourth has intervals wide enough that agreement is weak
evidence. And the objective does not choose its own plane: left to itself
it runs to a head density the meteorite collection forbids, and an
external bound at $\SI{194+-11}{\metre}$ is what stops it.
%=======================================================================
%  Paper V — Section 4 (The four weak-signal bodies)   FINAL
%
%  Requires the macros listed at the head of Section 2.
%  Additional macro used here:
%      \newcommand{\densregion}{\varrho_{R}}   % already defined as \densR
%
%  This version supersedes all earlier drafts of Section 4. Principal
%  changes against ver10:
%      contrast intervals    -> now include the induced change in
%                               rho_rest, as in Eq. (itokawaref)
%      Eros "1.000 +/- 0.017"-> 1.018 +/- 0.026 (central value corrected)
%      Bennu I = 0.000 %     -> I < 10^-3 %, with the floor stated
%      negative I values     -> identified as a grid artefact and bounded
%      amplification factor  -> tabulated; ranks the bodies identically
%                               to I (Spearman rho = -1 over five bodies)
%      S statistic           -> shown to FAIL on the Eros misplaced plane,
%                               and demoted to a necessary-not-sufficient
%                               condition within a three-part criterion
%
%  All values confirmed against paperV_final.json, paper4_recheck.json,
%  v9_section2_tests.json and arro_235/250/500.json, arrokoth_clean.json.
%  Corrections in this pass: S at the misplaced Eros plane 19 -> 2.8
%  (the region holds 88 per cent of the volume, not 5-20); Arrokoth run at
%  235 kg/m3 added, reversing the sign; Arrokoth table column count fixed.
%=======================================================================

\section{The four weak-signal bodies}
\label{sec:nulls}

Itokawa is one body out of five. This section reports the other four, and
its purpose is not to add results but to bound a failure rate: if the
inversion manufactured structure wherever it was pointed, the method would
be worthless whatever it returned on Itokawa. It does not. All four return
improvements at or below the mesh-resolution sensitivity envelope of
Paper~IV, and all four return contrasts consistent with unity once the
curvature sensitivity interval is propagated through the mass constraint.

Three further things are established here, and they matter more than the
null results themselves. The amplification factor $\ampfac$ of
Equation~\eqref{eq:amplification}, computed from densities alone and
without reference to the objective function, ranks all five bodies in
exactly the order of their dispersion improvements
(Section~\ref{sec:nulls:amplification}). A deliberately misplaced plane on
Eros produces a displacement statistic $\Sstat$ three times larger than
Itokawa's, which shows that $\Sstat$ cannot be used on its own
(Section~\ref{sec:nulls:misplaced}). And on Arrokoth the recovered
contrast reverses sign along two independent axes---the assumed bulk
density and the position of the dividing plane---which is the concrete
form of the statement that nothing is invariant when the signal is weak
(Section~\ref{sec:nulls:arrokoth}).

\subsection{The four results}
\label{sec:nulls:table}

Table~\ref{tab:nulls} collects them. Two columns require comment before
the bodies are taken individually.

The contrast intervals are computed as in
Equation~\eqref{eq:itokawaref}: the curvature sensitivity interval on
$\densR$ and the change it induces in $\densrest$ through the mass
constraint move the ratio in the same direction, so their fractional
contributions add rather than combining in quadrature,
%
\begin{equation}
\frac{\delta C}{C}
 = \frac{\delta\densR}{\densR}
 + \frac{f}{1-f}\,\frac{\delta\densR}{\densrest} .
\label{eq:contrastinterval}
\end{equation}
%
On Bennu, where the region is half the body, the second term is as large
as the first and the interval is twice what the first term alone would
give. Reporting only the first term, as earlier drafts of this section
did, understates every interval in the table.

The improvement for Bennu is quoted as a bound rather than as a value.
Because no sub-grid refinement is applied
(Section~\ref{sec:methods:minimisation}), the scan need not contain the
density at which $\potvarn$ is stationary, and the reported minimum can
therefore lie marginally \emph{above} the uniform solution, giving a
negative $\improve$. This is an arithmetic impossibility for the
underlying problem---the uniform body is the special case
$\densR = \densrest$ of the two-region model, so the minimum over
$\densR$ can never exceed it---and its appearance is a measure of the
grid rather than of the physics. It is observed on Arrokoth
(Section~\ref{sec:nulls:arrokoth}), where it reaches
$\improve = -0.006\%$, and that figure is the empirical floor of the
calculation. We therefore report improvements below $10^{-3}\%$ as
bounds and do not distinguish among them.

\begin{table}[htbp]
\centering
\small
\setlength{\tabcolsep}{4pt}
\caption{The four weak-signal bodies. Planes are those of
Table~\ref{tab:planes}; ``imposed'' marks the one plane not returned by
the detector. $f = \VR/\Vtot$. Contrast intervals follow
Equation~\eqref{eq:contrastinterval} and are therefore wider than the
curvature interval on $\densR$ alone. $\ampfac$ is the amplification
factor of Equation~\eqref{eq:amplification}, computed from the densities
without reference to the objective; its sign follows that of
$\densbulk - \densrest$ and is negative where the region is the less
dense of the two. $\Sstat$ is the displacement statistic of
Equation~\eqref{eq:sstat}. Itokawa is repeated from
Section~\ref{sec:itokawa} for comparison and is not part of this section's
argument.}
\label{tab:nulls}
\setlength{\tabcolsep}{10pt}
\begin{tabular}{llrrrrrr}
\toprule
Body & Plane & $f$ & $\densR$ & $\densrest$ & $C = \densR/\densrest$
 & $\improve$ & $\dCOM$ \\
 & (km) & (\%) & \multicolumn{2}{c}{(kg\,m$^{-3}$)} & & (\%) & (m) \\
\midrule
Bennu    & $0$ (imposed)   & $50.2$ & $1190$ & $1192$   & $0.999\pm0.071$ & $<10^{-3}$ & $-0.07$ \\
Eros     & $+4.4698$       & $34.6$ & $2700$ & $2652$   & $1.018\pm0.026$ & $0.418$    & $+56.4$ \\
67P      & $+0.9517$       & $26.6$ & $565$  & $527$    & $1.073\pm0.075$ & $0.574$ & $+29.1$ \\
Arrokoth & $-4.7863$       & $66.3$ & $390$  & $419.7$  & $0.929\pm0.038$ & $1.272$ & $-141$ \\
\midrule
\itshape Itokawa & \itshape $+0.1608$ & \itshape $16.4$ & \itshape $2800$
 & \itshape $1858$ & \itshape $1.51\pm0.09$ & \itshape $16.9$ & \itshape $+16.9$ \\
\bottomrule
\end{tabular}

\vspace{4pt}
\begin{minipage}{\textwidth}
\footnotesize
Arrokoth is quoted at an assumed bulk density of
$\SI{400}{\kilo\gram\per\cubic\metre}$; the dependence on that assumption
is the subject of Section~\ref{sec:nulls:arrokoth}. Its contrast interval
uses the curvature interval of the run with the buried facets retained,
$\pm\SI{5.6}{\kilo\gram\per\cubic\metre}$, the rerun that excludes them
having been made without it; the two runs differ by one scan step in
$\densR$. The 67P curvature interval is $\pm\SI{28}{\kilo\gram\per\cubic\metre}$.
\end{minipage}
\end{table}

\subsection{The amplification factor ranks the bodies}
\label{sec:nulls:amplification}

A result emerges from Table~\ref{tab:nulls} that we did not anticipate
and that costs nothing to obtain. The amplification factor
$\ampfac = \densrest/(\densbulk-\densrest)$ is computed from two
densities. It involves no surface average, no facet areas, no rotation
rate, and no evaluation of $\potvarn$; given a shape model, a mass and a
dividing plane, it is available before the objective function is ever
called. Yet it orders the five bodies exactly as the dispersion
improvement does.

\begin{table}[htbp]
\centering
\small
\caption{The amplification factor against the dispersion improvement. The
two orderings are identical, with $\ampfac$ increasing as $\improve$
decreases; over the five bodies, each at the plane its
detector returned, the Spearman rank correlation is $-1.000$ exactly
($p = 0.017$, two-tailed); the ordering does not survive the inclusion of a
misplaced plane, as the text below records. $\ampfac$ diverges as the
signal vanishes, because its denominator $\densbulk - \densrest$ is
proportional to the excess mass and the excess mass goes to zero.}
\label{tab:amplification}
\setlength{\tabcolsep}{10pt}
\begin{tabular}{lrrr}
\toprule
Body & $\densbulk-\densrest$ (kg\,m$^{-3}$) & $\lvert\ampfac\rvert$
 & $\improve$ (\%) \\
\midrule
Itokawa  & $+154.8$ & $12.0$   & $16.9$ \\
Arrokoth & $-19.7$  & $21.3$   & $1.272$ \\
67P      & $+10.1$  & $52.2$   & $0.574$ \\
Eros     & $+15.8$  & $168$    & $0.418$ \\
Bennu    & $-1.0$   & $1187$   & $<10^{-3}$ \\
\bottomrule
\end{tabular}
\end{table}

The reason is not mysterious once
Equation~\eqref{eq:excessidentity} is in view. Both quantities are
monotone functions of the excess mass: $\improve$ measures how much
dispersion the excess mass removes, and $\ampfac^{-1}$ measures how large
the excess mass is relative to the body. What is useful is that the
second can be computed without the first, so that $\ampfac$ serves as a
diagnostic on any body for which a mass and a plane are available---and
that a large $\ampfac$ is a warning, not a precision. A body returning
$\ampfac = 1187$, as Bennu does, is one on which the two-region model has
found essentially no excess mass, and on which any fractional stability
quoted for $\densrest$ corresponds to a fractional instability in
$\Dmass$ three orders of magnitude larger.

We report this as a correlation over five bodies and not as a law. Five
points, four of them nulls, cannot establish a functional form, and the
ordering could be coincidental at the $1.7\%$ level. It is offered as a
cheap first check, and Section~\ref{sec:guidelines} recommends it as one.

One limit on that check is already visible in this paper and should be
stated here rather than left to be discovered. The five bodies of
Table~\ref{tab:amplification} are evaluated at planes chosen by the rule of
Section~\ref{sec:methods:neck}. At the misplaced Eros plane of
Section~\ref{sec:nulls:misplaced} the same formula gives
$\densrest = \SI{2996.9}{\kilo\gram\per\cubic\metre}$ against a bulk
density of $2667.8$, so $\ampfac = -9.11$: in absolute value smaller than
Itokawa's $12.0$, which on the ordering above would rank a spurious result
as the strongest signal in the study. The excess mass it implies,
$-12.3\%$ of the total, is larger than Itokawa's $+7.7\%$ for the same
reason. $\ampfac$ therefore orders bodies whose planes are properly placed,
and says nothing about whether a plane is properly placed; used without the
provenance checks of Section~\ref{sec:failures:plane} it will rank a
misplaced plane above a real signal. The displacement statistic behaves
better on this case, returning $2.86$ against Itokawa's $5.8$
(Section~\ref{sec:failures:plane}), but it does not reject it either.

\subsection{Eros at the detected saddle}
\label{sec:nulls:eros}

The detector returns the Himeros saddle at $x = \SI{+4.4698}{\kilo\metre}$
with a prominence of $0.216$, the second highest in the study, placing
$34.6\%$ of the volume in the region. The inversion returns
$\densR = 2700$ against $\densrest = \SI{2652}{\kilo\gram\per\cubic\metre}$,
a contrast of $1.018$, and the dispersion falls from $0.03413$ to
$0.03399$, an improvement of $0.418\%$.

The curvature sensitivity interval is
$\pm\SI{44}{\kilo\gram\per\cubic\metre}$ on $\densR$, which through
Equation~\eqref{eq:contrastinterval} gives $C = 1.018 \pm 0.026$: the
contrast is $0.7$ standard intervals from unity and consistent with a
homogeneous body. We state the central value as $1.018$ rather than as
unity, because the two are not the same statement and the first is what
the inversion returned.

Three figures place the result. The improvement of $0.418\%$ lies below
the whole of the $0.9$--$5.9\%$ envelope, so no comparison with the
envelope's edges is needed. The amplification factor is $\ampfac = 168$,
fourteen times the Itokawa value. And the displacement statistic is
$\Sstat = 0.68$: the background density sits less than one curvature
interval from the bulk density, which is what a uniform body does.

One caution attaches to the improvement. The absolute change in
$\potvarn$ is $1.4\times10^{-4}$, and the field-point offset of
Section~\ref{sec:methods:model} perturbs the surface potential on Eros by
of order $0.04\%$ of $\lvert\Ubar\rvert$. Whether that perturbation
propagates into $\potvarn$ at the $10^{-4}$ level is exactly what the
offset-sensitivity test of that section measures, and the Eros null
should be read together with its outcome rather than before it. That
test returns $0.426\%$ with the offset scaled by $10$ and $0.418\%$ with it
scaled by $0.1$, against $0.419\%$ at the nominal offset: the improvement
moves by $0.007$ percentage points, and the recovered contrast not at all,
so the Eros null does not depend on where the field points sit.

\subsection{Eros with the plane misplaced, and the limits of \texorpdfstring{$\Sstat$}{S}}
\label{sec:nulls:misplaced}

Section~\ref{sec:methods:neck} rejected the rule that selects the
smallest interior half-width, on the grounds that on a body tapering at
both ends it lands on a flank. Eros is such a body, and the rejected rule
selects $x = \SI{-10.7619}{\kilo\metre}$. Running the inversion there is
the one worked failure case of this paper, and it is more instructive
than we expected.

At that plane the inversion returns a contrast of $0.876$---a density
deficit of over twelve per cent, of the opposite sign to the Itokawa
result and larger in magnitude than three of the four nulls---with an
improvement of $\improve = 5.526\%$ and a centre-of-mass offset of
$\SI{-214.8}{\metre}$, nearly four times the offset at the true saddle.
Nothing in the output marks it as spurious. It would be reported as a
detection by any criterion based on the size of the recovered contrast.

The improvement is the first thing that catches it, and it does so by a
narrow margin that we state rather than round away. At $5.526\%$ it falls
inside the $0.9$--$5.9\%$ envelope, but only $6\%$ below the upper edge,
and that upper edge is attained on a different body. The comparison that
matters is against the value the envelope takes on Eros itself, which is
$6.03\%$ over the ladder from $2{,}000$ to $42{,}000$ facets with quadric
decimation; the misplaced plane is $0.92$ times that value, and the test
rejects it with a margin of eight per cent. The margin depends on the
decimator. With vertex clustering, which Paper~IV used, the Eros envelope is
$5.40\%$, the misplaced plane is $1.02$ times it, and the test would
\emph{not} reject the case. We use the quadric value because it is the one
consistent with the meshes of this paper, but a reader should know that the
envelope test catches this artefact by a margin smaller than the difference
between two defensible ways of computing the envelope.

The second thing that catches it is provenance: the plane was selected by
a rule the paper had already rejected on geometric grounds, before any
density was computed. That is a strong argument, but it is not available
to a reader who is handed only the output.

The displacement statistic ranks it correctly but does not reject it. The
curvature sensitivity interval at the misplaced plane is
$\pm\SI{15}{\kilo\gram\per\cubic\metre}$, three times \emph{narrower}
than at the true saddle: the objective is more sharply curved about the
spurious minimum than about the real one. Propagating it through the mass
constraint,
%
\begin{equation}
\Sstat = \frac{\lvert\densbulk-\densrest\rvert}{\sigma(\densrest)}
       = \frac{\densbulk\,(1-C)f/[1-(1-C)f]}
              {[f/(1-f)]\,\delta\densR}
       = \frac{329}{116} = 2.8 ,
\label{eq:sfail}
\end{equation}
%
evaluated with the region holding $f = 0.884$ of the volume, $C = 0.876$
and $\densbulk = \SI{2668}{\kilo\gram\per\cubic\metre}$. The misplaced
plane leaves the region on the large side of the cut, so the small
complement makes $\sigma(\densrest)$ some eight times $\delta\densR$, and
that factor, not the sharpness of the minimum, sets the value. $\Sstat = 2.8$
is half the $5.8$ returned by the Itokawa reference solution, so the
statistic does rank the artefact below the genuine result. It does not
reject it: a displacement of nearly three curvature intervals is not one a
reader handed only this number would dismiss.

We draw the conclusion rather than leave it implicit. $\Sstat$ measures
displacement in units of curvature, and a sharply curved objective makes
a small displacement look significant; sharpness of the minimum is not
evidence that the minimum is in the right place. $\Sstat$ is therefore a
necessary condition and not a sufficient one, and
Sections~\ref{sec:methods:minimisation} and~\ref{sec:itokawa:ambiguity},
which use it, should be read with that qualification. We report a result
as a detection only where all three of the following hold:
%
\begin{enumerate}
\item the improvement exceeds the value the mesh-resolution envelope
takes on that body;
\item the displacement statistic $\Sstat$ is large, in the sense of
Equation~\eqref{eq:sstat}; and
\item the dividing plane carries geometric warrant---a saddle returned by
the detector of Section~\ref{sec:methods:neck} at a prominence above
threshold, rather than a plane chosen by the investigator.
\end{enumerate}
%
The misplaced Eros plane passes the second and fails the first and third.
Itokawa passes all three. No other combination arises in this study, and
we do not claim the criterion is complete; it is the weakest set of
conditions that rejects every case we have shown to be spurious.

\subsection{Bennu: a structure the model cannot represent}
\label{sec:nulls:bennu}

Bennu is the one body whose plane we imposed. The detector refuses it at
a prominence of $0.0004$, more than two orders of magnitude below the
lowest accepted value, and correctly: Bennu is an oblate spheroid with no
constriction. We placed a plane through the centre of figure, which
divides the volume $50.2$ to $49.8$, precisely to see what the objective
returns when the model has no geometric warrant at all.

It returns nothing. The recovered densities are $\densR = 1190$ and
$\densrest = \SI{1192}{\kilo\gram\per\cubic\metre}$, a contrast of
$0.999$; the curvature interval of
$\pm\SI{42}{\kilo\gram\per\cubic\metre}$ gives $C = 0.999 \pm 0.071$
through Equation~\eqref{eq:contrastinterval}, the widest interval in the
study because the region is half the body and the second term of that
equation is as large as the first. The improvement is below the
$10^{-3}\%$ floor established in Section~\ref{sec:nulls:table}. The
centre of mass moves by $\SI{7}{\centi\metre}$ on a body of mean radius
$\SI{245}{\metre}$. The amplification factor is $1187$ and the
displacement statistic is $\Sstat = 0.02$.

This is the cleanest null in the study, and it is worth being precise
about what it does and does not show. It does not show that Bennu is
homogeneous. Radio science and the OSIRIS-REx gravity field indicate the
opposite: an underdense equatorial bulge and an underdense core
\citep{scheeres2020}. A plane-bounded two-region model is blind to both,
because a radial or axisymmetric density variation contributes almost
nothing to the difference between the two half-spaces on either side of a
plane through the centre---the deficit is distributed symmetrically about
it. The return of unity is the correct response of a model to a structure
it cannot express, not a measurement of the structure.

Read the other way, this is the strongest available bound on the
false-positive rate of the method on a real body: given a real shape
model, a real mass, a real spin state, and a plane with no geometric
justification whatever, the inversion declined to invent a contrast. The
eight synthetic bodies of Section~\ref{sec:constrains} extend that bound;
Bennu anchors it.

\subsection{67P: a signal below every systematic}
\label{sec:nulls:67p}

The detector returns the neck at $x = \SI{+0.9517}{\kilo\metre}$ with a
prominence of $0.091$, the lowest accepted in the study, placing $26.6\%$
of the volume in the small lobe---the figure validated against
\citet{jorda2016} in Section~\ref{sec:methods:external}. The inversion
returns $\densR = 565$ against $\densrest = \SI{527}{\kilo\gram\per
\cubic\metre}$, a contrast of $1.073$, with $\improve = 0.574\%$.

The number is inside the envelope, and that alone would settle it. But on
67P a stronger statement is available, because the systematic
uncertainties on this body are all larger than the signal. Repeating the
inversion at the same plane on the second of the two acceptable shape
models returns $\densrest = \SI{521}{\kilo\gram\per\cubic\metre}$ and
$\improve = 0.725\%$: the same $\densR$, but an improvement $26\%$ larger
in relative terms. The decimation from $1.3$ million facets to $50{,}000$
moves the recovered contrast by $-3.9\%$
(Section~\ref{sec:methods:shapes}), a figure some seven times the
improvement itself. And the choice of shape model changes $\Vtot$ by
$1.1\%$ between SPG and SPC products. A signal smaller than the
model-to-model scatter, smaller than the decimation error, and smaller
than the volume difference between reductions of the same data is not a
signal.
% Confirmed: Table 4.1 reports the MSPCD model, a closed surface (chi = 2)
% after decimation. The SPC 2017 model, not closed as supplied, enters only
% the cross-model comparison and no tabulated result.

The neck model is the one place in this section where an alternative
region shape does better. A slab spanning the constriction returns a
contrast of $0.903$---a \emph{deficit}, the neck less dense than the
lobes---with $\improve = 1.132\%$, twice the head-model value. Both lie
inside the envelope and neither is a detection, but the ordering is worth
recording because it is the reverse of Itokawa's, where the head model
beat the neck model by a factor of seven
(Section~\ref{sec:itokawa:ambiguity}), and it is of the opposite sign,
Itokawa's neck model having returned a contrast of $1.300$. Which region
shape the objective prefers is a property of the individual body, not a
general feature of contact binaries, and a study that tested only one
shape would not learn this.

A lower-density neck is what a compressed contact between two lobes would
not produce and what a region of enhanced porosity would; the point is
raised only to note that the sign is physically interpretable and that
the magnitude does not support interpreting it.

\subsection{Arrokoth: two independent sign reversals}
\label{sec:nulls:arrokoth}

Arrokoth returns the largest improvement of the four, $1.272\%$, and is
the only body in the study without a measured mass. The two facts are
connected, and the connection is the subject of this subsection.

\paragraph{Dependence on the assumed bulk density.}
With no mass determination \citep{spencer2020}, $\Mtot$ must be assumed,
and the assumption propagates directly into the recovered contrast.
Table~\ref{tab:arrokothdensity} gives the four values run, and
Figure~\ref{fig:arrokoth} plots them.

\begin{figure}[htbp]
\centering
\includegraphics[width=0.72\textwidth]{fig05_arrokoth_mass.pdf}
\caption{Recovered contrast on Arrokoth against the assumed bulk density,
with and without the buried facets of
Section~\ref{sec:methods:shapes}. Excluding them, the contrast crosses unity
between $235$ and $\SI{250}{\kilo\gram\per\cubic\metre}$: at the central
published estimate of \citet{keane2022} the region is the denser of the two,
and at every higher assumption it is the lighter. Retaining them shifts the
whole curve upward and moves the crossing to about
$\SI{345}{\kilo\gram\per\cubic\metre}$. The shaded region is excluded by the
lower bound near $\SI{290}{\kilo\gram\per\cubic\metre}$ that
\citet{spencer2020} derive from the requirement that mutual gravity exceed
the centrifugal acceleration; the crossing lies just below it.}
\label{fig:arrokoth}
\end{figure}

\begin{table}[htbp]
\centering
\small
\caption{Arrokoth at four assumed bulk densities, at the detected plane
$x = \SI{-4.7863}{\kilo\metre}$ with the buried facets of
Section~\ref{sec:methods:shapes} excluded. The recovered contrast rises
toward unity as the assumed density falls, reaches it at $250$, and crosses
it: at $235$ the region is the denser of the two. The improvement passes
through the negative value that marks the grid floor of
Section~\ref{sec:nulls:table} at the crossing.}
\label{tab:arrokothdensity}
\setlength{\tabcolsep}{10pt}
\begin{tabular}{rrrrrr}
\toprule
Assumed $\densbulk$ & $\densR$ & $\densrest$ & $C$ & $\Dmass/\Mtot$ & $\improve$ \\
(kg\,m$^{-3}$) & \multicolumn{2}{c}{(kg\,m$^{-3}$)} & & (\%) & (\%) \\
\midrule
$235$ & $240.0$ & $225.2$ & $1.0659$ & $+4.19$ & $+1.224$ \\
$250$ & $250.0$ & $250.0$ & $0.9999$ & $-0.01$ & $-0.006$ \\
$400$ & $390.0$ & $419.7$ & $0.9293$ & $-4.92$ & $+1.272$ \\
$500$ & $485.0$ & $529.5$ & $0.9159$ & $-5.91$ & $+2.378$ \\
\bottomrule
\end{tabular}
\end{table}

The published bulk density is not well constrained, and the constraints do
not agree. \citet{keane2022} infer $\SI{235}{\kilo\gram\per\cubic\metre}$
with a one-sigma range of $155$ to $\SI{600}{\kilo\gram\per\cubic\metre}$;
\citet{spencer2020} require more than $290$ if the neck is not in tension;
\citet{mckinnon2020} find $250$ to $500$ for a comet-like strength. We ran the
central value of \citet{keane2022} as well as three values spanning the
range of \citet{mckinnon2020}.

At $\SI{235}{\kilo\gram\per\cubic\metre}$ the recovered contrast is
$1.066$, the region the denser; at $250$ it is $0.9999$; at $400$ and $500$ it
is $0.929$ and $0.916$, the region the less dense. The sense of the answer
reverses between $235$ and $250$, inside the published one-sigma range. A
linear extrapolation from the three higher values would have put the
contrast at $235$ near $1.006$; the measured value is ten times further from
unity, so the crossing could not have been located from runs on one side of
it. The inversion returns a surplus of seven per cent at the density
Arrokoth most probably has, and a deficit of six to eight per cent at
densities near the upper end of the published range. What the method
reports about this body is, to a good approximation, a restatement of what
was assumed about it.

\paragraph{Dependence on the buried facets.}
The Arrokoth model is supplied as two unjoined shells, and $455$ facets,
carrying $0.943\%$ of the surface area, lie inside the opposite shell
(Section~\ref{sec:methods:shapes}). They are not part of the body's
surface and the relaxation argument of \citet{richardson2014} does not
apply to them, so all figures above exclude them. With them
excluded the contrast crosses unity between $235$ and
$\SI{250}{\kilo\gram\per\cubic\metre}$, and linear interpolation between
the two runs puts the crossing within a kilogram per cubic metre of $250$.
Retaining them moves the crossing to approximately
$\SI{345}{\kilo\gram\per\cubic\metre}$---a shift of some $38\%$,
produced by facets carrying under one per cent of the area. At $235$ the two
treatments agree, both returning $1.066$, so the buried facets move the
crossing without changing the answer far from it.
A sub-percent geometric decision therefore controls the sign of the
result over a wide band of assumed densities. This is why
Section~\ref{sec:methods:model} recomputes $\Ubar$ and $\sum_i A_i$ over
the retained set rather than merely omitting terms from the numerator.

\paragraph{Dependence on the plane.}
The third dependence is the one that makes this body decisive for the
argument of Section~\ref{sec:constrains}. Sweeping the plane so that the
region fraction runs from $40.0\%$ to $75.6\%$, at a fixed assumed bulk
density of $\SI{400}{\kilo\gram\per\cubic\metre}$, the recovered density
difference $\Ddens$ runs from $+25.0$ to
$\SI{-61.5}{\kilo\gram\per\cubic\metre}$. It passes through zero at
$f = 61.2\%$, where the inversion returns
$\densR = \densrest = \SI{400}{\kilo\gram\per\cubic\metre}$ exactly, a
contrast of $1.0000$, and an improvement at the negative grid floor.

The consequence for the excess mass is the point. Since
$\Dmass/\Mtot = \Ddens f/\densbulk$, the sweep carries $\Dmass$ from
$+2.5\%$ of the total to $-11.6\%$: a swing of fourteen percentage points
that changes sign in the middle. On Itokawa the corresponding quantity
varied by $\pm5.12\%$ about a non-zero mean and never approached zero
(Table~\ref{tab:itokawasweep}). Here there is no invariant to speak of,
because the quantity that was invariant on Itokawa is identically zero
somewhere inside the sweep, and a quantity that passes through zero has
no meaningful fractional stability.

The detected plane, at $f = 66.3\%$, lies $5.1$ percentage points from
that null. The contrast of $0.929$ reported in Table~\ref{tab:nulls} is
therefore principally a measure of how far the detector's plane happens
to fall from the point at which the model degenerates, and not of any
property of the body. This is the clearest demonstration in the paper of
the claim made in the abstract: where the signal is weak, nothing is
invariant, and the apparent stability documented on Itokawa is a
property of the strong-signal regime rather than of the method.

\paragraph{Resolution and the neck model.}
Two further checks complete the body. Repeating the inversion on meshes
of $15{,}000$, $25{,}000$ and $40{,}960$ facets returns
$\densR = \SI{395}{\kilo\gram\per\cubic\metre}$ on all three, with the
complementary density between $409.8$ and
$\SI{410.2}{\kilo\gram\per\cubic\metre}$: a variation of
$\SI{0.4}{\kilo\gram\per\cubic\metre}$, or $0.1\%$, and, unlike the
Itokawa case of Section~\ref{sec:itokawa:mesh}, with the detected plane
displaced by no more than one scan step across the three meshes.
This series retains the buried facets, and that alone accounts for the
difference from Table~\ref{tab:arrokothdensity}. The plane is the same, and
so is the volume fraction: $66.330\%$, $66.328\%$ and $66.326\%$ on the three
meshes against $66.326\%$ in the table, the three-thousandth of a per cent
being the volume the decimation loses. Both answers satisfy the mass
constraint at that fraction---$0.66326\times395 + 0.33674\times409.8 = 400.0$
and $0.66326\times390 + 0.33674\times419.7 = 400.0$---because $\densrest$ is
solved from it rather than fitted, so the constraint cannot distinguish them.
What distinguishes them is the surface over which the dispersion is averaged:
with the $455$ buried facets included the minimum falls one scan step higher,
at $395$ rather than $\SI{390}{\kilo\gram\per\cubic\metre}$, and the contrast
reads $0.963$ rather than $0.929$. They are two objective functions evaluated
on the same body at the same plane, not two answers to the same question. We
keep the buried facets in this series because it compares meshes with one
another, and identifying the buried set separately on each decimated mesh
would introduce a second difference between the runs.

The neck model does not survive at all. Its minimum falls on the lower
limit of the scan interval, $\SI{25}{\kilo\gram\per\cubic\metre}$, with a
slab enclosing $1.0\%$ of the volume: the small-region limit in which the
mass constraint forces an arbitrarily large density contrast for an
arbitrarily small mass, and the objective is driven by the scan boundary
rather than by the data. The boundary-minimum guard of
Section~\ref{sec:methods:minimisation} flagged it and the result is
discarded. It is one of two occasions on which that guard fired: the
other, at $\xsplit = \SI{0.2200}{\kilo\metre}$ on Itokawa, was resolved
by widening the interval and is described in
Section~\ref{sec:failures}. Here no widening helps, because the limit is
physical---a slab of one per cent of the volume cannot carry a meaningful
mass anomaly---and the model is simply inapplicable.

\subsection{What the four cases establish}
\label{sec:nulls:summary}

The nulls are not symmetrical with the detection, and it is worth being
explicit about the asymmetry. Four bodies returning nothing bound the
rate at which this method invents structure; they do not bound the rate
at which it misses structure, because on none of them is the true
interior known. Bennu is the sharpest illustration: the inversion
returned unity on a body that independent gravity data show to be
heterogeneous, and that is a false negative by construction, deliberately
provoked.

What the four do establish is the following. The method does not
manufacture a contrast when handed a plane with no geometric warrant
(Bennu). It does not manufacture one on a body whose neck is real but
whose interior contrast, if any, is below the systematics of the
available shape models (67P). It returns an improvement below the
resolution envelope on a body with the second-highest neck prominence in
the study (Eros). And on a body without a measured mass it returns a
result whose sign is controlled by three separate modelling choices, none
of which is an observation (Arrokoth).

Against this, one case shows that the method \emph{can} be made to
produce a plausible-looking contrast: the misplaced plane on Eros, which
yields a twelve per cent density deficit, an improvement just inside the
envelope, and a displacement statistic three times the value returned by
the genuine Itokawa result. That case is rejected here on the provenance
of its plane and on the envelope comparison, not on any property of its
output, and it is the reason Section~\ref{sec:failures} catalogues
failure modes rather than merely listing null results.
%=======================================================================
%  Paper V — Section 5 (What the method constrains)   FINAL
%
%  Requires the macros listed at the head of Section 2, plus:
%      \newcommand{\Dipole}{D}
%      \newcommand{\geomslope}{g}
%
%  PRINCIPAL FINDING OF THIS REVISION -------------------------------
%  Every statistic in ver10 of this section (A = 11.8, s = -1.14,
%  +/-13.1 %, +/-5.1 %) belongs to the ELEVEN planes at
%  x_split <= 0.2000 km, i.e. those below the LL-chondrite grain-density
%  ceiling. ver10 labelled this subset as the one the BLOCK-density
%  bound permits, which is the nine planes at x_split <= 0.1900 km and
%  gives different numbers. The eleven-plane subset spans a factor
%  1.907 in V_R, which is what the synthetic family's factor 1.90 was
%  matched to; the nine-plane subset spans 1.68 and all thirteen span
%  4.15. All three are now reported side by side in Table 5.1.
%
%  Superseded values:
%      "+/-13.1 % over the block-density range" -> 13.20 %, grain ceiling
%      "+/-0.435 %"                             -> 0.434 %
%      "factor of about ninety"                 -> decomposed, 3 subsets
%      "dipole displaced in observed direction" -> OPPOSITE direction
%      "no explanation for excess steepness"    -> Eq. (steepidentity)
%      "rho_rest = 1855 +/- 8"                  -> plane-placement spread;
%                                                  budget is +/- 27
%      "Delta M ... to +/-5.1 %"                -> +/-17.5 % with budget
%      dipole "three times less steady"         -> 21 sigma slope test
%
%  All values confirmed against paperV_final.json and synth_final.json.
%  Changes in this pass: the stability ratio is decomposed into identity,
%  subtraction and compensation (32.1 = 11.78 x 7.52 / 2.77), replacing the
%  single 'residual' of 2.7, which was the rho_R-vs-Delta M comparison under
%  another name; per-body synthetic table with standard errors; correlation
%  reported as Pearson with its two-tailed p and the Spearman value beside it.
%=======================================================================

\section{What the method constrains}
\label{sec:constrains}

Sections~\ref{sec:itokawa} and~\ref{sec:nulls} report one strong signal and
four weak ones. This section asks a different question: where the method
does respond strongly, what quantity has it determined? The answer is not
the density of the region, and establishing that is the principal result of
this paper.

\subsection{The observation, and the subset it depends on}
\label{sec:constrains:observation}

Across the thirteen-plane Itokawa sweep of Table~\ref{tab:itokawasweep} the
recovered region density varies by $\pm38.3\%$ about its mean while the
background density varies by $\pm0.434\%$. Both figures depend on how far
the sweep is allowed to run, and because the planes at the far end are ones
Section~\ref{sec:itokawa:ceiling} excludes on compositional grounds, the
dependence has to be made explicit before anything is inferred from it.

Table~\ref{tab:subsets} therefore reports every statistic over three nested
subsets: all thirteen planes; the eleven at
$\xsplit \le \SI{0.2000}{\kilo\metre}$, which lie below the LL-chondrite
grain-density ceiling and are the widest set that is compositionally
defensible; and the nine at $\xsplit \le \SI{0.1900}{\kilo\metre}$, which
lie below the softer empirical block-density bound adopted in
Section~\ref{sec:itokawa:ceiling}. We use the eleven-plane subset for the
scaling analysis of Sections~\ref{sec:constrains:slope}
to~\ref{sec:constrains:synthetic}, for one reason and not for its
convenience: it spans a factor $1.907$ in region volume, and the synthetic
family of Section~\ref{sec:constrains:synthetic} was built to span $1.90$.
A scaling exponent compared between sweeps of different span is not a
comparison. Where a conclusion changes between subsets we say so; the
central ones do not.

\begin{table}[htbp]
\centering
\small
\setlength{\tabcolsep}{5pt}
\caption{Sweep statistics over three nested subsets of the planes of
Table~\ref{tab:itokawasweep}. Coefficients of variation use the $n-1$
denominator. $\ampfac$ is evaluated per plane from
Equation~\eqref{eq:amplification} and averaged. The last three
ratios decompose $\CV(\densR)/\CV(\densrest)$ exactly: the identity of
Equation~\eqref{eq:excessidentity} fixes $\CV(\Dmass)/\CV(\densrest)$,
the subtraction that defines $\Ddens$ fixes $\CV(\Ddens)/\CV(\densR)$,
and only $\CV(\Ddens)/\CV(\Dmass)$, the compensation between contrast and
region volume, is a property of the objective function. The scaling exponent $s$ is
defined by Equation~\eqref{eq:slope} and fitted by ordinary least squares;
$\geomslope$ is the purely geometric slope of
Equation~\eqref{eq:geomslope}.}
\label{tab:subsets}
\begin{tabular}{lrrr}
\toprule
 & All planes & Below grain ceiling & Below block bound \\
$\xsplit$ range (km) & $0.1300$--$0.2400$ & $0.1300$--$0.2000$ & $0.1300$--$0.1900$ \\
$n$ & $13$ & $11$ & $9$ \\
span in $\VR$ & $4.15$ & $\mathbf{1.907}$ & $1.68$ \\
\midrule
$\densR$ mean (kg\,m$^{-3}$)      & $3169$   & $2936$   & $2856$ \\
\quad half-range                  & $38.26\%$ & $13.20\%$ & $10.07\%$ \\
\quad CV                          & $20.61\%$ & $8.60\%$  & $6.89\%$ \\
\midrule
$\densrest$ mean (kg\,m$^{-3}$)   & $1855.2$ & $1855.4$ & $1856.6$ \\
\quad half-range                  & $0.434\%$ & $0.434\%$ & $0.404\%$ \\
\quad CV                          & $0.251\%$ & $0.268\%$ & $0.252\%$ \\
\midrule
$\Dmass$ mean ($10^{9}$~kg)       & $2.805$  & $2.800$  & $2.779$ \\
\quad half-range                  & $5.12\%$  & $5.12\%$  & $4.80\%$ \\
\quad CV                          & $2.95\%$  & $3.16\%$  & $3.00\%$ \\
\midrule
$\dCOM$ mean (m)                  & $17.93$  & $17.51$  & $17.12$ \\
\quad half-range                  & $15.70\%$ & $12.82\%$ & $10.92\%$ \\
\quad CV                          & $9.31\%$  & $8.15\%$  & $7.42\%$ \\
\midrule
$\Ddens$ CV                       & $49.85\%$ & $23.79\%$ & $20.13\%$ \\
\midrule
$\bar{\ampfac}$                   & $11.77$  & $11.80$  & $11.90$ \\
$\CV(\densR)/\CV(\densrest)$      & $82.2$   & $32.1$   & $27.3$ \\
\quad identity, $\CV(\Dmass)/\CV(\densrest)$ & $11.76$ & $11.78$ & $11.88$ \\
\quad subtraction, $\CV(\Ddens)/\CV(\densR)$ & $2.42$ & $2.77$ & $2.92$ \\
\quad compensation, $\CV(\Ddens)/\CV(\Dmass)$ & $16.9$ & $\mathbf{7.52}$ & $6.71$ \\
\midrule
$s$                               & $-1.041\pm0.020$ & $\mathbf{-1.141\pm0.018}$ & $-1.159\pm0.025$ \\
$\geomslope$                      & $-0.177$ & $-0.244$ & $-0.262$ \\
$\mathrm{d}\log \Dipole/\mathrm{d}\log\VR$
                                  & $-0.217\pm0.020$ & $-0.385\pm0.018$ & $-0.420\pm0.025$ \\
\quad departure from dipole       & $10.9\sigma$ & $21.4\sigma$ & $16.8\sigma$ \\
\midrule
$\max \improve$                   & $26.4\%$ & $22.4\%$ & $21.0\%$ \\
\bottomrule
\end{tabular}
% Confirmed against paperV_final.json. The three values of geomslope were
% recomputed directly from the region centroids and x_cf = 0.36 m and agree.
\end{table}

Three things in that table carry the argument.

First, the background density is better determined than the region density
by a factor of $82$ over all thirteen planes and $32$ over the eleven. That
comparison is between $\densR$ and $\densrest$, which are measured in the
same units and are the two densities of the same model; they are not,
however, independent parameters, since the mass constraint of
Equation~\eqref{eq:massbalance} leaves exactly one degree of freedom, and
the pair $(\densR,\VR)$ is degenerate along $\Ddens\VR = \mathrm{const}$.
The inversion selects a point on that hyperbola only once the dividing
surface has been fixed by other means.

Second, only part of that factor is a result. Over the eleven planes it
decomposes exactly as $32.1 = 11.78 \times 7.52 / 2.77$.
Equation~\eqref{eq:excessamplify} fixes the first factor as algebra,
present whatever the objective function does. The third is also algebra:
$\Ddens$ is the difference between $\densR$ and a nearly constant
$\densrest$, so its coefficient of variation exceeds that of $\densR$ by the
ratio of their means. What remains, the factor of $7.5$ by which $\Dmass$ is
steadier than $\Ddens$, is the compensation between contrast and region
volume; since $\Dmass = \Ddens\VR$, it measures directly how far $\Ddens$
falls as $\VR$ grows, and it is the part we report. Dividing the net ratio
by $\bar{\ampfac}$ alone, as an earlier draft of this section did, gives
$\CV(\densR)/\CV(\Dmass) = 2.7$, which is the comparison of $\densR$ with
$\Dmass$ under another name: it leaves the subtraction factor in place and
understates the compensation by $2.8$.

Third, the ordering is not a property of the method. On 67P the region
density varies by $\pm3.1\%$ across its sweep while the excess mass varies
by $\pm34.4\%$---the reverse of Itokawa. On Arrokoth the density difference
$\Ddens$ passes through zero inside the sweep
(Section~\ref{sec:nulls:arrokoth}), so the sweep mean of $\Dmass$ lies
inside its own spread and no fractional range for it is meaningful; we
report instead that $\lvert\Dmass\rvert/\Mtot$ stays below $12\%$ at every
plane on both bodies. Figure~\ref{fig:rho_dm} shows the three sweeps on a
common scale.
These are half-ranges over the eight planes of the 67P sweep, from
$\xsplit = 0.70$ to $\SI{1.20}{\kilo\metre}$, which span a factor $1.504$ in
$\VR$ against the $1.907$ used for Itokawa.

\begin{figure}[p]
\centering
\includegraphics[height=0.86\textheight]{fig06_rho_vs_dM.pdf}
\caption{Recovered region density and excess mass across each split-plane
sweep, both as deviations from their own sweep mean and on a common
symmetric-log scale, against the region volume fraction. On Itokawa the
excess mass is nearly flat while the density slopes; on 67P and Arrokoth the
order is reversed. The shaded band marks $\pm20\%$. Arrokoth's excess mass
changes sign across its sweep, so its fractional range is not meaningful
and no scaling exponent exists. The Itokawa planes beyond
$\SI{0.2000}{\kilo\metre}$, which the grain-density ceiling of
Section~\ref{sec:itokawa:ceiling} excludes, are drawn open.}
\label{fig:rho_dm}
\end{figure}

\subsection{A diagnostic that does not break down}
\label{sec:constrains:slope}

Comparing fractional ranges fails when a quantity passes through zero,
because the mean then sits inside the spread. On the Arrokoth sweep
$\Ddens$ runs from $+25.0$ to $\SI{-61.5}{\kilo\gram\per\cubic\metre}$, so
its fractional range is meaningless; a ratio-based analysis of the synthetic
bodies produces apparent ratios of up to $1623$ that are artefacts of
dividing by nearly zero.

Since $\Dmass = \Ddens\,\VR$, the behaviour is captured instead by the
exponent
%
\begin{equation}
s \;=\; \frac{\mathrm{d}\log\lvert\Ddens\rvert}{\mathrm{d}\log \VR} ,
\label{eq:slope}
\end{equation}
%
fitted by ordinary least squares across the planes of a sweep. A slope of
$-1$ means $\Dmass$ is held fixed as the plane moves; a slope of $0$ means
the density difference is held fixed instead. Two conditions are required
for the exponent to be measurable: $\Ddens$ must keep its sign, and the
sweep must span an adequate range in $\VR$. We adopt a minimum span of
$1.35$, below which the logarithmic fit is unconstrained. The threshold
does no work in this paper: the Itokawa sweep spans $1.907$, the 67P sweep
$1.504$ and every synthetic sweep $1.90$ by construction, so no reported
exponent lies near it, and we report the standard error with every
exponent so that the reader can judge each fit directly.

Equation~\eqref{eq:excessidentity} gives the exponent a second reading, and
it is worth stating in its differential form rather than only at the single
point $s=-1$. Since $\Dmass = \Ddens\VR$ and
$\Dmass = (\densbulk - \densrest)\Vtot$ with $\densbulk$ and $\Vtot$ fixed,
%
\begin{equation}
s + 1 \;\equiv\; \frac{\mathrm{d}\log\Dmass}{\mathrm{d}\log\VR}
      \;\equiv\; \frac{\mathrm{d}\log\!\left(\densbulk-\densrest\right)}
                      {\mathrm{d}\log\VR} ,
\label{eq:steepidentity}
\end{equation}
%
of which the special case
%
\begin{equation}
s = -1 \iff \Ddens \VR = \mathrm{const} \iff \densrest = \mathrm{const}
\label{eq:slopeequiv}
\end{equation}
%
is the statement made in earlier drafts. Equation~\eqref{eq:steepidentity}
is an identity, not a result, and two consequences follow from it that
should be kept in view. The exponent is not independent evidence about
$\densrest$: it is a re-parametrization of the same information. What it
adds is direction rather than magnitude---it distinguishes a systematic
drift in the background density from scatter about a fixed value, which a
coefficient of variation cannot do---and it remains finite where fractional
ranges do not.

On the eleven Itokawa planes below the grain-density ceiling, $s = -1.141 \pm 0.018$. On 67P,
$s = -2.80 \pm 0.27$ over eight planes spanning a factor $1.504$ in $\VR$;
the uncertainty is fifteen times that on Itokawa, as expected on a body whose
dispersion improvement is $0.574\%$. On Arrokoth the density difference
changes sign and no exponent exists.

The Itokawa exponent is steeper than $-1$, and
Equation~\eqref{eq:steepidentity} says what that steepness is. It is the
systematic drift of the background density already documented in
Section~\ref{sec:itokawa:sweep}: $\densbulk - \densrest$ rises from $147.0$
to $\SI{163.1}{\kilo\gram\per\cubic\metre}$ as the region shrinks, giving
$\mathrm{d}\log(\densbulk-\densrest)/\mathrm{d}\log\VR = -0.16$ between the
endpoints, against the $-0.141$ that the fitted exponent implies. Earlier
drafts reported that the excess steepness was unexplained; it is not a
separate phenomenon from the shallow minimum in $\densrest$, and
Equation~\eqref{eq:steepidentity} makes the two the same statement.

\subsection{Excluding conservation of the mass dipole}
\label{sec:constrains:dipole}

A distinct hypothesis has to be disposed of before the excess mass can be
called the conserved quantity. Suppose the objective function held fixed not
$\Dmass$ but the mass dipole
$\Dipole = \Dmass\,(x_R - x_{\mathrm{cf}})$, which is $\Mtot\dCOM$ by
Equation~\eqref{eq:excessmass} and is the quantity the YORP measurement of
\citet{lowry2014} constrains. The two hypotheses make different and
quantitatively separated predictions, because the geometric factor is not
constant along a sweep.

Write
%
\begin{equation}
\geomslope \;\equiv\;
 \frac{\mathrm{d}\log\!\left(x_R - x_{\mathrm{cf}}\right)}{\mathrm{d}\log\VR} ,
\label{eq:geomslope}
\end{equation}
%
which involves no density and no objective function: it is fixed by the
shape model and the plane positions alone. Then
$\mathrm{d}\log \Dipole/\mathrm{d}\log\VR = (s+1) + \geomslope$, and dipole
conservation requires that combination to vanish, hence
$s = -1 - \geomslope$. On the eleven planes below the grain ceiling
$\geomslope = -0.244$, so dipole conservation predicts $s = -0.756$.

The measured value is $s = -1.141 \pm 0.018$. The prediction and the
measurement therefore lie on \emph{opposite} sides of $-1$: the observed
exponent is steeper, dipole conservation requires shallower. An earlier
draft of this section asserted the reverse, that dipole conservation would
displace the exponent in the observed direction, and that assertion was
wrong in sign.

Stated directly, the dipole itself scales as
%
\begin{equation}
\frac{\mathrm{d}\log \Dipole}{\mathrm{d}\log\VR}
 = -0.385 \pm 0.018 ,
\label{eq:dipoleslope}
\end{equation}
%
which is twenty-one standard errors from the zero that conservation
requires. The conclusion does not depend on the choice of subset: the same
calculation over all thirteen planes gives $-0.217 \pm 0.020$, and over the
nine planes below the block-density bound $-0.420 \pm 0.025$, or $10.9$ and
$16.8$ standard errors respectively (Table~\ref{tab:subsets}). A second and
cruder statement points the same way: across the eleven planes the
centre-of-mass offset has a coefficient of variation of $8.15\%$ against
$3.16\%$ for the excess mass, so the dipole is the less steady of the two
by a factor of $2.6$ rather than the more steady.
% Confirmed: geomslope recomputed directly from the region centroids and
% x_cf = 0.36 m reproduces -0.177, -0.244 and -0.262.


\subsection{A controlled experiment}
\label{sec:constrains:synthetic}

Three real bodies with three different behaviours do not settle why. A
fourth real body could not settle it either: we cannot choose in advance
whether it will carry a signal, and no body's true interior is known, so
agreement would establish only that something had happened twice.

We therefore constructed eight synthetic contact binaries as unions of
prolate lobes of revolution, meshed directly, with the strength of the
inversion signal varied by geometry alone. It should be stated at the outset
that no density anomaly is imposed: every synthetic body is of uniform
density by construction, and what varies between them is shape. The family
divides into three groups. A near-spherical control stands in for a body
with no interior structure. A symmetric family of equal lobes is symmetric
about its own neck, so the dispersion-minimizing contrast should be close to
unity and the response weak by construction. An asymmetric family of unequal
lobes has a geometric mass asymmetry, in the sense that a uniform body of
that shape carries its volume unequally about the dividing plane, and the
inversion has something to respond to. Within each family the neck is
deepened while the enclosed volume varies by no more than $7\%$.

Two design points matter. Planes are selected by bisection on the volume
fraction rather than by position, so that every body spans the same factor
$1.90$ in region volume as the eleven Itokawa planes below the grain-density ceiling, which span
$1.907$; selection by fixed distance produces spans of only $1.13$ to
$1.54$, which is not a comparison on equal terms. The full thirteen-plane
Itokawa sweep spans $4.15$ and is not comparable with the synthetic family
on this measure, which is the reason Table~\ref{tab:subsets} exists. And the
total mass is held at $\SI{2000}{\kilo\gram\per\cubic\metre}$ times the
enclosed volume with a rotation period of $\SI{12.13}{\hour}$ throughout, so
that only the shape differs.

It should be said clearly what this experiment does not test.
Potential-dispersion minimization does not fit observed data: it assumes the
true density is the one that flattens the surface geopotential. Imposing an
arbitrary density on a synthetic body and asking the method to recover it
would fail by construction and would test nothing. What is examined here is
the structure of the objective function---which quantity it holds fixed as
the dividing plane moves, and how that depends on the strength of its
response. Because no known density anomaly is injected, the experiment
characterises the regimes of the objective function; it does not by itself
establish that an exponent near $-1$ implies a physical mass anomaly in a
real body.

\begin{table}[htbp]
\centering
\small
\setlength{\tabcolsep}{5pt}
\caption{The synthetic family, one row per body. Improvements are the
maximum over each sweep, and the Itokawa figures given for comparison are
therefore also maxima---$22.4\%$ over the eleven planes below the grain
ceiling---rather than the $16.9\%$ reported at the detected neck in
Section~\ref{sec:itokawa:reference}. Every sweep uses seven planes spanning
a factor $1.90$ in $\VR$. Standard errors on $s$ are those of the
least-squares fit. $\bar{\ampfac}$ is the amplification factor of
Equation~\eqref{eq:excessamplify} averaged over the sweep; it is what
separates the $\densR$ and $\Dmass$ half-ranges, which are therefore not
directly comparable, and it is not meaningful on the three symmetric bodies,
where $\Dmass$ passes through zero. The last column is the compensation
ratio, the half-range of $\Ddens$ divided by that of $\Dmass$.}
\label{tab:synthetic}
\begin{tabular}{lrrrrrr}
\toprule
Body & $\improve_{\max}$ & $s$ & $\bar{\ampfac}$ & $\densR$ range
 & $\Dmass$ range & $\Ddens/\Dmass$ \\
\midrule
\texttt{sphere}        & $1.24\%$  & none & $44$ & $\pm3.10\%$ & sign change & --- \\
\texttt{sym\_waist}    & $3.68\%$  & none & ---  & $\pm5.61\%$ & sign change & --- \\
\texttt{sym\_neck}     & $4.35\%$  & none & ---  & $\pm6.22\%$ & sign change & --- \\
\texttt{sym\_deep}     & $4.47\%$  & none & ---  & $\pm5.60\%$ & sign change & --- \\
\addlinespace
\texttt{asym\_waist}   & $16.67\%$ & $-2.020\pm0.109$ & $15.0$ & $\pm9.28\%$  & $\pm31.3\%$ & $1.9$ \\
\texttt{asym\_neck}    & $17.63\%$ & $-2.366\pm0.165$ & $16.2$ & $\pm9.37\%$  & $\pm39.2\%$ & $1.7$ \\
\texttt{asym\_deep}    & $19.50\%$ & $-2.069\pm0.094$ & $14.9$ & $\pm8.81\%$  & $\pm29.1\%$ & $2.0$ \\
\texttt{asym\_extreme} & $43.71\%$ & $-1.217\pm0.023$ & $8.1$  & $\pm12.78\%$ & $\pm6.2\%$  & $6.1$ \\
\midrule
\itshape Itokawa, eleven planes & \itshape $22.4\%$ & \itshape $-1.141\pm0.018$
 & \itshape $11.8$ & \itshape $\pm13.20\%$ & \itshape $\pm5.12\%$
 & \itshape $7.1$ \\
\bottomrule
\end{tabular}
\end{table}

The pattern is monotonic. Where the response is weak the density difference
changes sign across the sweep and neither quantity is invariant, because
neither has anything to be invariant about. At moderate response the
exponent sits near $-2$, so the excess mass rises as the region shrinks and
the density is the steadier of the two. Only at the strongest response does
the exponent approach $-1$ and the excess mass become the steadier quantity.
Figure~\ref{fig:slope} plots the relation with the three real bodies
overlaid.

The Pearson correlation between signal strength and exponent is $0.95$
across the $n = 4$ synthetic bodies for which an exponent exists
($p = 0.05$ two-tailed, $0.03$ one-tailed); the four bodies without an
exponent cannot enter it. The Spearman rank correlation over the same four
is only $0.40$ ($p = 0.6$), because the three clustered bodies are not in
rank order: the Pearson value is carried by the single body at $43.7\%$. Four points, three of them clustered between $16.7\%$ and
$19.5\%$, cannot establish a functional form, and we return to this in
Section~\ref{sec:constrains:limits}.


\begin{figure}[htbp]
\centering
\includegraphics[width=\textwidth]{fig07_slope_vs_signal.pdf}
\caption{Scaling exponent $s=\mathrm{d}\log\lvert\Ddens\rvert/
\mathrm{d}\log\VR$ against signal strength for the synthetic family, with
the real bodies overlaid. A slope of $-1$ means the excess mass is held
fixed and $0$ that the density difference is. Bodies whose density
difference changes sign have no exponent and are placed on the lower row:
the symmetric synthetic family and Arrokoth. The exponent approaches $-1$
only as the signal strengthens, with a Pearson correlation of $0.95$ over the four
synthetic bodies that have one. Itokawa, at $-1.141\pm0.018$, is matched
only by the strongest synthetic body. 67P does not follow the pattern: its
exponent of $-2.80$ is steeper than the moderate group while its dispersion
improvement, $0.574\%$, places it among the weak bodies, which in the
synthetic family have no exponent at all.}
\label{fig:slope}
\end{figure}

The single synthetic body that reproduces Itokawa's behaviour is the one
with the strongest response, and it reproduces it on every measure:
exponent $-1.22$ against $-1.141$, density half-range $\pm12.78\%$ against
$\pm13.20\%$, excess-mass half-range $\pm6.22\%$ against $\pm5.12\%$. The
agreement is close on the eleven-plane subset and would not be on the
others---the thirteen-plane density half-range is $\pm38.3\%$---which is
the sharpest illustration of why the subset has to be declared before the
comparison is made.

One real body does not follow the pattern and we record it rather than
absorb it. 67P returns a measurable exponent of $-2.80$ while its
dispersion improvement, $0.574\%$, lies in the range where every synthetic
body has no exponent at all. Either the correspondence between improvement
and exponent is not one-to-one, or the 67P exponent is being measured on a
sweep whose $\Ddens$ is small enough throughout that the fit is dominated by
the systematics catalogued in Section~\ref{sec:nulls:67p}. We cannot
separate the two with the present data.

\subsection{The claim, restricted to the strong-signal regime}
\label{sec:constrains:claim}

Where the dispersion responds strongly to a contrast, the quantity the
inversion determines is the background density $\densrest$, and equivalently
through Equation~\eqref{eq:excessidentity} the mass anomaly $\Dmass$; the
pair $(\densR, \VR)$ is degenerate along $\Ddens \VR = \mathrm{const}$.
Where the response is weak, $\densrest$ drifts with the dividing plane and
nothing is determined. This accounts for why Itokawa behaves as it does and
why Arrokoth and the symmetric synthetic bodies do not, without appeal to
anything peculiar to any of them; 67P is the case it does not account for.

What the synthetic experiment establishes is a property of the objective
function in a regime, not a physical inference. It shows that when the
minimization responds strongly, the quantity it pins down is $\Dmass$ rather
than $\densR$. Whether a strong response on a real body \emph{implies} a
physical mass anomaly is a separate question that these data do not settle,
since no known anomaly was injected. On Itokawa the case for a physical
anomaly rests elsewhere: on the independent YORP determination of
\citet{lowry2014}, which makes no appeal to surface relaxation---and which,
by Section~\ref{sec:constrains:dipole}, constrains a quantity the objective
function demonstrably does not hold fixed, so the two lines of evidence are
genuinely independent rather than two readings of the same one.

The practical consequence is a reporting rule. Publish $\densrest$ and
$\Dmass$, which the method constrains, each with the amplification factor
$\ampfac$ and the plane at which it was evaluated; publish $\densR$ only
together with the dividing surface that produced it, since it is the ratio
of a constrained mass to a volume the investigator chose.

For Itokawa the numbers are these, and the distinction between the two kinds
of interval matters more than either value. Over the eleven planes below
the grain ceiling the background density has a mean of
$\SI{1855}{\kilo\gram\per\cubic\metre}$ with a plane-placement half-range of
$\SI{8}{\kilo\gram\per\cubic\metre}$, and the excess mass a mean of
$2.80\times10^{9}$~kg with a half-range of $5.1\%$. Those spreads describe
how little the answer moves when the plane moves; they are not
uncertainties, and quoting them as such---as an earlier draft of this
section did, in the form $1855\pm8$---understates the interval by more than
a factor of three. The uncertainty is the one assembled in
Equation~\eqref{eq:budget}, dominated by the curvature sensitivity interval:
$\densrest = 1858\pm\SI{27}{\kilo\gram\per\cubic\metre}$ and
$\Dmass = (2.75\pm0.48)\times10^{9}$~kg, or $7.7\pm1.3\%$ of the total, both
evaluated at the reference plane.

This also frames the head--neck ambiguity of
Section~\ref{sec:itokawa:ambiguity} correctly. A mass anomaly does not
specify the region that carries it. Two structures placing comparable excess
mass at comparable distance from the centre of figure will produce
comparable centre-of-mass offsets and comparable dispersion reductions, and
the method cannot separate them. That is not a defect of this implementation
but a property of what the data determine. It is also why the background
densities returned by the two region shapes differ by
$\SI{85}{\kilo\gram\per\cubic\metre}$: the invariance established here holds
under a change of plane position within a fixed region shape, and not under
a change of region shape.

\subsection{Limits of this evidence}
\label{sec:constrains:limits}

Five limits should be stated plainly.

The experiment varies shape, not density. Every synthetic body is uniform,
so what the exponent tracks is the strength of the inversion's response, and
the identification of that response with a physical mass anomaly is an
interpretation rather than a measured result. A version of this experiment
that built a body whose surface had relaxed onto an equipotential of a
planted interior would test the interpretation directly;
Section~\ref{sec:discussion:rate} explains why we do not report one.

The exponent is not independent of the background density.
Equation~\eqref{eq:steepidentity} makes $s+1$ identical to the logarithmic
drift of $\densbulk-\densrest$, so the scaling analysis and the stability
analysis of Section~\ref{sec:constrains:observation} are two readings of the
same data and should not be counted as two pieces of evidence. What the
exponent adds is the distinction between drift and scatter, and a
well-defined quantity where fractional ranges diverge. The dipole test of
Section~\ref{sec:constrains:dipole} is not subject to this objection,
because it compares the measured exponent against a value fixed by geometry
alone.

Four of the eight synthetic bodies produce a measurable exponent, and three
of those cluster between $16.7\%$ and $19.5\%$ improvement while the fourth
sits at $43.7\%$. The correlation of $0.95$ therefore rests on a single
high-leverage point, and with $n=4$ its nominal $p$ of $0.03$ should not be
read as strong. The \emph{direction} of the effect is supported; its
functional form is not determined, and a threshold anywhere between roughly
$20\%$ and $44\%$ fits the data equally well.

The numerical threshold does not transfer. The synthetic bodies require
about $44\%$ improvement to reach excess-mass invariance while Itokawa
reaches it at $22\%$ over the eleven planes below the grain ceiling, or $26\%$ if the
compositionally excluded planes are counted. They are smooth surfaces of
revolution; real bodies carry roughness, which Paper~IV found contributes an
error floor of order two percentage points. No threshold value should be
quoted as universal, and we quote none.

The experiment uses the same forward model in both directions. It therefore
examines the structure of the objective function, not the accuracy of the
forward model. That accuracy was established separately in Paper~IV against
an independent field, and is tested again here against an analytic ellipsoid
in Section~\ref{sec:discussion:ellipsoid}.

\subsection{Sensitivity to the assumed mass as a candidate diagnostic}
\label{sec:constrains:mass}

The mass-assumption tests of Sections~\ref{sec:itokawa:mass}
and~\ref{sec:nulls:arrokoth} were run for a different purpose, but placed
side by side they suggest a second, provisional way of separating a
strong-signal case from a weak one. The comparison has to be made carefully,
because the obvious reading of it is the wrong one.

\begin{table}[htbp]
\centering
\small
\caption{Response of the recovered contrast to the assumed total mass. The
final column is the logarithmic sensitivity
$\mathrm{d}\log C/\mathrm{d}\log\Mtot$, which is what makes the two rows
comparable: the mass ranges tested differ by a factor of two, so the raw
changes in the third column do not. On the raw change Itokawa appears the
\emph{more} sensitive of the two, which is the opposite of the intended
reading; per e-fold of assumed mass, Arrokoth is the more sensitive by a
factor of $1.6$. The property that separates the two bodies is not the size
of the change but whether the resulting interval reaches unity.}
\label{tab:massdiag}
\begin{tabular}{llllr}
\toprule
Body & Mass range & Contrast range & Behaviour
 & $\mathrm{d}\log C/\mathrm{d}\log\Mtot$ \\
\midrule
Itokawa  & factor $4$ & $1.404$--$1.569$ & excludes unity throughout
 & $0.080$ \\
Arrokoth & factor $2$ & $0.9159$--$0.9999$ & reaches unity
 & $0.127$ \\
\bottomrule
\end{tabular}
\end{table}

On the body with a strong response the inference is robust to the mass
assumption, the contrast staying clear of unity across a factor of four in
assumed mass; on the one without, a factor of two carries the answer to the
threshold of a sign reversal. The diagnostic is therefore the position of
the interval relative to unity and not its width. The test is cheaper than a
split-plane sweep---three or four runs at different assumed masses, with no
bisection and no plane search---though a dividing region is of course still
required.

We put this forward as a candidate diagnostic rather than an established
one, and with a specific reservation beyond the obvious ones. It rests on
two bodies that differ in several respects besides signal strength, among
them spin rate, size and whether the mass is measured at all, and those are
not separated here. The last of them is awkward: the diagnostic is of use
precisely where the mass is unknown, and on such a body there is no measured
value against which to calibrate the interval, while on a body where the
mass is known---Itokawa---the diagnostic is not needed. Whether it
generalises is untested, and the synthetic family of
Section~\ref{sec:constrains:synthetic}, where the mass is known exactly by
construction and the signal strength can be dialled, is the natural place to
test it.
%=======================================================================
%  Paper V — Section 6 (The dividing plane is not self-determined) FINAL
%
%  Requires the macros listed at the head of Section 2.
%
%  PRINCIPAL CHANGE IN THIS REVISION ---------------------------------
%  The argument no longer rests on the extrapolated turning point. The
%  turning point quoted in ver10 (0.253-0.255 km from cubic and quartic
%  fits) falls within a few metres of the plane at which the 8000 kg/m3
%  scan ceiling of Section 2.9 begins to truncate the recovered
%  contrast (x ~ 0.257 km), and local extrapolation of the gradient
%  puts it at 0.262-0.278 km instead. It is therefore demoted to a
%  remark with its caveats stated. The argument is instead closed on
%  observed data alone: the improvement is monotone across all thirteen
%  planes, and the outermost plane exceeds the grain-density ceiling by
%  11.2 sigma, so the maximum over the admissible set is attained at a
%  boundary fixed by composition rather than by the objective.
%
%  Superseded values:
%      sigma quoted to 6 dp      -> 5 dp (grid floor is 3.5e-6)
%      I quoted to 3 dp          -> 2 dp
%      "a turning point must exist" -> not proved; limit at the tip is
%                                   finite and the minimum may be at the
%                                   endpoint
%      "genuine convergence"     -> consistency with a one-sided bound
%      "7.5 kg/m3 per metre"     -> 10.0 at the reference plane; and the
%                                   sensitivity of rho_rest, 0.28, is
%                                   the one that matters
%      eq:massconstraint         -> eq:massbalance
%
%  All values confirmed against paperV_final.json. The rate-of-fall list
%  is reconciled below: twelve gradients, not thirteen.
%=======================================================================

\section{The dividing plane is not self-determined}
\label{sec:plane}

The dividing plane is the one ingredient of the method that is neither
measured nor derived. The mass comes from radio science, the shape from a
spacecraft, the spin from lightcurves and imaging; the plane is placed by
the investigator. Section~\ref{sec:itokawa:sweep} showed that the recovered
region density depends strongly on where it is placed. This section asks
whether the objective function can place it, and finds that it cannot---not
in the weak sense that the optimum is shallow, but in the strong sense that
within the range the composition permits there is no interior optimum at
all, and the objective's preferred plane is one the meteorite collection
forbids.

\subsection{The improvement rises monotonically to the end of the sweep}
\label{sec:plane:monotonic}

Across the whole Itokawa sweep the dispersion at the minimum falls
monotonically, from $0.07549$ at $\xsplit = \SI{0.130}{\kilo\metre}$ to
$0.06373$ at $\SI{0.240}{\kilo\metre}$, against a uniform-body value of
$0.08663$: an improvement rising from $12.86\%$ to $26.44\%$ without a
single reversal in thirteen planes (Table~\ref{tab:itokawasweep},
Figure~\ref{fig:plane}).

We quote the dispersion to five decimal places and the improvement to two,
and not further, although the computation is deterministic and returns more
digits. The reason is the scan grid of
Section~\ref{sec:methods:minimisation}. Near its minimum the objective is
quadratic, $\potvarn(\densR) \simeq \potvarn^{\ast} +
\tfrac{1}{2}k(\densR-\densR^{\ast})^{2}$, and the curvature sensitivity
interval---a rise of $0.60\%$ over
$\pm\SI{130}{\kilo\gram\per\cubic\metre}$---fixes
$k = 7.1\times10^{-7}\,\potvarn^{\ast}$ per $(\mathrm{kg\,m^{-3}})^{2}$.
A half-step of $\SI{12.5}{\kilo\gram\per\cubic\metre}$ therefore inflates
the reported minimum by at most $5.5\times10^{-5}$ in relative terms, or
$3.5\times10^{-6}$ absolute. That is below the step-to-step change along the
sweep, roughly $1.1\times10^{-4}$ per ten metres, so the monotonicity is
safe; but it is above the sixth decimal place, so the sixth decimal place is
grid noise and we do not print it.

The rate of fall is not uniform. It steepens from about
$-1.11\times10^{-3}$ per ten metres at the start of the sweep to a maximum
near $\SI{0.185}{\kilo\metre}$, then slackens to $-1.06\times10^{-3}$
between $0.200$ and $\SI{0.220}{\kilo\metre}$ and to
$-6.9\times10^{-4}$ between $0.220$ and $\SI{0.240}{\kilo\metre}$. The
slackening is the first sign that the fall is not indefinite, and
Section~\ref{sec:plane:impossible} takes it up.
% Resolved: there is no fourteenth plane. The thirteen-entry list came from
% a run record in which x = 0.1608 km appears twice, the reference run and
% the sweep run at the same detected plane, differing only in the sixth
% digit of the excess mass. De-duplicated, the thirteen planes give twelve
% gradients, in units of 1e-6 per ten metres:
%   -1107 -1130 -1163 -1178 -1196 -1222 -1237 -1230 -1223 -1194 -1061 -692
% The same duplicate explains the two earlier discrepancies noted in the
% working file; the sweep driver emitted the plane set Table 3.1 prints.

\begin{figure}[htbp]
\centering
\includegraphics[width=\textwidth]{fig08_improvement_vs_plane.pdf}
\caption{Dispersion improvement against dividing-plane position on
Itokawa. The improvement rises monotonically across all thirteen planes
with no interior turning point. The vertical bands mark the two external
bounds of Section~\ref{sec:itokawa:ceiling}: the empirical block-density
bound at $\SI{194+-11}{\metre}$ and the LL-chondrite grain-density ceiling
at $\SI{207+-5}{\metre}$. The plane returned by the neck detector, at
$\SI{160.8}{\metre}$, is marked; the objective prefers the block-density
boundary to it by $4.6$ percentage points of improvement. The shaded
horizontal band is the $0.9$--$5.9\%$ mesh-resolution envelope of
Paper~IV, shown on the same scale to indicate that the objective's
preference among admissible planes is of the same order as its own
resolution sensitivity.}
\label{fig:plane}
\end{figure}

\subsection{Why this settles the question without extrapolation}
\label{sec:plane:impossible}

The monotonicity by itself does not show that the objective cannot choose a
plane; a function rising to the edge of the range examined might turn just
beyond it. What settles the question is that the planes the objective
prefers are excluded on grounds that have nothing to do with the objective.

At $\xsplit = \SI{0.240}{\kilo\metre}$ the recovered head density is
$\SI{5000}{\kilo\gram\per\cubic\metre}$. The LL-chondrite grain density is
$3540 \pm \SI{130}{\kilo\gram\per\cubic\metre}$
(Section~\ref{sec:itokawa:ceiling}), so the objective's best available plane
sits $11.2$ standard deviations above the density of solid LL material with
no pore space at all. It is not a marginal exclusion and it does not depend
on which of the two bounds is adopted: the same plane exceeds the empirical
block-density bound of $3220 \pm 220$ by $8.1$ standard deviations.

The argument then closes on observed data alone. The improvement is strictly
increasing over every plane in the sweep; the admissible set is the interval
below $\xsplit = \SI{194+-11}{\metre}$; therefore the improvement attains its
maximum over the admissible set at the upper boundary of that set, and the
boundary is fixed by the meteorite collection rather than by the
minimization. No statement about the behaviour of $\potvarn$ beyond the
sweep is required.

It is worth putting a number on what the objective's preference costs.
Interpolating Table~\ref{tab:itokawasweep} to the block-density boundary
gives $\improve \simeq 21.5\%$ there, against $16.89\%$ at the neck the
detector returns: the objective prefers the boundary by $4.6$ percentage
points, or $27\%$ in relative terms. That difference is real, but it lies
inside the $0.9$--$5.9\%$ mesh-resolution envelope of Paper~IV, so it is not
a discrimination the method can defend on its own terms. Over the full
admissible range the improvement varies by $8.6$ percentage points, which
does exceed the envelope; the objective discriminates among admissible
planes, but only in the direction of the boundary.

\paragraph{The turning point, and why we do not use it.}
A cubic fit to $\potvarn(\xsplit)$ over the sweep places a minimum at
$\SI{0.253}{\kilo\metre}$ and a quartic fit at $\SI{0.255}{\kilo\metre}$.
Earlier drafts of this section reported those figures as the location of the
turning point. We now report them only as a remark, for three reasons, and
the argument above is constructed so as not to depend on them.

First, they are extrapolations of $13$~m to $15$~m beyond the last plane
computed, and the agreement between the two fit orders is not evidence of
accuracy: two adjacent polynomial orders extrapolate similarly by
construction. Extrapolating the measured gradient instead places the zero at
$\SI{0.262}{\kilo\metre}$ on a quadratic through the last three planes and
$\SI{0.278}{\kilo\metre}$ on a straight line through the last two, which is
$8$ to $25$~m further out.

Second, the extrapolated position coincides with the plane at which the scan
ceiling begins to bind. Holding $\Dmass$ at $2.80\times10^{9}$~kg and
$\densrest$ near $\SI{1855}{\kilo\gram\per\cubic\metre}$, the ceiling of
$\SI{8000}{\kilo\gram\per\cubic\metre}$ used at the outer planes
(Section~\ref{sec:itokawa:sweep}) permits a region no smaller than
%
\begin{equation}
f_{\min} = \frac{\Dmass}{(\densR^{\max}-\densrest)\,\Vtot}
         = \frac{2.80\times10^{9}}
                {6145 \times 1.77856\times10^{7}}
         = 2.56\% ,
\label{eq:capbinds}
\end{equation}
%
which the sweep reaches at $\xsplit \approx \SI{0.257}{\kilo\metre}$. A
turning point predicted at $0.253$ to $\SI{0.255}{\kilo\metre}$ therefore
lies within a few metres of the plane beyond which the recovered contrast
would be truncated by the scan rather than determined by the data, and the
two cannot be separated in the present configuration. Re-running the outer
planes with the ceiling raised to, say,
$\SI{20000}{\kilo\gram\per\cubic\metre}$ would separate them; we have not
done so, because the result would lie far outside the compositional bounds
either way.

Third, we do not assert that a turning point must exist at all. As the plane
advances toward the tip, $\VR \to 0$ while the excess mass need not: the
region approaches a point mass of free magnitude $\Dmass$ located at the
tip, and the dispersion approaches the finite value obtained by optimizing
over that magnitude. The objective is therefore continuous on a closed
interval with a finite limit at its right endpoint, and its minimum may be
attained at that endpoint rather than at an interior stationary point.
Whether it is attained interior or at the limit is a question the present
data do not answer and the argument of this section does not need answered.

\subsection{Why this happens}
\label{sec:plane:why}

The explanation is Section~\ref{sec:constrains}. What the minimization
determines in the strong-signal regime is the background density, and with
it, through Equation~\eqref{eq:excessidentity}, the excess mass. Write the
region density as the background density plus the excess mass spread over
the region,
%
\begin{equation}
\densR = \densrest + \frac{\Dmass}{\VR} ,
\label{eq:plane:decompose}
\end{equation}
%
which is Equation~\eqref{eq:excessidentity} rearranged. As the plane
advances, $\VR$ falls, the mass constraint~\eqref{eq:massbalance} carries
$\densrest$ toward $\Mtot/\Vtot$, and the objective compensates by raising
$\densR$. Because the perturbation to the surface potential enters
Equation~\eqref{eq:superposition} through the product
$(\densR-\densrest)\UR$, and $\UR$ scales with $\VR$ for a small region, the
objective can hold that product nearly fixed while $\VR$ shrinks. It has no
reason to stop, and within the range examined it does not.

This is also the reconciliation between this section and
Section~\ref{sec:constrains}, which at first reading point in opposite
directions: one says the plane is unconstrained, the other that the answer
is stable under moving it. Both are true, and they are compatible precisely
because the stable quantity is not the one being optimized. Moving the plane
from the detected neck to the block-density boundary---a displacement of
$\SI{32.7}{\metre}$, a fifth of the plane's own position---changes the
recovered quantities as follows:
%
\begin{multline}
\denshead:\ 2800 \to 3220\ (+15.0\%),
\qquad
\Dmass:\ 2.753 \to 2.881\ (+4.65\%),
\qquad \\
\densrest:\ 1858.1 \to 1850.9\ (-0.387\%).
\label{eq:plane:move}
\end{multline}
%
The ratio of the last two is $4.65/0.387 = 12.0$, which is the amplification
factor $\ampfac$ at the reference plane, recovered here from a finite
displacement rather than from the differential
Equation~\eqref{eq:excessamplify}; the ratio of the first and last is
$38.8$, or $\ampfac$ times the residual factor of $3.2$ that measures
compensation between contrast and region volume. A plane displacement large
enough to move the region density by fifteen per cent moves the quantity the
method constrains by four parts in a thousand.

\subsection{What must supply the plane instead}
\label{sec:plane:external}

Two external criteria are available, and on Itokawa they do not contradict
one another. We put it that way deliberately, because an earlier draft of
this section described their relation as a convergence, and that is more
than the evidence supports.

The first is geomorphological. The neck detector of
Section~\ref{sec:methods:neck} returns
$\xsplit = \SI{160.8}{\metre}$ at a prominence of $0.170$, using only the
cross-sectional profile of the shape model and no density information
whatever. The second is compositional: the block-density bound of
Section~\ref{sec:itokawa:ceiling} excludes planes beyond
$\SI{194+-11}{\metre}$, and the grain-density ceiling those beyond
$\SI{207+-5}{\metre}$, using meteorite measurements and no shape information
beyond the gradient used to convert a density into a plane position.

The two are genuinely independent in their inputs, and the detected neck
falls $\SI{33}{\metre}$ inside the tighter of the two bounds, about three
times that bound's own uncertainty. But the compositional criterion is
one-sided: it excludes planes that are too far out and says nothing about
planes that are too far in. Every plane in the sweep below
$\SI{194}{\metre}$ satisfies it, which is $63.5$ of the $110$~m swept, so
under a uniform null a plane chosen at random would have satisfied it $58\%$
of the time. The correct statement is that an independently chosen plane
turns out not to be excluded, which is reassurance rather than corroboration,
and not that two methods have selected the same plane. We adopt the detected
neck because it has geometric warrant, and we use the compositional bound to
truncate the sweep, not to locate it.

One apparent inconsistency should be disposed of. The neck slab of
Section~\ref{sec:itokawa:ambiguity} has its upper edge at
$\SI{202.9}{\metre}$, beyond the block-density bound of
$\SI{194}{\metre}$. There is no conflict: the bound is a bound on the
recovered \emph{density}, converted into a plane position by
$\mathrm{d}\denshead/\mathrm{d}\xsplit$ for the \emph{head} model, and the
conversion does not transfer to a different region shape. The neck model
returns $\densneck = \SI{2525}{\kilo\gram\per\cubic\metre}$, comfortably
below both bounds, so it violates neither.

The geomorphological criterion carries its own resolution dependence, and it
propagates asymmetrically. Section~\ref{sec:itokawa:mesh} found that
decimating to $20{,}000$ facets displaces the detected plane by
$\SI{9.4}{\metre}$. At the local gradients of
Table~\ref{tab:itokawasweep}---$\SI{10.0}{\kilo\gram\per\cubic\metre\per
\metre}$ for $\denshead$ and
$-\SI{0.28}{\kilo\gram\per\cubic\metre\per\metre}$ for $\densrest$---that
displacement moves the region density by $\SI{94}{\kilo\gram\per
\cubic\metre}$, or $3.4\%$, which is comparable to the curvature sensitivity
interval of $4.6\%$; and the background density by
$\SI{2.6}{\kilo\gram\per\cubic\metre}$, or $0.14\%$, against a combined
budget of $1.45\%$, where it adds nothing in quadrature. The plane detector
is a significant error source for the quantity the method does not constrain
and a negligible one for the quantity it does.

We note that this is a weaker footing than the rest of the method. The
resolution stability of the inferred density
(Section~\ref{sec:itokawa:mesh}) and the invariance of the background
density in the strong-signal regime (Section~\ref{sec:constrains}) are
properties measured within the method. The placement of the plane is not,
and rests on judgement informed by two independent lines of evidence, only
one of which is two-sided. Any result reported from this method should state
the plane used, the criterion by which it was chosen, and the sensitivity of
the answer to moving it---and should state that sensitivity for the quantity
being reported. On Itokawa at the reference plane the three sensitivities
are $\SI{10.0}{\kilo\gram\per\cubic\metre\per\metre}$ for $\denshead$
(rising to $\SI{55}{\kilo\gram\per\cubic\metre\per\metre}$ at the outer
planes and falling to $7.5$ over the first four),
$\SI{0.28}{\kilo\gram\per\cubic\metre\per\metre}$ for $\densrest$, and
$0.18\%$ per metre for $\Dmass$. Quoting the first alone, as earlier drafts
did, states the sensitivity of the quantity
Section~\ref{sec:constrains} concludes the method does not determine.
%=======================================================================
%  Paper V — Section 7 (Failure modes)   FINAL
%
%  Requires the macros listed at the head of Section 2.
%
%  PRINCIPAL CORRECTION IN THIS REVISION -----------------------------
%  Section 7.3 fixes the region fraction at the misplaced Eros plane at
%  f = 0.8844 (complement 11.56 %), recovered exactly from the reported
%  rho_rest = 19,259 at rho_R = 500. This overturns the claim made in
%  Section 4.4 of an earlier draft that the displacement statistic S
%  reaches 19 at that plane: the correct value is 2.86, BELOW the 5.8
%  returned by Itokawa, so S ranks the artefact correctly. What fails
%  instead is the amplification factor, |A| = 9.12 < 12.0, which ranks
%  the artefact as the stronger signal of the two. Section 4.4 and
%  Table 4.2 must be revised accordingly.
%
%  Superseded values:
%      "2.1e-14 in relative terms" volume closure -> 2.0e-14 per cent
%      "S = 19 at the misplaced plane" -> 2.86 (Section 4.4)
%      complement-too-small "survives: background density"
%                                      -> destroys it; 3.83 % vs 1.38 %
%      "7.5 kg/m3 per metre"           -> 10.0 at the reference plane,
%                                         and 0.28 for rho_rest
%      3700/4125 as a flatness demonstration
%                                      -> a sub-grid location of the
%                                         minimiser at rho_R ~ 3912
%
%  All values confirmed against paperV_final.json and paper4_recheck.json.
%  Resolved in this pass: closure residual is a percentage, not a relative
%  value; the "-355/+1114" pair conflated the 0.60 % and 5.9 % levels;
%  the Eros envelope is 6.03 % (quadric) or 5.40 % (clustering); the window
%  diagnostic reproduces 10.3 % once the bounds are applied to both
%  densities, and Itokawa recomputes to 87.6 %; the Arrokoth crossover with
%  the buried facets excluded lies within 1 kg/m3 of 250.
%=======================================================================

\section{Failure modes}
\label{sec:failures}

Most of what went wrong in this work went wrong silently. A region can
enclose the wrong solid and still report an exact volume; a minimum can be
pinned to the edge of a scan and still report an improvement; a set of
facets can lie inside the body and still enter a surface average. This
section catalogues the seven failures encountered, with the quantity each
destroys and the check that catches it, because a reader applying the
method needs the failure list more than the result list.

Two of them are recorded not because they are subtle but because they were
ours: the wrong solid of Section~\ref{sec:methods:cone} survived a
volume-closure test, and the buried facets of
Section~\ref{sec:methods:shapes} survived every geometric check the
implementation then had. The others are properties of the method that any
implementation will meet.

\begin{table}[htbp]
\centering
\scriptsize
\setlength{\tabcolsep}{3.5pt}
\caption{The seven failure modes. ``Destroys'' names the quantity that
becomes untrustworthy, and is stated in terms of the quantities
Section~\ref{sec:constrains} identifies as constrained ($\densrest$ and
$\Dmass$) rather than in terms of $\densR$, which the method does not
determine in any case. Note the asymmetry between the fourth and fifth
rows: a region that is too small degrades only $\densR$, while a
complement that is too small degrades $\densrest$, and it is the second
that matters.}
\label{tab:failures}
\setlength{\tabcolsep}{5pt}
\small
\begin{tabular}{p{2.3cm}p{2.6cm}p{3.6cm}p{2.4cm}p{3.0cm}}
\toprule
Failure & Encountered & Symptom & Destroys & Caught by \\
\midrule
Wrong solid built
 & Itokawa, all planes
 & volume exact, mesh watertight, contrast inverted
 & everything
 & vertex predicate; mass-balance audit \\
\addlinespace
Degenerate facets
 & Itokawa, $\xsplit = \SI{0.130}{\kilo\metre}$
 & NaN at every field point, not merely near the plane
 & everything
 & self-announcing \\
\addlinespace
Facets not on the surface
 & Arrokoth, two unjoined shells
 & $0.943\%$ of the area buried; a sign reversal appears that is not there
 & sign of $\Dmass$
 & per-component Euler check; containment test \\
\addlinespace
Region too small
 & Itokawa, $\SI{0.240}{\kilo\metre}$, $f = 5.0\%$
 & curvature interval widens from $\pm4.6\%$ to $\pm7.1\%$
 & $\densR$ only
 & curvature interval \\
\addlinespace
Complement too small
 & Eros, $\SI{-10.76}{\kilo\metre}$, complement $11.56\%$
 & curvature interval in $\densR$ \emph{narrows} to $\pm15$, which looks like precision
 & $\densrest$, at $3.83\%$ against $1.38\%$
 & interval propagated to $\densrest$, not read in $\densR$ \\
\addlinespace
Minimum at a scan edge
 & Arrokoth neck, $f = 1.0\%$
 & $\improve = 13.6\%$ from a density pinned to the scan floor
 & everything
 & boundary-minimum guard; specific improvement \\
\addlinespace
Plane misplaced
 & Eros, $\SI{-10.76}{\kilo\metre}$
 & $C = 0.876 \pm 0.039$, excluding unity at $3.2$ intervals; $\dCOM$ reversed and four times larger
 & everything
 & envelope comparison; plane provenance; admissible-scan fraction \\
\bottomrule
\end{tabular}
\end{table}

\subsection{Building the wrong solid}
\label{sec:failures:cone}

The first failure is the one described in Section~\ref{sec:methods:cone},
and it is placed first because it is the only one that produced a
publishable-looking wrong answer across every plane simultaneously.

Closing a clipped region by fanning its boundary edges to a single interior
point produces a watertight mesh whose divergence-theorem volume agrees
with the sum of its own tetrahedra to machine precision. The
volume-closure test then in use compared exactly those two quantities, and
reported a residual of $2.0\times10^{-14}$ per cent, that is
$2.0\times10^{-16}$ in relative terms---round-off, and nothing else. That
figure lies below the double-precision epsilon because the two half-space
meshes share most of their facets with the whole-body mesh, so the round-off
is common to both sides of the comparison and cancels
(Section~\ref{sec:methods:verify}). The solid enclosed was nevertheless a radial cone over
the surface patch rather than the plane-cut lobe, too large by a factor of
$1.600$ on a unit sphere and with its centroid displaced by $0.1125\,R$.
% Resolved from the code: density_estimate.py records
% complement_closure_pct = 100 x |V_R + V_complement - V_tot| / V_tot, so the
% archived 1.95e-14 is a percentage. Section 2.6 and this subsection now agree.

On Itokawa the consequence was a sign reversal in the recovered contrast,
because raising the density of a cone adds mass along a wedge reaching
back to the centre of the body rather than in the head. Two checks would
have caught it and neither was present: the vertex predicate of
Section~\ref{sec:methods:verify}, which the cone fails outright with a
violation of $0.5$ in units where the plane sits at $x = 0.5$; and the
mass-balance audit, which would have found
$\Dmass = \Ddens\,\VR$ and $\Dmass = (\densbulk-\densrest)\Vtot$
disagreeing by the same factor $1.600$, since the volume passed to the
mass constraint and the volume generating the potential were different
solids.

The general lesson is that a closure test comparing a mesh against its own
tetrahedra tests self-consistency, not correctness. Only a test that
appeals to the region's \emph{definition}---every vertex satisfying its
defining inequality---distinguishes the two constructions.

\subsection{Degenerate facets and total NaN}
\label{sec:failures:nan}

At $\xsplit = \SI{0.130}{\kilo\metre}$ on Itokawa a vertex of the source
mesh lies on the cutting plane. Without the classification rule of
Section~\ref{sec:methods:tol}, both edges incident on that vertex register
as crossings, two coincident intersection points are created, and
zero-area triangles appear on both the clipped surface and the cap.

The failure mode is worth stating precisely because its scope is
counter-intuitive. A single null-normal facet makes the Werner--Scheeres
evaluation return NaN at \emph{every} field point on the body, not merely
at points near the plane, because the facet's contribution enters the sum
for all field points. The symptom is therefore total rather than local,
and it cannot be mistaken for a numerical inaccuracy or worked around by
discarding nearby points.

The present implementation classifies vertices within $10^{-10}$ of the
bounding-box diagonal as lying on the plane, snaps them onto it exactly,
and never duplicates them; only strict sign changes generate an
intersection. On Itokawa that tolerance is $\SI{64}{\nano\metre}$. The row
reported at $\SI{0.1300}{\kilo\metre}$ in
Table~\ref{tab:itokawasweep} comes from the fixed implementation and is
not affected; the failure is recorded here because the plane at which it
occurs is one the sweep uses, so a reader reproducing this work without
the tolerance rule will meet it at the first plane. At the reference plane
of $\SI{0.1608}{\kilo\metre}$ no degenerate facet arises with or without
the rule.

The rule does not remove slivers, only null-normal facets. Slivers remain
and are harmless, which Section~\ref{sec:methods:tol} establishes by
measurement rather than by assertion.

\subsection{Regions and complements that are too small}
\label{sec:failures:small}

As the dividing plane advances the region shrinks, and the objective
becomes progressively less able to determine its density. The curvature
sensitivity interval on $\denshead$ widens from
$\pm\SI{130}{\kilo\gram\per\cubic\metre}$ at $f = 16.43\%$, which is
$4.6\%$ of the recovered value, to $\pm242$ at $f = 7.86\%$ ($6.2\%$) and
to $\pm355$ at $f = 5.00\%$, or $7.1\%$ of the recovered value. Extracted
from the scan grid rather than from a parabola, the intervals are
$-125/+125$ at the reference plane, $-200/+225$ at $f = 7.86\%$ and
$-350/+325$ at $f = 5.00\%$: asymmetric by one grid step at the two outer
planes and symmetric at the reference plane to the resolution of the scan.
% Resolved. The pair "-355/+1114" conflated two confidence levels rather
% than describing an asymmetry: at f = 5.00 % the grid gives -350/+325 at the
% 0.60 % level and -1125/+1100 at the 5.9 % level, and 1114 is the half-width
% at the latter. The level sets are indeed the roots of a quadratic and are
% asymmetric in general, but at these planes the asymmetry is one grid step.

A second observation at $\xsplit = \SI{0.220}{\kilo\metre}$ makes the
flatness concrete in an unexpected way. Head densities of $3700$ and
$\SI{4125}{\kilo\gram\per\cubic\metre}$, differing by $425$, return
dispersions differing by $4\times10^{-6}$. That is not by itself a
demonstration that the minimum is flat: two points placed nearly
symmetrically about a parabolic minimum return nearly equal values however
sharp the minimum is. Read as such a pair, however, the observation
locates the minimiser below the scan grid. With the local curvature fixed
by the $\pm242$ interval, equality to $4\times10^{-6}$ requires the two
displacements to differ by $\SI{1.4}{\kilo\gram\per\cubic\metre}$, placing
the stationary point at $\densR^{\ast} \simeq
\SI{3912}{\kilo\gram\per\cubic\metre}$. The grid reports $3900$, which is
$\SI{12}{\kilo\gram\per\cubic\metre}$ away and so within the half-step
bound claimed in Section~\ref{sec:methods:minimisation}---an inadvertent
confirmation of that bound, and a further argument for solving the
stationarity condition in closed form rather than scanning.

The opposite failure is the more dangerous of the two, and it is the one
this paper nearly reported. At $\xsplit = \SI{-10.76}{\kilo\metre}$ on
Eros the region is not small: it holds $88.44\%$ of the volume, and it is
the \emph{complement} that is small, at $11.56\%$. The mass constraint
then forces the complement density across an enormous range as $\densR$ is
scanned, from $\SI{19259}{\kilo\gram\per\cubic\metre}$ at
$\densR = 500$ to $\SI{127}{\kilo\gram\per\cubic\metre}$ at
$\densR = \SI{3000}{\kilo\gram\per\cubic\metre}$, so that most of the scan
describes a body that cannot exist. The recovered minimum has a curvature
interval of only $\pm\SI{15}{\kilo\gram\per\cubic\metre}$, three times
narrower than at the true saddle, and that apparent precision is the
symptom.

It is a coordinate artefact. The interval is narrow in $\densR$ because
the mass constraint makes $\densrest$ enormously sensitive to $\densR$:
by Equation~\eqref{eq:massbalance},
$\delta\densrest = [f/(1-f)]\,\delta\densR$, and at $f = 0.8844$ that
factor is $7.65$. Expressed in the quantity
Section~\ref{sec:constrains} concludes the method determines, the
interval is
$\delta\densrest = \SI{115}{\kilo\gram\per\cubic\metre}$, or $3.83\%$ of
the recovered $\SI{2996}{\kilo\gram\per\cubic\metre}$---the worst
determination anywhere in this study.

Table~\ref{tab:propagated} sets the four cases side by side, and the
ordering is the point of this subsection.

\begin{table}[htbp]
\centering
\small
\caption{The curvature sensitivity interval propagated through the mass
constraint. The interval on $\densR$ narrows as the region grows and
widens as it shrinks; the propagation factor $f/(1-f)$ moves the opposite
way and by more, so the interval on $\densrest$ is \emph{smallest} where
the region is smallest. A region that is too small therefore damages only
the quantity the method does not constrain, while a complement that is
too small damages the one it does.}
\label{tab:propagated}
\setlength{\tabcolsep}{10pt}
\begin{tabular}{llrrrr}
\toprule
Body & Plane (km) & $f$ (\%) & $\delta\densR$ & $f/(1-f)$
 & $\delta\densrest/\densrest$ \\
\midrule
Itokawa & $+0.1608$  & $16.43$ & $\pm130$        & $0.197$ & $1.38\%$ \\
Itokawa & $+0.2200$  & $7.86$  & $\pm242$        & $0.085$ & $1.11\%$ \\
Itokawa & $+0.2400$  & $5.00$  & $-355/+1114$    & $0.053$ & $-3.16\%/+1.01\%$ \\
Eros    & $-10.7619$ & $88.44$ & $\pm15$         & $7.654$ & $\pm3.83\%$ \\
\bottomrule
\end{tabular}
\end{table}

The practical rule that follows is stated in
Section~\ref{sec:guidelines}: where a body admits a choice, place the
\emph{smaller} of the two pieces in the region. It costs precision in a
quantity the method does not determine and buys it in the quantity it
does.

\subsection{Minima at a scan edge}
\label{sec:failures:edge}

A minimisation reports the smallest value it finds, and if the smallest
value is at the edge of the interval searched, it reports the edge. The
guard of Section~\ref{sec:methods:minimisation} exists because this
failure produces the most convincing-looking output in the catalogue.

The clear case is the Arrokoth neck model. A slab enclosing $1.0\%$ of the
volume drives $\densR$ to the scan floor of
$\SI{25}{\kilo\gram\per\cubic\metre}$ and reports an improvement of
$13.6\%$---four-fifths of the Itokawa figure, on a body that returns
$1.272\%$ under the head model. The guard flagged it and the result is
discarded. Widening the interval does not help here, because the limit is
physical rather than numerical: with $f = 0.01$ the mass constraint gives
$\densrest = \SI{403.8}{\kilo\gram\per\cubic\metre}$ and an excess mass of
only $-0.947\%$ of the total, so an extreme contrast is being purchased
with almost no mass.

That last figure suggests a diagnostic, which we put forward as a
candidate on the same footing as the mass-sensitivity test of
Section~\ref{sec:constrains:mass}. The ratio of the improvement to the
fractional excess mass---the dispersion reduction bought per unit of
claimed anomaly---is $14.4$ for the Arrokoth neck against $2.20$ for the
Itokawa head, $0.71$ for the Eros saddle, $0.66$ for the Itokawa neck,
$0.45$ for the misplaced Eros plane, $0.31$ for 67P and $0.26$ for the
Arrokoth head. The edge minimum exceeds the next highest value by a factor
of $6.5$ and is the only case that does. The diagnostic costs nothing, but
its limitation should be stated with it: it does not flag the misplaced
plane of Section~\ref{sec:failures:plane}, whose ratio is unremarkable.

The second case is the one this paper resolved rather than discarded. At
$\xsplit = \SI{0.220}{\kilo\metre}$ on Itokawa an early run used an upper
scan limit of $\SI{3900}{\kilo\gram\per\cubic\metre}$ and returned a
minimum at that limit. Widening the interval to
$\SI{8000}{\kilo\gram\per\cubic\metre}$ returned an interior minimum at
the same tabulated value. Two things should be said about that agreement.
The run with the lower ceiling predates the archived configuration, in
which the limit at that plane is $8000$, so it cannot be reproduced from
the released driver scripts and is described here only because the
diagnostic it illustrates is the point. And the agreement is one of grid
rather than of arithmetic: with no sub-grid refinement the minimiser is
reported to the nearest $\SI{25}{\kilo\gram\per\cubic\metre}$, and
Section~\ref{sec:failures:small} places the true stationary point near
$\SI{3912}{\kilo\gram\per\cubic\metre}$, which happens to round onto the
old boundary. A boundary minimum is a warning that the interval may be too
narrow; it is not a demonstration that it is.

A third consequence of the ceiling belongs here rather than in
Section~\ref{sec:plane}. Holding $\Dmass$ near its sweep mean, the ceiling
of $8000$ admits no region smaller than $f_{\min} = 2.56\%$
(Equation~\eqref{eq:capbinds}), which the sweep would reach at
$\xsplit \approx \SI{0.257}{\kilo\metre}$. Planes beyond that point are
not merely inadmissible on compositional grounds; they are outside what
the present scan configuration can resolve at all, and no statement about
the behaviour of $\potvarn$ there is supported.

\subsection{A misplaced plane}
\label{sec:failures:plane}

The plane is the one ingredient neither measured nor derived
(Section~\ref{sec:plane}), and misplacing it is the failure that produces
the most plausible output. On Eros the rejected selection rule of
Section~\ref{sec:methods:neck}---smallest interior half-width rather than
saddle---returns $\xsplit = \SI{-10.7619}{\kilo\metre}$, on a shoulder.

The output at that plane is worth setting out in full, because every
individual figure is unremarkable and the combination is wrong:
%
\begin{equation}
C = 0.876 \pm 0.039, \qquad
\improve = 5.526\%, \qquad
\Dmass/\Mtot = -12.32\%, \qquad
\dCOM = \SI{-214.8}{\metre} .
\label{eq:misplaced}
\end{equation}
%
The contrast excludes unity by $3.2$ of its own intervals. The excess mass
is larger in magnitude than the $7.7\%$ recovered on Itokawa. The
centre-of-mass offset is nearly four times that at the true saddle and of
the opposite sign. Nothing in the numbers marks them as spurious.

Four diagnostics were available and they do not agree, which is why this
subsection exists.

\emph{The envelope comparison catches it, narrowly.} At $5.526\%$ the
improvement lies inside the $0.9$--$5.9\%$ envelope of Paper~IV, but only
$6\%$ below its upper edge, and that edge is attained on a different body.
The comparison that matters is against the value the envelope takes on
Eros itself, which is $6.03\%$ over the ladder from $2{,}000$ to $42{,}000$
facets with quadric decimation. The misplaced plane reaches $0.92$ times that
value, so the envelope test rejects it, but with vertex clustering, which
Paper~IV used, the Eros envelope is $5.40\%$ and the same result would be
$1.02$ times it and would not be rejected: the margin is smaller than the
difference between two defensible ways of computing the envelope.

Two things should be said about the figure of $6.03\%$ itself. It exceeds
the upper edge of the $0.9$--$5.9\%$ range quoted throughout this paper, and
there is no contradiction in that: the published range was measured under
vertex clustering, whereas $6.03\%$ is measured under the quadric edge
collapse used here, and Appendix~\ref{app:envelope} gives both. Where a
body-specific value is available it supersedes the cross-body range, which is
attained on a different body in any case. The corresponding Itokawa value
under quadric collapse is $1.35\%$, so the $16.9\%$ improvement reported
there is $12.5$ times its own body's envelope, and the two comparisons are
made on the same footing.
% Resolved from paper4_recheck.json: the Eros-specific envelope is 6.03 %
% with quadric decimation and 5.40 % with vertex clustering. The misplaced
% plane reaches 5.526 %, so the diagnostic catches it under the first and
% not under the second; both values are now stated in the text.

\emph{Plane provenance catches it.} The plane was selected by a rule the
paper had already rejected on geometric grounds before any density was
computed, and the saddle prominences of Section~\ref{sec:methods:neck}
separate cleanly from $0.490$ on Arrokoth to $0.0004$ on Bennu. This is a
strong argument and it is not available to a reader handed only the
output.

\emph{The displacement statistic ranks it correctly, but weakly.} With
$\densrest = \SI{2996}{\kilo\gram\per\cubic\metre}$ and the propagated
interval of Table~\ref{tab:propagated},
%
\begin{equation}
\Sstat = \frac{\lvert\densbulk-\densrest\rvert}{\sigma(\densrest)}
       = \frac{328.6}{114.8} = 2.86 ,
\label{eq:smisplaced}
\end{equation}
%
against $5.8$ for the Itokawa reference solution. The artefact therefore
scores below the result, which is the right ordering; but it also scores
above all four genuine nulls, whose values run from $0.02$ on Bennu to
$0.68$ at the Eros saddle, so $\Sstat$ places it between the two classes
rather than with the nulls. It is a necessary condition and not a
sufficient one.

\emph{The amplification factor fails outright.} At this plane
$\ampfac = \densrest/(\densbulk-\densrest) = -9.12$, so
$\lvert\ampfac\rvert = 9.1$ against $12.0$ for Itokawa. By the ordering
established over five bodies in Section~\ref{sec:nulls:amplification}, in
which a smaller $\lvert\ampfac\rvert$ accompanies a stronger response,
this plane ranks as the strongest signal in the study. It is the one case
we have found that breaks that ordering, and it bounds what the
correlation reported there can be used for: $\ampfac$ orders bodies
evaluated at properly placed planes, and says nothing about whether the
plane is properly placed.

A fifth diagnostic was used during development and should be defined
before it is relied on. The fraction of the scan interval over which the
implied $\densrest$ remains within physically plausible bounds is
%
\begin{equation}
W = \frac{1}{\Delta\densR}\,\frac{1-f}{f}\,
    \left(\densrest^{\max}-\densrest^{\min}\right) ,
\label{eq:window}
\end{equation}
%
clipped to unity, where $\Delta\densR$ is the width of the scan. It is
computable from the mass constraint alone, before any gravity evaluation,
and it collapses as $f \to 1$, which is exactly the small-complement
regime of Section~\ref{sec:failures:small}. Evaluated with one rule for both bodies---each density confined to
$0$--$\SI{3540}{\kilo\gram\per\cubic\metre}$, the grain density of the
ordinary-chondrite analogue, and to the scan---it is $87.6\%$ for the
Itokawa reference plane, where the constraint binds only through
$\densR \le 3540$, and $10.3\%$ for the misplaced Eros plane, where the
small complement confines $\densR$ to $2554$--$\SI{3017}{\kilo\gram\per\cubic\metre}$
out of a scan of $500$--$5000$. The earlier figure of $90.0\%$ for Itokawa
was quoted against a slightly different upper bound and is superseded.
% Resolved: 10.3 % is reproduced exactly by applying the bounds to BOTH
% densities, which is what the working notes did; the ratio quoted in the
% query assumed the bounds act on rho_rest alone. Recomputed under the single
% rule stated in the text, Itokawa gives 87.6 % rather than 90.0 %.

\subsection{Facets that are not on the surface}
\label{sec:failures:buried}

The Arrokoth model of \citet{porter2024} is supplied as two unjoined
shells, geometrically in contact and topologically separate, with
$\chi = 4$ over two connected components. Four hundred and fifty-five
facets, carrying $0.943\%$ of the total area, have their centroids inside
the opposite shell, and the shells interpenetrate over $0.044\%$ of the
volume.

Those facets are not on the body's surface, and the relaxation argument of
\citet{richardson2014} does not apply to them: no regolith flows there and
no geopotential is levelled there. Including them in
Equation~\eqref{eq:dispersion} is a violation of the geometry, not an
approximation within it.

The consequence is larger than the area warrants. At an assumed bulk
density of $\SI{250}{\kilo\gram\per\cubic\metre}$, retaining the buried
facets returns a contrast of $1.062$ and excluding them returns $0.9999$:
under one per cent of the area moves the answer by $6.2\%$ and creates an
apparent density excess where there is none. Across the assumed-density
range the same decision moves the crossover---the assumed bulk density at
which the contrast passes through unity---from between $235$ and
$\SI{250}{\kilo\gram\per\cubic\metre}$, and within one kilogram per cubic
metre of $250$ by interpolation between those two runs, to approximately
$\SI{345}{\kilo\gram\per\cubic\metre}$
(Section~\ref{sec:nulls:arrokoth}).


All Arrokoth results reported in this paper exclude them, with $\Ubar$ and
$\sum_i A_i$ recomputed over the retained set rather than terms merely
dropped from the numerator. The check that finds them is a containment
test of each facet centroid against the other components, and it is
required only because the source model has more than one component---which
the per-component Euler criterion of Section~\ref{sec:methods:verify}
reports, and a single global $\chi = 2$ test would not.

\subsection{What the catalogue implies}
\label{sec:failures:summary}

Three of the seven failures announce themselves: the degenerate facets
produce NaN, the edge minima are caught by the boundary guard, and the
region-too-small case declares itself through a curvature interval that
reaches $22\%$ of the recovered value. Four do not.

Of the four silent ones, the wrong solid and the buried facets are
geometric and are caught by checks on the region's \emph{definition}
rather than on its self-consistency---the vertex predicate, the
per-component Euler criterion, the containment test, and the mass-balance
audit. The small complement is caught only by propagating the curvature
interval through the mass constraint before reading it, since in the
scanned coordinate it presents as unusual precision. The misplaced plane
is caught only by the provenance of the plane and, marginally, by
comparison against the body's own resolution envelope; two of the four
statistical diagnostics available rank it below the genuine result, and
one ranks it above.

The general shape of the catalogue is that no single number distinguishes
a result from an artefact. This is why Section~\ref{sec:nulls:misplaced}
states a three-part criterion rather than a threshold, why
Section~\ref{sec:guidelines} asks for the plane, its provenance and the
propagated interval to be reported alongside any density, and why the
quantity we recommend reporting is $\densrest$ rather than $\densR$: of
the seven failures, one degrades $\densrest$ and six leave it intact or
destroy the calculation outright, whereas five of the seven corrupt
$\densR$ while leaving the output superficially reasonable.

%=======================================================================
%  Paper V — Section 8 (Discussion)   FINAL
%
%  Requires the macros listed at the head of Section 2.
%  NOTE: ver10 used \dens in Eq. (macmillan). That macro is not defined
%  anywhere in this manuscript and would break the compile. It is
%  written out as \varrho below.
%
%  PRINCIPAL ADDITION IN THIS REVISION -------------------------------
%  Section 8.2 is new. It shows, from the mass constraint alone, that
%  the first condition of Richardson and Bowling -- that the rotational
%  term be a significant fraction of the gravitational one -- is the
%  condition under which THEIR objective is not identically flat, and
%  that it has no counterpart in a fixed-mass two-region inversion,
%  where the mass constraint fixes the scale instead. This converts the
%  discomfort acknowledged in Section 8.1 from an unresolved concern
%  into a bounded one, and it costs nothing: the omega = 0 control run
%  needs no new gravity evaluation.
%
%  Superseded values and statements:
%      "2e-14 per cent" volume closure  -> 2.1e-14 in relative terms
%      "2.2e-14 per cent" sphere limit  -> relative; and it is a check
%                                          on the analytic quadrature,
%                                          not on the polyhedral code
%      "19 to 59 times more accurate"   -> quoted at the working
%                                          resolution as well, ~16-20
%      "our 16.5 m lies within it"      -> 16.5 +/- 2.9 m, and the
%                                          Lowry interval contains the
%                                          whole sweep, so the test does
%                                          not discriminate
%      head-neck as the fundamental ambiguity
%                                       -> the method separates them by
%                                          a factor 7.3 in I; the real
%                                          degeneracy is elsewhere
%      "no numerical diagnostic detects a misplaced plane"
%                                       -> one of four fails outright,
%                                          two rank it correctly
%      "no change across the three Itokawa meshes"
%                                       -> bounded at half a scan step
%
%  Resolved in this pass: the three missing dispersion pairs of Table 8.2
%  (Itokawa does carry the largest residual, 83.1 % against 98.7 %); the
%  seven rungs of the planted ladder; the Paper IV slope range, 1.8-25.6 %
%  as published against 1.8-21.9 % recomputed with clustering here.
%  Also resolved: the omega = 0 control and the equatorial column
%  (v10_s8_tests.json), and the standard errors of Tables 8.3 and 8.4. The
%  49,152-facet row was abandoned as disproportionately expensive; see the
%  note at that table.
%=======================================================================

\section{Discussion}
\label{sec:discussion}

\subsection{The standing of the relaxation premise on these bodies}
\label{sec:discussion:premise}

Everything in this paper rests on a premise that is not established for any
of the bodies examined. \citet{richardson2014} state four conditions under
which a surface may be taken as relaxed toward an equipotential: the mean
rotational force must be a significant fraction of the mean gravitational
force; a sufficiently thick, low-cohesion, mobile regolith layer must cover
most of the surface; a downslope flow disturbance source must be present;
and sufficient time must have passed since the last major surface
alteration. They are explicit that the result is not a density measurement
but the density corresponding to the most eroded state of the body given its
present shape, topography and spin.

On Itokawa three of the four are supported and the first is not. Mobile
regolith is present: the Muses Sea and the other smooth terrains cover about
a fifth of the surface \citep{fujiwara2006,kanamaru2019pss}. A disturbance
source is present: \citet{tatsumi2018} invoke impact-induced seismic shaking
to explain the depletion of small quasi-circular depressions. Time is
available: they place the crater retention age at $3$--$33$~Myr in the main
belt, in agreement with the cosmic-ray exposure ages of the returned
samples, and conclude that global resurfacing occurred long after the
disruption of the parent body. The first condition, however, fails.
\citet{richardson2014} tabulate the mean rotational potential on Itokawa as
$2.4\%$ of the mean gravitational potential, the second lowest of the bodies
they examine, and their global minimization returns a bulk density of
$\SI{330}{\kilo\gram\per\cubic\metre}$ against a measured $1950$---an error
of a factor of six, which they attribute directly to that first criterion.
% Supporting evidence for "seven": the companion conference abstract
% (Richardson, Graves & Bowling, ACM 2014) states that the study used 22
% radar-derived shape models and SEVEN small-body shape models from
% spacecraft observations. The seven quoted here are the spacecraft set, the
% only ones with measured masses. Confirm against the Icarus paper itself
% before submission; the abstract is a companion, not the paper.

This is uncomfortable, and it should be stated rather than left for a reader
to find. The single strong signal reported here is on the body where the
premise is weakest, while Bennu, which rotates fast enough to satisfy the
condition comfortably, returns a weak-signal result.
Section~\ref{sec:discussion:conditioning} shows that the discomfort is
narrower than it first appears, because the condition that fails is a
condition on the conditioning of a different inverse problem; this
subsection first sets out where the premise stands body by body.

\begin{table}[htbp]
\centering
\small
\caption{The four conditions of \citet{richardson2014}, assessed body by
body. The rotation parameter in the second column is
$q = \omega^{2}R_{v}/(G\Mtot/R_{v}^{2})$, evaluated at the
volume-equivalent radius $R_{v}=(3\Vtot/4\pi)^{1/3}$, which reduces
identically to $q = 3\omega^{2}/(4\pi G\densbulk)$: it depends on the spin
period and the bulk density alone and carries no information about the
shape. It is therefore \emph{not} the ratio at the equator of an elongated
body, which is larger by the ratio of the maximum distance from the spin
axis to $R_{v}$---a factor of $1.89$ on Itokawa and $2.09$ on
Eros---and it is not the quantity \citet{richardson2014} tabulate, which is
a ratio of mean potentials. The three are listed or distinguished
separately and must not be conflated. Arrokoth has no measured mass, so its
$q$ is a function of the assumed bulk density and is quoted at all three
values run.}
\label{tab:conditions}
\setlength{\tabcolsep}{4.5pt}
\begin{tabular}{lcccccc}
\toprule
& \multicolumn{3}{c}{(1) rotation} & (2) mobile & (3) disturbance & (4) time \\
\cmidrule(lr){2-4}
Body & $q$ at $R_{v}$ & $q$ at equator & potential
 & regolith & source & since \\
\midrule
Itokawa  & $0.037$ & $0.070$ & $0.024$ & yes & yes & yes \\
Eros     & $0.147$ & $0.308$ & $0.111$ & yes & yes & yes \\
Bennu    & $0.496$ & $0.586$ & ---     & partly & yes & uncertain \\
67P      & $0.132$ & $0.211$ & ---     & partly & yes & no \\
Arrokoth & $0.173/0.108/0.086$ & $0.337/0.211/0.169$ & ---
         & unknown & weak & yes \\
\bottomrule
\end{tabular}

\vspace{4pt}
\begin{minipage}{\textwidth}
\footnotesize
The Arrokoth entries correspond to assumed bulk densities of $250$, $400$
and $\SI{500}{\kilo\gram\per\cubic\metre}$ respectively; at the central
published estimate near $\SI{235}{\kilo\gram\per\cubic\metre}$ the value is
$0.18$, which exceeds Eros's. The equatorial column is measured on the shape models used here, as
$\omega^{2}R_{\rm eq}/(G\Mtot/R_{v}^{2})$ with $R_{\rm eq}$ the largest
distance of any vertex from the spin axis: $306.1$ and $\SI{17.6}{\kilo\metre}$
on Itokawa and Eros, against $R_{v}$ of $161.9$ and
$\SI{8.43}{\kilo\metre}$, and ratios $R_{\rm eq}/R_{v}$ of $1.89$, $2.09$,
$1.18$, $1.60$ and $1.96$ for the five bodies in order.
\end{minipage}
% Measured on the meshes (v10_s8_tests.json): R_eq = 306.1 m, 17628.6 m,
% 289.4 m, 2633.5 m, 19482.7 m for Itokawa, Eros, Bennu, 67P and Arrokoth.
% The earlier Itokawa entry used a semi-axis of 267.5 m; the mesh gives 306.1.
\end{table}

The premise is not equally plausible across the five, and the reader should
be able to see where it is weak. Table~\ref{tab:conditions} sets the four
conditions against each body; the paragraphs below give the reasoning. One
caution applies throughout: \citet{richardson2014} do not state a numerical
threshold for ``a significant fraction'', so no body is judged against a
stated cut. The evidence that Itokawa falls below it is empirical---their own
bulk-density inversion on that body errs by a factor of six---and the
evidence that Eros lies above it is that the same inversion works there. The
threshold in mean potential ratio is therefore bracketed between $0.024$ and
$0.111$ and is not located.

\emph{Itokawa} satisfies three and fails the first, as described above.

\emph{Eros} is the most favourable case on the premise as a whole. Its
rotation parameter is four times Itokawa's, its regolith is deep enough to
have formed ponded deposits, impacts supply a disturbance source, and its
surface is old. That the method returns no contrast on the body where the
premise is strongest, and where the measured gravity field independently
indicates near-uniformity \citep{yeomans2000}, is the cleanest control in
the study.

\emph{Bennu} satisfies the first condition most comfortably of all five, at
a rotation parameter approaching one half. The other three are less clear.
Its surface is boulder-rich rather than smoothly regolith-covered, though
particles are observed to be ejected and to move, so material is mobile in
some sense. And its surface may be young: \citet{scheeres2020} attribute the
reported low-density equator to recent migration and redistribution of
material, which is a statement that the surface has changed recently rather
than settled. A surface still in motion is not a relaxed surface.

\emph{67P} fails the fourth condition outright. A comet nucleus is
resurfaced continuously by sublimation, and its shape is not the product of
gravitational relaxation alone. Whatever mobile material exists is airfall
dust redeposited between perihelion passages rather than a regolith that has
flowed downslope to equilibrium. This is a further reason, independent of
the model-class argument of \citet{jorda2016}, to expect no interpretable
signal.

\emph{Arrokoth} is the least constrained, and its entry in
Table~\ref{tab:conditions} illustrates why. Its rotation parameter is not a
measured quantity at all but a function of the assumed bulk density, running
from $0.173$ to $0.086$ across the range tried; an earlier draft described
it as comparable to Eros's, which is true only at the middle of that range,
and at the central published estimate it exceeds Eros's. Beyond that,
nothing is known of its regolith, no disturbance source is established for a
primordial Kuiper Belt object, and its surface is thought to have been
little modified since accretion \citep{mckinnon2020}---which satisfies the
fourth condition in the sense of elapsed time while removing any reason to
think relaxation ever operated. With no measured mass on top of this, its
result carries the least weight of the five.

The pattern is worth stating plainly. The premise is strongest on Eros,
which returns no signal; it is weakest on Itokawa, which returns the only
strong one. The four weak-signal cases are therefore not equally
informative: Eros's is a result about Eros, while 67P's and Arrokoth's are
as much about the premise as about the bodies.

Two of the three considerations that bear on how much weight this should
carry are empirical, and are taken up here; the third is structural and is
the subject of Section~\ref{sec:discussion:conditioning}.

The first empirical consideration is that the result does not stand alone,
and the second is that this support is weaker than it first appears.
\citet{lowry2014} derive a centre-of-mass offset of $\SI{21+-12}{\metre}$
from an observed spin-up, an argument that makes no appeal to surface
relaxation at all, and our $\SI{16.5+-2.9}{\metre}$ at the same dividing
plane lies within it. Two inferences resting on different premises agree,
and by Section~\ref{sec:constrains:dipole} they are genuinely independent
rather than two readings of one quantity, since the objective function
demonstrably does not hold the mass dipole fixed. But the agreement should
not be asked to carry more than it can. The Lowry determination is itself
only $1.75$ standard deviations from zero, so it does not exclude a
homogeneous Itokawa at the two-sigma level; and its interval,
$9$ to $\SI{33}{\metre}$, contains the centre-of-mass offset at
\emph{every} plane of Table~\ref{tab:itokawasweep}, whose range is
$14.89$ to $\SI{20.52}{\metre}$, together with the $\pm2.9$ interval on each.
The external check therefore confirms the sign and the order of magnitude of
the anomaly and discriminates neither the plane nor the value.

The third empirical consideration cuts the other way, and is recorded here.
Even after the best two-region fit, Itokawa's surface remains $7.2\%$ from
an equipotential in the normalized dispersion, against $3.4\%$ for Eros with
no contrast fitted at all. Three things should be said about that
comparison before it is read as evidence that Itokawa is the less relaxed.
It is not a resolution effect: a $5.9\%$ change in $\potvarn$, the top of
the Paper~IV envelope, is $0.2$ percentage points on Eros's residual against
the $3.8$ percentage points separating the two bodies. It is, however, not
controlled for shape, since a more irregular body has a larger $\potvarn$
whatever its state of relaxation, and Itokawa is a contact binary while Eros
is not. And it is incomplete: the residuals for 67P, Bennu and Arrokoth are
not reported anywhere in this paper, so the claim that Itokawa's is the
largest of the five is not yet supported. The missing entries are the
decisive ones, because 67P and Arrokoth are also bilobate: if their
residuals are comparable to Itokawa's, the $7.2\%$ is shape, and if they are
not, it is relaxation.

\begin{table}[htbp]
\centering
\small
\caption{Normalized geopotential dispersion before and after the two-region
fit, at the plane of Table~\ref{tab:planes} for each body. The final column
is the fraction of the departure from an equipotential that the best
two-region density model does \emph{not} remove, and it is large on every
body: on Itokawa, where the fit is by far the most effective, five-sixths of
the departure remains and is not of a form any model of this class can
absorb.}
\label{tab:residual}
\setlength{\tabcolsep}{10pt}
\begin{tabular}{lrrrr}
\toprule
Body & $\potvarn$ (uniform) & $\potvarn$ (best fit) & $\improve$
 & residual \\
\midrule
Itokawa  & $0.08663$ & $0.07200$ & $16.9\%$  & $83.1\%$ \\
Eros     & $0.03413$ & $0.03399$ & $0.418\%$ & $99.6\%$ \\
Bennu    & $0.06158$ & $0.06158$ & $<10^{-3}\%$ & $100.0\%$ \\
67P      & $0.07615$ & $0.07571$ & $0.574\%$ & $99.4\%$ \\
Arrokoth & $0.05679$ & $0.05607$ & $1.272\%$ & $98.7\%$ \\
\bottomrule
\end{tabular}
\end{table}
% Filled from paperV_final.json; the Arrokoth pair is from the run with the
% buried facets excluded (arrokoth_clean.json), matching the 1.272 % quoted
% throughout. Itokawa does carry the largest residual: 83.1 % against 98.7 %
% for the next body. The superseded note read:
% overturn it.

\subsection{Why the first condition does not transfer to this inversion}
\label{sec:discussion:conditioning}

The two inference problems are not the same, and the difference can be made
exact rather than asserted. \citet{richardson2014} vary the total mass,
which changes the balance between the gravitational and rotational terms;
this work holds the measured mass fixed and varies only the internal
distribution. That distinction is what the first condition turns on.

Consider their problem first. With a single uniform density,
Equation~\eqref{eq:geopotential} reads
$\Upot_i = \varrho\,\Ufull_i + \Urot_i$, so
%
\begin{equation}
\potvarn^{2}(\varrho)
 = \frac{\varrho^{2}\operatorname{Var}(\Ufull)
        + 2\varrho\operatorname{Cov}(\Ufull,\Urot)
        + \operatorname{Var}(\Urot)}
        {\left(\varrho\,\overline{\Ufull}+\overline{\Urot}\right)^{2}} ,
\label{eq:uniformobjective}
\end{equation}
%
where variances and covariances are area-weighted over the surface. If
$\Urot \equiv 0$ this collapses to
$\potvarn = \sqrt{\operatorname{Var}(\Ufull)}\big/\lvert\overline{\Ufull}\rvert$,
which is independent of $\varrho$: the objective is identically flat and the
bulk density is not determined at all. Their first condition is precisely
the condition that Equation~\eqref{eq:uniformobjective} is not flat, and a
rotational potential ratio of $0.024$, as on Itokawa, leaves a minimum too
shallow to locate---which is why their inversion there returns $330$ against
a measured $1950$.

Now substitute the mass constraint~\eqref{eq:massbalance} into
Equation~\eqref{eq:geopotential}. With
$a \equiv \Mtot/(\Vtot-\VR)$ and $b \equiv \VR/(\Vtot-\VR)$, so that
$\densrest = a - b\densR$, the surface geopotential is affine in the single
remaining free parameter,
%
\begin{equation}
\Upot_i = \alpha_i + \beta_i\,\densR ,
\qquad
\alpha_i = a\left(\Ufull_i - \UR_i\right) + \Urot_i ,
\qquad
\beta_i = (1+b)\,\UR_i - b\,\Ufull_i .
\label{eq:affine}
\end{equation}
%
Set $\Urot \equiv 0$, as the first condition's failure would in the limit
require. Then $\alpha_i = a(\Ufull_i - \UR_i)$, and $\alpha$ and $\beta$ are
proportional only if $a(1+b) = ab$, that is only if $a = 0$. Since
$a = \Mtot/(\Vtot-\VR)$ is strictly positive for any region smaller than the
body, they are never proportional, and $\potvarn(\densR)$ remains a
non-constant ratio of quadratics with an interior stationary point. The
degeneracy that the first condition protects against does not arise here,
because the mass constraint fixes the scale of the density field that the
rotational term fixes in their problem.

This is a statement about conditioning and not about physics, and the
distinction matters. It shows that a slow rotator does not make the
fixed-mass two-region problem ill-posed; it does not show that the surface of
a slow rotator has relaxed, which remains the premise and remains unverified.
What it removes is the specific inference that because Richardson and
Bowling's method fails on Itokawa by a factor of six, a method built on their
premise must fail there too. The two failures would have different causes,
and only one of the causes is present.

Two empirical checks bear on the argument and we report the first. A
fourfold change in the assumed total mass moves the recovered contrast on
Itokawa by $11\%$ and never near unity
(Section~\ref{sec:constrains:mass}), whereas the bulk-density inference on
the same body lands at $330$ against $1950$: the fixed-mass problem is
demonstrably better conditioned on this body than the variable-mass one, by
a wide margin. The second check is direct, and we have now run it. It confirms the
derivation in the form that matters and corrects the form in which an earlier
draft anticipated it. With $\Urot$ set identically to zero the scan returns
an interior minimum at
$\denshead = \SI{2950}{\kilo\gram\per\cubic\metre}$ and
$\densrest = \SI{1828.6}{\kilo\gram\per\cubic\metre}$, a contrast of
$1.613$ and an improvement of $17.16\%$, against $2800$, $1858.1$, $1.507$
and $16.89\%$ with the rotational term retained. The structural claim is
therefore established in the output and not only in the derivation: the
two-region objective does not collapse without rotation, whereas the
single-density objective of \citet{richardson2014} becomes identically flat
under the same substitution and determines nothing at all. What supplies the
scale there is rotation; what supplies it here is the mass constraint.

What the control does not show is that rotation is unimportant. The contrast
moves by $7.1\%$, which exceeds its own error budget of $6.0\%$, and the
background density by $\SI{29.5}{\kilo\gram\per\cubic\metre}$, which
exceeds the $\pm\SI{27}{\kilo\gram\per\cubic\metre}$ quoted for it in
Section~\ref{sec:itokawa}. A rotational potential that is $2.4\%$ of the
gravitational one still carries a measurable share of the inference, and it
is twice what the choice of spin axis contributes
(Section~\ref{sec:methods}). One detail of the control is worth recording:
the improvement is \emph{higher} without rotation, $17.16\%$ against
$16.89\%$. The centrifugal term contributes dispersion that no two-region
density model can remove, so its presence lowers the ceiling on what the fit
can reach rather than raising the floor. The correct summary is that a slow
rotator is not ill-conditioned for this inversion, not that rotation may be
neglected in it.
% Run (v10_s8_tests.json): omega = 0 gives rho_R 2950, rho_rest 1828.6,
% contrast 1.6133, improvement 17.162 %, interior minimum. The conditioning
% claim is confirmed; the 7.1 % shift in contrast is reported beside it.

\subsection{What the method measures}
\label{sec:discussion:measures}

Section~\ref{sec:constrains} establishes that where this method responds
strongly, the quantity determined is the background density of a body, and
through the identity of Equation~\eqref{eq:excessidentity} the mass anomaly
that accompanies it, rather than the density of the region. It is more
accurate to describe the technique as bounding the mass asymmetry of a body,
from which a density contrast follows only once a dividing surface has been
fixed by other means.

That reframing bears on a question that has stood since 2008, though it does
not dissolve it, and an earlier draft of this subsection claimed more than
the data support. \citet{scheeres2008} inferred a dense neck,
\citet{lowry2014} a dense head, and \citet{kanamaru2019pss} found the neck
model fitting marginally better on the smooth terrains while reporting that
a three-region model is largely unconstrained. The reframing explains why
such a disagreement is possible: a mass anomaly does not specify the region
carrying it, and two structures placing comparable excess mass at comparable
distance from the centre of figure produce comparable centre-of-mass offsets
and comparable dispersion reductions. The degeneracy is the one already
identified in Section~\ref{sec:constrains}: the pair $(\densR,\VR)$ is free
along $\Ddens\VR = \mathrm{const}$ at fixed first moment, and no surface
measurement resolves it.

What must be said, however, is that the head and neck models of
Section~\ref{sec:itokawa:ambiguity} are \emph{not} an instance of that
degeneracy, and citing them as one would misdescribe our own result. They do
not place comparable excess mass: the head model returns
$\Dmass/\Mtot = 7.69\%$ and the neck model $3.47\%$, a factor of $2.2$, and
the slab's centroid lies much closer to the centre of figure, so the two
mass dipoles differ by more still. The objective distinguishes them
decisively, by $16.9\%$ against $2.3\%$ in improvement---a factor of
$7.3$---and by $5.8$ against $2.0$ on the displacement statistic
$\Sstat$. On the global surface the head model is preferred and the neck
model does not clear the resolution envelope at all. The genuine ambiguity
is between structures the method cannot separate; head and neck on Itokawa
are structures it can.

\subsection{An analytic test of the Paper~IV resolution result}
\label{sec:discussion:ellipsoid}

Paper~IV compared the sensitivity of two surface quantities to mesh
resolution, and concluded that the normalized geopotential dispersion is far
the more stable. That comparison, like every comparison of its kind, took the
finest available mesh as the reference. The premise is reasonable but it is
not proof: if both quantities were biased in the same direction, or if the
decimation were introducing a systematic of its own, mesh-to-mesh variation
would not reveal it.

A homogeneous ellipsoid removes the assumption. Its interior gravity, and
hence its surface gravity as the limit from within, has an exact closed form
\citep{macmillan1930}:
%
\begin{equation}
g_i = -2\pi G \varrho\, abc\, A_i x_i , \qquad
A_i = \int_0^\infty \frac{\mathrm{d}u}
     {(a_i^2+u)\sqrt{(a^2+u)(b^2+u)(c^2+u)}} ,
\label{eq:macmillan}
\end{equation}
%
which reduces to $GM/R^{2}$ in the sphere limit. Our evaluation of
Equation~\eqref{eq:macmillan} reproduces that limit to $2.2\times10^{-14}$
in relative terms. That figure is a check on the numerical quadrature of
$A_i$ and on nothing else; it is not a statement about the polyhedral code,
which cannot reproduce a sphere to any such precision, since an inscribed
icosphere of the facet counts used here falls short of $4\pi/3$ in volume by
$0.05\%$ (Section~\ref{sec:methods:verify}). Because the field points of
Section~\ref{sec:methods:model} are displaced inward from the facet
centroids, the interior form of Equation~\eqref{eq:macmillan} is the
appropriate reference at every point compared, and no exterior continuation
is required. The slope at any surface point follows from the analytic gravity
and the analytic normal of the ellipsoid, and the area-weighted means follow
by quadrature on the exact surface. Nothing in the reference depends on a
mesh, on the polyhedral code, or on any decimation.

We tested three ellipsoids, three spin states and five facet counts, ninety
configurations in all. The mean absolute error against the analytic value is
given in Table~\ref{tab:ellipsoid}.

\begin{table}[htbp]
\centering
\small
\caption{Mean absolute error against the analytic ellipsoid, averaged over
three shapes, three spin states and both decimators at each facet count. The
two decimators are separated in Table~\ref{tab:decimator}; the entries here
are their mean, which is why the column means reproduce the decimator
figures quoted below. The final two columns give the error as a fraction of
the effect this paper measures: the dispersion error is between three and
four orders of magnitude smaller than the $16.9\%$ improvement recovered on
Itokawa, which is the figure of merit that matters for the inversion, as
distinct from the accuracy of either quantity in its own right.}
\label{tab:ellipsoid}
\setlength{\tabcolsep}{10pt}
\begin{tabular}{rrrrr}
\toprule
Facets & Mean slope & Dispersion & Ratio
 & dispersion error $/\improve$ \\
\midrule
$2{,}000$  & $11.62\pm3.33\%$ & $0.196\pm0.023\%$ & $59\pm18$ & $1.2\times10^{-2}$ \\
$4{,}000$  & $4.74\pm1.35\%$  & $0.117\pm0.014\%$ & $40\pm13$ & $6.9\times10^{-3}$ \\
$8{,}000$  & $2.30\pm0.63\%$  & $0.049\pm0.004\%$ & $47\pm14$ & $2.9\times10^{-3}$ \\
$16{,}000$ & $0.93\pm0.25\%$  & $0.028\pm0.003\%$ & $33\pm9$  & $1.7\times10^{-3}$ \\
$32{,}000$ & $0.33\pm0.07\%$  & $0.017\pm0.002\%$ & $19\pm4$  & $1.0\times10^{-3}$ \\
\bottomrule
\end{tabular}
\end{table}
% The 49,152-facet row was attempted and abandoned. It needs a source mesh
% finer than the target, and the next subdivision gives 327,680 facets, at
% which a single configuration costs some 256 times the work of one at
% 20,480; nine configurations would have run for about a week. One
% configuration completed before the run was stopped, the triaxial ellipsoid
% without spin: slope errors of +0.15 % (quadric) and +0.19 % (clustering)
% against dispersion errors of -0.02 %, a ratio near 8. It is one shape and
% one spin state and is not comparable with the pooled rows above, so it is
% reported in the text as an indication and not added to the table.

The dispersion is the more accurate at every facet count, by a factor that
runs from $59$ at the coarsest mesh to $19$ at the finest, and the trend
continues beyond the table. At $49{,}152$ facets, the resolution at which the
Itokawa and Eros results are computed, a single configuration gives slope
errors of $+0.15\%$ for quadric collapse and $+0.19\%$ for clustering
against dispersion errors of $-0.02\%$: a ratio near $8$. That is the figure
a summary of this paper should quote, since it is the resolution the results
use, and it carries two qualifications. It is one shape and one spin state
rather than the average of nine, so it has none of the averaging the rows
above have; and the dispersion error it is formed against is quoted to one
significant figure, so the ratio is determined only to within roughly $6$ to
$12$. The direction is what matters:
the advantage of the dispersion narrows as the mesh refines, and the figures
that should be carried into any judgement about the meshes used in this
paper are the smaller ones at the foot of the table, not the $59$ at its
head. Three qualifications belong with those figures.

The advantage narrows with resolution, and the direction is systematic
rather than accidental. Fitting each column against facet count in
logarithms gives $E_{\mathrm{slope}} \propto N^{-1.26}$ and
$E_{\mathrm{disp}} \propto N^{-0.91}$, so the ratio falls as $N^{-0.35}$:
the slope converges faster. Extrapolated to the $49{,}152$-facet meshes on
which the Itokawa and Eros results are computed, the advantage is of order
$16$ to $20$ rather than $59$, and the figure quoted in an abstract should be
the one at the working resolution. The ratio column is noisy, and quantifiably so. With eighteen
configurations per row the standard errors of the mean are large: the ratio
is $59.3\pm18.4$ at $2{,}000$ facets, $46.8\pm13.5$ at $8{,}000$ and
$19.3\pm4.4$ at $32{,}000$. The departure from the monotone trend at
$8{,}000$ facets is under half a standard error and is not a feature. The
spread arises because three ellipsoid shapes and three spin states are
pooled in each row, and the error of a coarse mesh depends strongly on
both. ``$19$ to $59$'' therefore describes the range of the row means, not a
determination of how the ratio varies: only the two endpoints differ by more
than their errors, and the ordering in between is not resolved.
% Computed (n = 18 per row, both decimators pooled):
%   2000  slope 11.62+-3.33  disp 0.1958+-0.0233  ratio 59.3+-18.4
%   4000  slope  4.74+-1.35  disp 0.1175+-0.0139  ratio 40.4+-12.5
%   8000  slope  2.30+-0.63  disp 0.0492+-0.0044  ratio 46.8+-13.5
%  16000  slope  0.93+-0.25  disp 0.0281+-0.0026  ratio 33.2+- 9.4
%  32000  slope  0.33+-0.07  disp 0.0172+-0.0019  ratio 19.3+- 4.4
% The 8,000-facet ratio departs from the trend by less than half a standard
% error and is not a genuine feature.

Paper~IV inferred a factor of roughly $3$ to $14$ from mesh-to-mesh
variation, and the two numbers should not be read as measurements of the
same thing. Mesh-to-mesh variation measures \emph{stability}---how much a
quantity moves when the discretization changes---while
Table~\ref{tab:ellipsoid} measures \emph{accuracy} against a reference
computed independently of the code. A quantity can be stable and biased, or
unstable and unbiased. What the ellipsoid establishes is that Paper~IV's
ordering of the two quantities is not an artefact of taking the finest mesh
as truth, which is the assumption it could not test; it does not establish
that Paper~IV's factor was conservative, and an earlier draft that said so
was comparing across the two meanings.

The ellipsoid also bounds what it can be used for. Decimating from
$32{,}000$ to $2{,}000$ facets moves the dispersion by about $0.18$
percentage points on these smooth bodies, against the $0.9$ to $5.9\%$
Paper~IV measured over a comparable range on real ones---a factor of five to
thirty. Surface roughness, which an ellipsoid does not have, is the
difference. The mesh-resolution envelope used throughout this paper is
therefore a property of rough bodies that the analytic test cannot
reproduce, and the test validates the envelope's ordering and not its
magnitude.

The same experiment answers a related question about the choice of
decimator. Papers~III and~IV reduced their meshes by vertex clustering; this
paper uses quadric edge collapse, because clustering does not preserve
topology and cannot reach small facet targets
(Section~\ref{sec:methods:shapes}). Against the analytic reference the two
are equivalent in accuracy: the paired differences over the ninety
configurations are $+0.12\pm0.20\%$ on the slope and
$+0.019\pm0.012\%$ on the dispersion, neither of them significant, so the
apparent lead of clustering on the dispersion is not established
(Table~\ref{tab:decimator}). The numbers reported in Papers~III and~IV are
therefore not in question, and the change of decimator in this paper is made
for topology rather than accuracy.

\begin{table}[htbp]
\centering
\small
\caption{The two decimators against the analytic ellipsoid, averaged over
all ninety configurations. The differences are $3\%$ on the slope and $29\%$
on the dispersion, both in favour of clustering; whether either is
resolved requires the standard errors, which are not yet reported.}
\label{tab:decimator}
\setlength{\tabcolsep}{12pt}
\begin{tabular}{lrr}
\toprule
Decimator & Mean slope error & Dispersion error \\
\midrule
Vertex clustering     & $3.93 \pm 1.22\%$ & $0.072 \pm 0.009\%$ \\
Quadric edge collapse & $4.05 \pm 1.16\%$ & $0.091 \pm 0.016\%$ \\
\midrule
Paired difference     & $+0.12 \pm 0.20\%$ & $+0.019 \pm 0.012\%$ \\
\bottomrule
\end{tabular}
\end{table}
% Computed (n = 45 each): clustering slope 3.925+-1.223, dispersion
% 0.0722+-0.0088; quadric slope 4.049+-1.156, dispersion 0.0909+-0.0155.
% The configurations pair exactly, so the paired differences are the sharper
% test: quadric minus clustering is +0.0187+-0.0115 on the dispersion, 1.6
% standard errors, and +0.124+-0.200 on the slope, 0.6. Neither is
% significant, which is what makes "equivalent" defensible.

One observation from the same comparison is not explained. On the four real
bodies of the Paper~IV resolution study, the mean slope under clustering
rises monotonically with facet count on every one, by $+21.9\%$, $+14.2\%$,
$+5.8\%$ and $+1.8\%$, while under quadric collapse there is no monotonic
trend. On the smooth ellipsoid neither method drifts: both overestimate the
slope on coarse meshes by the same amount and converge smoothly. Surface
roughness is the natural explanation, but no analytic reference exists for a
rough body and we make no claim about which decimation is closer to the
truth there.
% Reconciled: 1.8-25.6 % is what Paper IV published, from its own mesh
% ladder. Recomputing the same four bodies for this paper gives 1.8-21.9 %
% with vertex clustering and 3.1-6.2 % with quadric collapse. The lower ends
% agree because Gaspra sets both; the upper end differs because the ladders
% differ. The text below now quotes the recomputed range and attributes the
% published one to Paper IV.

\subsection{Relation to Paper~IV}
\label{sec:discussion:paper4}

Paper~IV showed that the normalized geopotential dispersion is far less
sensitive to mesh resolution than the area-weighted mean slope. That is a
necessary condition for the present inference but not a sufficient one: the
minimum of a stable function need not itself be stable.
Sections~\ref{sec:itokawa:mesh} and~\ref{sec:nulls:arrokoth} supply the
missing step, and the form in which they supply it should be stated
precisely, because it is weaker than the bare phrase ``no change'' suggests.

Across the three Itokawa meshes that share a dividing plane the recovered
$\denshead$ is unchanged, which---since no sub-grid refinement is
applied---bounds the variation at half a scan step,
$\SI{12.5}{\kilo\gram\per\cubic\metre}$ or $0.45\%$, rather than resolving
it. On the coarsest mesh the recovered density is not the informative
quantity at all: the detector moves the plane by $\SI{9.4}{\metre}$, which
at the local gradient displaces $\denshead$ by about $3.4\%$ and $\densrest$
by $0.14\%$. Resolution sensitivity enters this calculation through the
plane detector rather than through the objective function, and it enters
asymmetrically, mattering for the quantity the method does not constrain and
not for the one it does. On Arrokoth the recovered density is displaced by no
more than one scan step across three meshes. The two claims---that the
objective is stable, and that its minimizer is---are separate, and both now
have evidence, of different strengths.

One correction to Paper~IV should be recorded here. It attributes the
relaxation reasoning to \citet{kanamaru2019}, who in fact cite
\citet{richardson2014} for it in their own introduction. The argument, the
technique of minimizing potential variance, and the quantity Paper~IV calls
the normalized geopotential dispersion are all due to Richardson and
Bowling; Kanamaru and Sasaki extended the technique from a single mean
density to an interior distribution. Paper~IV's reference list contains no
citation to \citet{richardson2014}, and it should have.

\subsection{One strong signal in five bodies}
\label{sec:discussion:rate}

Whether one strong signal in five bodies reflects the rarity of large
interior contrasts or the sensitivity floor of the method is not settled by
this work, and we do not claim to settle it. The four weak-signal cases are
individually explicable. Eros is measured to be nearly uniform
\citep{yeomans2000,scheeres2020}. Bennu is heterogeneous but radially so
\citep{scheeres2020}, in a geometry a plane cut cannot see.
\citet{jorda2016} state that 67P's centre-of-mass position cannot be
explained by different but individually homogeneous lobe densities, which is
precisely the model class fitted here. Arrokoth has no measured mass, and
its answer sits at the threshold of a sign reversal across the plausible
range of assumptions.

What is missing is a calibrated detection limit: the contrast a body of
given size and spin must carry before this method would see it above the
mesh-resolution sensitivity envelope. The synthetic experiment of
Section~\ref{sec:constrains:synthetic} was designed for a different purpose
and does not supply one, since it varies geometry rather than injecting a
contrast of known size.

Establishing that limit is harder than it appears. Planting a contrast in a
body and leaving its shape alone tests nothing, because the objective
function is built from the shape and the trial density and never reads a
field computed beforehand: the inversion would return whatever the shape
alone would have given. The premise asserts that the surface has relaxed
onto an equipotential of the true interior, so a body with a genuine
contrast is not the same shape as a uniform one, and a ground-truth test
must build the shape the premise describes before inverting it.

We attempted this and do not report the outcome, because the construction did
not converge. Relaxing a synthetic surface toward an equipotential of a
planted two-region interior, by displacing vertices in proportion to the
local departure of the geopotential from its area-weighted mean, reduces the
dispersion but does not settle: on a planted contrast of $1.5$ the residual
fell from $0.066$ to $0.030$ over $140$ iterations and then rose again, and
the relaxed body was no longer a contact binary, the plane holding a given
volume fraction having moved by six times its original offset. Recovered
contrasts from a surface that has not relaxed are not a test of the premise,
and their correlation with the residual---$-0.84$ across the seven rungs of the
planted ladder, at contrasts of $1.00$, $1.05$, $1.10$, $1.25$, $1.50$,
$1.75$ and $2.00$---indicates that the numbers measure the failure of the
relaxation rather than the behaviour of the inversion.
% Seven rungs, listed in the text; the correlation is over n = 7.

The mode of failure is itself informative and points to what a working
scheme would have to do. The plane holding a fixed volume fraction moved as
the surface relaxed, so the region and the shape co-evolve: relaxing toward
an equipotential of a \emph{fixed} interior is not a well-posed target once
the interior is defined by a surface that is itself moving. A fixed point
requires solving for the shape and the dividing surface together, which is
the same coupling that Section~\ref{sec:plane} reports from the other
direction---the objective does not determine its own plane. A stable scheme,
displacing along surface normals with a step controlled to reduce the
dispersion monotonically and with the volume distribution constrained, is a
piece of work in its own right. We regard it as the most useful extension of
the present study, and as a paper rather than an appendix.

A second extension needs no new machinery at all. Every body examined here
was chosen because its mass is already known, or, in Arrokoth's case,
because the consequences of not knowing it are instructive. The reverse case
is available: a body with a resolved shape model and a spin state but no
mass, which a spacecraft will visit and weigh. For such a target the
inversion can be run now as a function of the assumed bulk density, giving
a curve of contrast against assumed bulk density, as in
Section~\ref{sec:nulls:arrokoth}, rather than a number, together with
the amplification factor and the saddle prominence that say how sharply the
answer depends on the assumption. When the mass arrives it selects a point
on a curve published beforehand. That is the one form of external test this
method has not had: the comparisons with \citet{lowry2014} and
\citet{kanamaru2019} in Section~\ref{sec:itokawa} were made against
measurements that already existed when we made them.

\subsection{On verification practice}
\label{sec:discussion:verification}

This series has returned repeatedly to the question of how numerical results
in this area should be checked, and this paper adds the sharpest instance.
The cone construction of Section~\ref{sec:methods:cone} passed a
volume-closure test at $2.1\times10^{-14}$ in relative terms while enclosing
a solid wrong by a factor of $1.6$, because the test compared the mesh
against the same tetrahedra that built it. The test is not vacuous---it would
catch an orientation error or a gross failure of the tetrahedral sum---but it
is blind by construction to the one property that was wrong, which is which
solid the mesh encloses. The predicate test that does distinguish the two
solids costs three lines of code.

The other failures of Section~\ref{sec:failures} share a property: each was
invisible in the configuration that mattered and appeared only when a
parameter was varied that the analysis did not require. The NaN failure was
total at one plane position and absent at another $\SI{30}{\metre}$ away.
The edge minima appeared only when a scan was widened. The misplaced
dividing plane was found by comparing two selection rules on the same
cross-sectional profile---the saddle criterion of
Equation~\eqref{eq:prominence} against the absolute interior minimum---and
noticing that they disagreed by $\SI{15}{\kilo\metre}$; a study committed to
either rule alone would not have discovered it. A study reporting one run at
one plane on one mesh would have met none of these and would have reported
plausible numbers throughout. In this method, parameter sweeps are not
supplementary material; they are the only way the failure modes become
visible.

The buried facets of Section~\ref{sec:failures:buried} make a sharper
version of the same point. The manifold criterion was added precisely to
catch what a volume comparison cannot, and it passes them without complaint,
because a mesh in two closed pieces is topologically sound whatever their
spatial arrangement. Each new check inspects the property its author had in
mind and no other, and on this occasion the omission would have carried a
qualitative conclusion: a sign reversal that is not there.

The analytic ellipsoid of Section~\ref{sec:discussion:ellipsoid} illustrates
the complementary point. A reference computed independently of the code under
test can settle in one experiment what mesh-to-mesh comparison can only
suggest, and it can overturn a conclusion drawn from such comparison: the
provisional reading of the decimator evidence, that quadric collapse is the
more accurate, did not survive contact with it.

\subsection{Limitations}
\label{sec:discussion:limits}

The model admits two uniform regions divided by a plane perpendicular to a
principal axis. Real interiors need not take that form, and for two of the
bodies here it is known that they do not: \citet{jorda2016} exclude it for
67P, and \citet{scheeres2020} report a radial structure on Bennu that the
parameterisation cannot represent. \citet{kanamaru2019pss} show that
relaxing the model to three regions on Itokawa leaves it largely
unconstrained, so enrichment is not a straightforward remedy.

The relaxation premise itself is unverified on every body examined, and
fails its first stated condition on the one that yields a strong signal
(Section~\ref{sec:discussion:premise}). Section~\ref{sec:discussion:conditioning}
argues that this condition governs the conditioning of a different inverse
problem and does not transfer to a fixed-mass two-region inversion, but that
argument is structural: it removes a reason to expect failure and does not
establish that the surface has relaxed.

The dividing plane rests on judgement informed by two external criteria
rather than on anything measured within the method
(Section~\ref{sec:plane:external}), and the numerical diagnostics respond to
a misplaced one unevenly (Section~\ref{sec:failures:plane}). Of the four
tested, the amplification factor fails outright, ranking the misplaced plane
as the strongest signal in the study; the displacement statistic ranks it
correctly at $\Sstat = 2.86$ against $5.8$ but places it above all four
genuine nulls; the envelope comparison rejects it by a margin of $6\%$, and
only if the body-specific envelope value is above the $5.526\%$ improvement
it returns; and the decisive argument is the provenance of the plane, which
is not available from the output.

The recovered densities are quantized at the $\SI{25}{\kilo\gram\per
\cubic\metre}$ scan step, with no sub-grid refinement
(Section~\ref{sec:methods:minimisation}). The consequences are visible
throughout---two planes $\SI{0.8}{\metre}$ apart returning the same
$\denshead$, two non-monotone entries in the $\dCOM$ column, and an
improvement that can come out marginally negative where the grid does not
contain the uniform solution---and one of them, reported in
Section~\ref{sec:failures:small}, locates the true minimizer at a plane where
the grid reports a value $\SI{12}{\kilo\gram\per\cubic\metre}$ away. Solving
the stationarity condition of Equation~\eqref{eq:affine} in closed form would
remove the grid from the error budget entirely, since $\potvarn^{2}$ is a
ratio of quadratics in $\densR$ and its stationarity condition is linear with
a single root. We have not done so, and it is the clearest improvement
available to a future implementation.

The synthetic calibration of the background-density invariance rests on four
bodies carrying a measurable exponent, three of them clustered in signal
strength, so the direction of the effect is supported while its functional
form is not determined; and because those bodies are uniform, the
identification of a strong inversion response with a physical mass anomaly
remains an interpretation (Section~\ref{sec:constrains:limits}).

One shape model required a correction that others did not. The Arrokoth
model is supplied as two unjoined shells with $0.943\%$ of the surface area
lying inside the opposite shell, and excluding those facets changes the
recovered contrast at the detected plane, under the head model, from
$1.062$ to $0.9999$ at the lowest assumed bulk density
(Section~\ref{sec:nulls:arrokoth}). No other model in the study required
this, but it is not a defect peculiar to Arrokoth: any shape model delivered
as more than one component may have the same property, and the test costs one
ray-casting pass.

Finally, we have not obtained \citet{motooka2012}, which
\citet{kanamaru2019} describe as an earlier attempt to fit the global shape
of Itokawa to a single equipotential surface, and which may therefore be
closer to the approach taken here than any other prior work. It is cited as
described by Kanamaru and Sasaki, and no claim is made about its contents.

Section~\ref{sec:guidelines} sets out what these limitations imply for
anyone applying the method.

%=======================================================================
%  Paper V — Section 8b / 9 (Practical guidelines)   FINAL
%
%  Requires the macros listed at the head of Section 2.
%
%  PRINCIPAL ADDITIONS IN THIS REVISION ------------------------------
%  (a) S is shown to equal Delta M / delta(Delta M) identically, so it
%      is the signal-to-noise ratio of the excess mass and is invariant
%      under exchanging the labels "region" and "complement". A and the
%      raw curvature interval are not. This resolves the concern raised
%      in Section 7.5 that S might be a coordinate artefact, and
%      explains why A fails there: A is label-dependent.
%  (b) The reporting rule is restated in terms of the LARGER piece.
%      From the mass constraint, |d rho_rest| / |d rho_R| = f/(1-f), so
%      the better-determined density is always that of the larger
%      piece, with the threshold exactly at f = 1/2. On Arrokoth the
%      region holds 66.3 %, so "report rho_rest" as previously worded
%      points at the worse-determined of the two by a factor 1.83.
%  (c) The three-part detection criterion of Section 4.4 and the
%      diagnostic-versus-failure table are brought in, since no single
%      number separates a result from an artefact.
%
%  Superseded statements:
%      "independent parameters of the same fit"
%              -> affinely related; one free parameter (Sections 2.2, 5.1)
%      "a fractional range in Delta M is SMALLER than in rho_rest by
%       rho_rest V_tot / Delta M"
%              -> LARGER, by that factor; Eq. (excessamplify)
%      "7.5 kg/m3 per metre in the linear regime"
%              -> 10.0 at the reference plane, 55.0 at the outer planes;
%                 and the quantity to quote is rho_rest, at 0.28
%      "place the smaller of the two pieces in the region" (Section 7.3)
%              -> relabelling buys no precision; it determines which
%                 density the reporting rule names
%      "the case for a real anomaly rests on Lowry et al."
%              -> that determination is 1.75 sigma from zero and its
%                 interval contains every plane of the sweep
%      "0.94 % of the surface area"  -> 0.943 %
%
%  All values confirmed. Resolved in this pass: the 1.35 span threshold is
%  replaced by reporting span and standard error, consistently with Section
%  5.2; the envelope row now carries the Eros-specific value, 6.03 % with
%  quadric decimation and 5.40 % with clustering, and the note records that
%  the entry turns into a cross under the second.
%=======================================================================

\section{Practical guidelines}
\label{sec:guidelines}

The results above are mostly negative or qualified, which makes them
awkward to act on. This section states what they imply for anyone applying
the method to a new body. Nothing here is new evidence with one exception,
noted where it arises: the relation between the displacement statistic and
the excess mass in Section~\ref{sec:guidelines:report} is an identity we
had not previously drawn out, and it changes which quantity we recommend
reporting. The rest is the preceding sections read as instructions.

\subsection{When the method is worth applying}
\label{sec:guidelines:when}

Four conditions should hold before the result will mean much.

\emph{The mass must be measured.} On a body without one, the assumed bulk
density can carry the recovered contrast to the threshold of a sign
reversal, as it does on Arrokoth within the plausible range
(Section~\ref{sec:nulls:arrokoth}), where the recovered contrast runs from
$0.9999$ at $\SI{250}{\kilo\gram\per\cubic\metre}$ to $0.9159$ at $500$ and
the plane sweep at fixed assumed mass carries $\Dmass$ from $+2.5\%$ to
$-11.6\%$ of the total. A result that depends on an assumption in that way
is not a measurement. Where no mass exists the method may still be run, but
the honest report is a range that includes both senses.

\emph{The body should have a dividing surface that something other than the
objective function can locate.} The dispersion does not constrain the plane
(Section~\ref{sec:plane}), so a saddle in the cross-sectional profile, or
some comparable geomorphological feature, has to supply it. Bennu, which has
no saddle, is the case where this fails: a plane can be imposed, but nothing
justifies its position.

\emph{The relaxation premise should be plausible for the particular body.}
Section~\ref{sec:discussion:premise} assesses it for the five studied here,
and finds it strong on Eros, weak on 67P and Arrokoth, uncertain on Bennu,
whose surface may still be in motion, and failing its first stated condition
on Itokawa. One qualification belongs with that last item, because it is
easily over-read. Section~\ref{sec:discussion:conditioning} shows that the
first condition---that the rotational term be a significant fraction of the
gravitational one---is the condition under which the \emph{bulk-density}
inversion of \citet{richardson2014} is not identically flat, and that it has
no counterpart in a fixed-mass two-region inversion, where the mass
constraint fixes the scale instead. A slow rotator is therefore not
ill-conditioned for the problem solved here. What remains true, and is the
reason the condition still appears in this list, is that the premise itself
is a statement about whether the surface has relaxed, and the empirical
evidence that it has is weakest where the rotation is slowest. A body
resurfaced by activity, or one whose surface is still in motion, is not a
good candidate whatever its spin.

\emph{The anticipated structure must be one a plane can represent.} A radial
or axisymmetric arrangement, of the kind \citet{scheeres2020} report for
Bennu, is invisible to this model however well the other conditions hold.

\subsection{What to run}
\label{sec:guidelines:what}

A single inversion at a single plane on a single mesh is not enough to know
whether its output means anything. The following are cheap, and each answers
a question the others cannot.

\emph{Vary the mesh resolution, and report both the density and the plane.}
Two or three meshes establish whether the recovered density is stable, which
the stability of the dispersion does not by itself guarantee
(Section~\ref{sec:itokawa:mesh}). Run this twice. Holding the plane fixed
isolates the sensitivity of the objective function; letting the detector
choose it on each mesh measures the sensitivity of the plane, and on Itokawa
the second is much the larger of the two---the detected plane moves by
$\SI{9.4}{\metre}$ between the finest and coarsest meshes, displacing
$\denshead$ by about $3.4\%$ against the $0.45\%$ to which the fixed-plane
comparison bounds it. Report the detected plane on every mesh, not the
density alone; a study that fixed the plane externally would not see this
error source at all.

\emph{Sweep the dividing plane, selecting by volume fraction.} Select planes
by volume fraction rather than by position, so that the sweep spans a
comparable range of region volume on bodies of different shape, and report
the scaling exponent of Eq.~\eqref{eq:slope} together with the fractional
ranges of $\densR$, $\densrest$ and $\Dmass$ and the subset of planes over
which each was computed. Three conditions govern whether the exponent means
anything. The density difference must keep its sign across the sweep, since
a quantity passing through zero has no fractional range and no logarithmic
slope; the sweep must span an adequate range in $\VR$, which we judge by
reporting the span and the standard error of the fit with every exponent
rather than by imposing a threshold---the spans used here are $1.907$ on
Itokawa, $1.504$ on 67P and $1.90$ on every synthetic body, all far above
the factor of about $1.35$ below which the logarithmic fit ceases to be
constrained; and the exponent should not be presented as
independent evidence about $\densrest$, since Eq.~\eqref{eq:steepidentity}
makes $s+1$ identically the logarithmic drift of $\densbulk-\densrest$. What
the exponent adds is the distinction between a systematic drift and scatter
about a fixed value, which a coefficient of variation cannot make, and a
well-defined quantity where fractional ranges diverge. Report the
centre-of-mass offset across the same sweep as well: comparing its scaling
against the purely geometric slope of Eq.~\eqref{eq:geomslope} is what
excludes conservation of the mass dipole
(Section~\ref{sec:constrains:dipole}), and it costs nothing, the offset
being a by-product of the same solutions.
% Resolved consistently with Section 5.2: the threshold does no work, since
% no sweep reported here comes near it, and the standard error is now quoted
% with every exponent (Itokawa -1.141 +/- 0.018, 67P -2.80 +/- 0.27, and the
% four synthetic exponents in Table 5.2).

\emph{Vary the assumed mass.} Three or four runs over a factor of a few show
whether the answer is driven by the data or by the assumption
(Section~\ref{sec:constrains:mass}). Read the outcome as the position of the
resulting contrast interval relative to unity and not as its width: on the
raw change in contrast Itokawa is the \emph{more} sensitive of the two
bodies tested, and it is only per e-fold of assumed mass, and in whether the
interval reaches unity, that the two separate. We put this forward as a
candidate diagnostic rather than an established one, since it rests on two
bodies here, and note the awkwardness that it is of use precisely where the
mass is unknown and there is nothing to calibrate it against.

\emph{Check that no minimum sits on a scan boundary, and widen the interval
until none does.} A minimisation reports the smallest value it finds, and a
minimum at the edge of the interval searched is the edge, not a result. This
guard fired twice in the present study: on the Arrokoth neck model, where a
slab of $1.0\%$ of the volume drove the density to the scan floor and
returned an improvement of $13.6\%$ that had to be discarded, and at one
Itokawa plane, where widening the interval returned an interior minimum
(Section~\ref{sec:failures:edge}). Record the interval used at every plane,
since it may have to change along a sweep, and check separately that the
upper limit is not implicitly restricting the region size: holding $\Dmass$
fixed, a ceiling $\densR^{\max}$ admits no region smaller than
$\Dmass/[(\densR^{\max}-\densrest)\Vtot]$, which on Itokawa is $2.56\%$ and
bounds the range of planes the configuration can resolve at all.

\emph{Check the geometry every time, in five layers.} The region must agree
with a closed-form solid where one exists, satisfy its defining inequality
at every vertex, close against its complement, form a closed surface with
the Euler characteristic its component count implies, and satisfy the mass
identity $\Dmass=(\densbulk-\densrest)\Vtot=\Ddens\VR$ on every solution.
The last of these is the one that operates on the output rather than on the
geometry, and it is the layer that would have caught the construction error
of Section~\ref{sec:methods:cone}, which passed a volume-closure test while
enclosing a solid wrong by a factor of $1.6$. Where the shape model has more
than one component, test facet centroids for containment in the others as
well: none of the other checks detects a facet that lies inside the body
(Sections~\ref{sec:methods:verify} and~\ref{sec:failures:buried}).

\emph{Where the signal is weak, test the field-point offset.} The inward
displacement of field points from the facet centroids shifts the surface
potential by of order $0.04\%$ of $\lvert\Ubar\rvert$ on the bodies here
(Section~\ref{sec:methods:model}), which is comparable to the dispersion
improvements returned by all four weak-signal bodies. Repeating the
inversion with the offset scaled by ten and by a tenth establishes whether a
sub-per-cent improvement is a property of the body or of the discretisation.

\subsection{What to report}
\label{sec:guidelines:report}

\emph{Report the density of the larger of the two pieces, and the excess
mass, not the region density as such.} This is a correction to the rule
stated in earlier drafts, which was to report the background density
$\densrest$, and the correction matters on one of the five bodies here.
Differentiating the mass constraint~\eqref{eq:massbalance} at fixed total
mass gives
%
\begin{equation}
\frac{\lvert\delta\densrest\rvert}{\lvert\delta\densR\rvert}
 = \frac{\VR}{\Vtot-\VR} = \frac{f}{1-f} ,
\label{eq:whichpiece}
\end{equation}
%
so the density of the larger piece is always the better determined, in
absolute terms, and the threshold is exactly $f = \tfrac{1}{2}$. On Itokawa,
Eros and 67P the region is the smaller piece and $\densrest$ is accordingly
the constrained quantity, which is the case the earlier wording had in mind.
On Arrokoth the detector places $66.3\%$ of the volume in the region, so
$\densrest$ is the density of the \emph{smaller} piece and is the worse
determined of the two by a factor of $1.83$ in fractional terms. Reporting
it there, as Table~\ref{tab:nulls} does, points at the wrong number. Which
piece is called the region is a labelling choice and carries no physics: it
does not change the uncertainty on either density, only which of them the
reporting rule names.

Report $\Dmass$ alongside, since Eq.~\eqref{eq:excessidentity} makes it the
same statement and it is the quantity that survives the relabelling. The
region density is a derived quantity depending on a dividing surface the
investigator chose (Section~\ref{sec:constrains:claim}), and a region
density quoted without its plane is not reproducible.

\emph{Quote the amplification factor with any stability claim, and the plane
at which it was evaluated.} A fractional range in $\densrest$ corresponds to
a fractional range in $\Dmass$ that is \emph{larger} by
$\ampfac = \densrest\Vtot/\Dmass$, not smaller---on Itokawa,
$\CV(\Dmass)=2.95\%$ against $\CV(\densrest)=0.251\%$, a ratio of $11.76$
which reproduces the mean $\ampfac$ of $11.77$ over the same planes. Any
comparison of how well two quantities are determined must divide by it
first. When comparing $\densR$ against $\densrest$, note that the two are
comparable because they are densities of the same model in the same units,
and \emph{not} because they are independent: the mass constraint leaves
exactly one free parameter and relates them affinely as
$\densrest = a - b\,\densR$ (Section~\ref{sec:methods:identity}). Quote
$\ampfac$ with its plane, since Eq.~\eqref{eq:ampsensitivity} makes it move
by $\pm5.6\%$ across a sweep over which $\densrest$ moves by $\pm0.434\%$,
and note that a large $\ampfac$ is a warning rather than a precision: a body
returning $\ampfac\simeq1200$, as Bennu does, is one on which the model has
found essentially no excess mass.

\emph{State the component count of the shape model, and whether any facets
were excluded.} A model delivered as several unjoined shells may have part
of one inside another; on Arrokoth this affected $0.943\%$ of the surface
area and changed a qualitative conclusion
(Section~\ref{sec:nulls:arrokoth}).

\emph{State the plane, the criterion that placed it, and the sensitivity of
the reported quantity to moving it.} The sensitivity must be quoted for the
quantity being reported, and the three differ by two orders of magnitude. At
the Itokawa reference plane they are
$\SI{10.0}{\kilo\gram\per\cubic\metre\per\metre}$ for $\denshead$---rising
to $\SI{55}{\kilo\gram\per\cubic\metre\per\metre}$ at the outer planes and
falling to $7.5$ over the first four, so no single figure describes the
sweep---$\SI{0.28}{\kilo\gram\per\cubic\metre\per\metre}$ for $\densrest$,
and $0.18\%$ per metre for $\Dmass$. Earlier drafts quoted the first alone,
which states the sensitivity of the quantity
Section~\ref{sec:constrains} concludes the method does not determine.

\emph{State the saddle prominence, and do not read it as a figure of merit.}
It lets a reader judge the provenance of the plane, and the values here
separate the accepted from the refused by more than two orders of magnitude,
from $0.091$ on 67P down to $0.0004$ on Bennu. Among the accepted bodies it
carries no information about the strength of the result: Arrokoth has the
highest prominence in the study, $0.490$, and returns an improvement of
$1.272\%$, while Itokawa has the third, $0.170$, and returns $16.9\%$.

\emph{Give the facet count used, not the one requested, and the decimation
method.} A slope or dispersion statistic depends on both
(Section~\ref{sec:discussion:ellipsoid}).

\emph{Bound the answer by the material, and state that the bound is
one-sided.} A recovered density above the grain density of the meteorite
analogue is not a result but a signal that the plane has advanced too far:
on Itokawa the outermost plane of the sweep exceeds the LL-chondrite grain
density by $11.2$ standard deviations
(Sections~\ref{sec:itokawa:ceiling} and~\ref{sec:plane:impossible}). Convert
the bound into a plane position through the measured gradient and carry its
uncertainty: on Itokawa the block-density bound of
$3220\pm\SI{220}{\kilo\gram\per\cubic\metre}$ becomes
$\xsplit \le \SI{194+-11}{\metre}$, and quoting it to the tenth of a metre
would overstate it. The bound excludes planes that are too far out and says
nothing about planes that are too far in, so it truncates a sweep and does
not locate one.

\subsection{How to read the outcome}
\label{sec:guidelines:read}

\emph{Require three things before calling a result a detection.} No single
number separates a result from an artefact, which is the lesson of
Section~\ref{sec:failures} and the reason for stating a criterion rather
than a threshold. The improvement must exceed the value the mesh-resolution
envelope takes on that body, not the upper edge of a range attained on some
other body; the displacement statistic $\Sstat$ must be large; and the
dividing plane must carry geometric warrant, meaning a saddle returned by a
detector at a prominence above threshold rather than a plane chosen by the
investigator. The misplaced plane of Section~\ref{sec:failures:plane}
satisfies the second and fails the first and third. Itokawa satisfies all
three.

\emph{Use $\Sstat$ rather than a curvature interval, because it is the only
one of the available diagnostics that does not depend on which piece was
called the region.} Its definition in Eq.~\eqref{eq:sstat} is a displacement
measured in units of an uncertainty, but multiplying numerator and
denominator by $\Vtot$ turns it into something more directly meaningful:
%
\begin{equation}
\Sstat = \frac{\lvert\densbulk-\densrest\rvert}{\sigma(\densrest)}
       = \frac{\lvert\densbulk-\densrest\rvert\,\Vtot}
              {\sigma(\densrest)\,\Vtot}
       = \frac{\lvert\Dmass\rvert}{\delta(\Dmass)} ,
\label{eq:ssnr}
\end{equation}
%
which is exact. $\Sstat$ is the signal-to-noise ratio of the excess mass,
and since $\Dmass$ merely changes sign when the labels ``region'' and
``complement'' are exchanged, $\Sstat$ is invariant under that exchange. The
raw curvature interval is not, and neither is $\ampfac$: at the misplaced
Eros plane the interval is $\pm15$ under one labelling and $\pm115$ under
the other, and $\ampfac$ is $-9.12$ under one and $+61.0$ under the other.
That is why $\ampfac$ ranks the misplaced plane as the strongest signal in
the study, and it is a limitation of the quantity rather than a statistical
accident: $\ampfac$ is a cheap first check on bodies whose planes are
properly placed (Section~\ref{sec:nulls:amplification}) and says nothing
about whether a plane is properly placed. Read the other way,
Eq.~\eqref{eq:ssnr} says that $\Sstat = 5.8$ on Itokawa is a statement that
the excess mass is determined to $1/5.8 = 17\%$, which is the figure the
error budget of Eq.~\eqref{eq:budget} arrives at independently.

\emph{A strong response is not by itself evidence of a physical anomaly.}
The synthetic bodies of Section~\ref{sec:constrains:synthetic} are uniform
by construction and still produce strong responses and a pinned background
density when their shape is sufficiently asymmetric. What that pinning
identifies is a regime of the objective function. On Itokawa the case for a
real anomaly rests on the independent YORP determination of
\citet{lowry2014}, not on the inversion alone---but that support should not
be asked to carry more than it can. The Lowry offset of
$\SI{21+-12}{\metre}$ is itself only $1.75$ standard deviations from zero,
so it does not exclude a homogeneous Itokawa at two sigma, and its interval
contains the centre-of-mass offset returned at \emph{every} plane of
Table~\ref{tab:itokawasweep}. It corroborates the sign and the order of
magnitude of the anomaly and discriminates neither the plane nor the value.

\emph{A weak response is not evidence of homogeneity.} It is consistent with
a uniform interior, with a structure the planar model cannot see, and with
the relaxation premise failing on that body. These are not distinguishable
from within the method (Section~\ref{sec:nulls}).

\emph{A sharp minimum is not a precise answer; propagate the interval before
reading it.} The tightest curvature bound in this study,
$\pm\SI{15}{\kilo\gram\per\cubic\metre}$, belongs to the one run known to be
wrong: a plane placed on a flank of Eros, where the complement holds only
$11.56\%$ of the volume and the mass constraint amplifies the objective on
both sides of the minimum. Carried through Eq.~\eqref{eq:whichpiece} that
interval becomes $\SI{115}{\kilo\gram\per\cubic\metre}$ on the complement
density, or $3.83\%$---the worst determination anywhere in this study, and
correspondingly $\Sstat = 2.86$ against $5.8$ on Itokawa. The apparent
precision is a property of the coordinate in which the scan was performed.

\emph{Watch for an improvement bought with almost no mass.} The ratio of the
dispersion improvement to the fractional excess mass is $14.4$ for the
Arrokoth neck model, which the boundary guard discarded, against $2.20$ for
the Itokawa head and below unity for every other case in the study
(Section~\ref{sec:failures:edge}). It separates an edge minimum from a
genuine result by a factor of $6.5$ and costs nothing to compute. Its
limitation should be stated with it: it does not flag a misplaced plane,
whose ratio is unremarkable.

\emph{Corroboration should be sought outside the method.} Itokawa is
interpretable here because a YORP measurement, resting on different
premises, gives a compatible centre-of-mass offset, and because
Section~\ref{sec:constrains:dipole} establishes that the objective function
demonstrably does not hold the mass dipole fixed, so the two lines of
evidence are independent rather than two readings of one. Where no such
check exists, the result should be reported as what the objective function
prefers, not as what the interior is.

\begin{table}[htbp]
\centering
\scriptsize
\setlength{\tabcolsep}{3pt}
\caption{Which diagnostic responds to which failure, from the cases in
Section~\ref{sec:failures}. A tick marks a diagnostic that identifies the
failure, a cross one that is silent or actively misleading, and a dash one
that does not bear on it. The table is the argument for a three-part
criterion rather than a threshold: no row catches everything, and two
diagnostics rank the misplaced plane \emph{above} the genuine result. The
last two columns are the two silent failures that no statistical diagnostic
reaches, and that only geometric checks or the provenance of the plane will
find.}
\label{tab:diagnostics}
\setlength{\tabcolsep}{0.2pt}
\begin{tabular}{lccccc}
\toprule
Diagnostic & Region too small & Complement too small & Edge minimum
 & Misplaced plane & Wrong solid \\
\midrule
Curvature interval in $\densR$
 & \checkmark\ $+22.3\%$ & $\times$\ $0.57\%$ & --- & $\times$\ $0.57\%$ & --- \\
Interval propagated to the larger piece
 & $\times$\ $1.01\%$ & \checkmark\ $3.83\%$ & --- & \checkmark\ $3.83\%$ & --- \\
Boundary-minimum guard
 & --- & --- & \checkmark & --- & --- \\
$\improve$ per unit $\lvert\Dmass\rvert/\Mtot$
 & --- & --- & \checkmark\ $14.4$ & $\times$\ $0.45$ & --- \\
$\Sstat$
 & --- & \checkmark & --- & partly, $2.86$ & --- \\
$\ampfac$
 & --- & --- & --- & $\times$\ $9.12$ & --- \\
Envelope comparison$^{a}$
 & --- & --- & --- & narrowly, $5.53\%$ & --- \\
Plane provenance
 & --- & --- & --- & \checkmark & --- \\
Vertex predicate; mass-balance audit
 & --- & --- & --- & --- & \checkmark \\
\bottomrule
\end{tabular}
\end{table}
\vspace{2pt}
\begin{minipage}{\textwidth}
\footnotesize
\textsuperscript{a}\,The Eros-specific envelope is $6.03\%$ with quadric
decimation, so the $5.53\%$ returned at the misplaced plane is $0.92$ times
it and the comparison rejects the case. With vertex clustering the envelope
is $5.40\%$, the same result is $1.02$ times it, and the entry becomes a
cross. The margin is therefore smaller than the difference between two
defensible ways of computing the envelope, which is why the row reads
``narrowly'' and why plane provenance, the row below it, is the check we
recommend relying on.
\end{minipage}

%=======================================================================
%  Paper V — Section 9 (Conclusions)   FINAL
%
%  Requires the macros listed at the head of Section 2.
%
%  Superseded values against ver10:
%      "2e-14 per cent"          -> 2.1e-14 in relative terms
%      0.11 R centroid offset    -> 0.1125 R
%      26.68 % -> 26.6 %, the value the run records (26.59 %)
%      I = 16.889 %              -> 16.9 %   (grid floor, Section 6.1)
%      C = 1.507                 -> 1.51 +/- 0.09
%      Bennu I = 0.000 %         -> < 10^-3 % (negative I observed)
%      Dx_COM = 16.5 m           -> 16.5 +/- 2.9 m; Lowry is 1.75 sigma
%      +/- 0.435 %               -> +/- 0.434 %
%      s = -1.14                 -> -1.141 +/- 0.018, eleven planes
%      correlation 0.95          -> 0.95 with n = 4, p ~ 0.03
%      rho_rest = 1855 +/- 8     -> 1858 +/- 27  (Eq. budget)
%      Delta M = 2.75e9, 7.7 %   -> (2.75 +/- 0.48)e9, 7.7 +/- 1.3 %
%      sigma to six decimals     -> five (Section 6.1)
%      turning point at 0.25 km  -> REMOVED; argument rebuilt on the
%                                   11.2 sigma compositional exclusion
%      Arrokoth lower bound 290  -> central ~235, range 155-600
%      7.5 kg/m3 per metre       -> 10.0 at the reference plane, and
%                                   0.28 for rho_rest
%      "no diagnostic detects a misplaced plane"
%                                -> three of four respond; A fails
%      2.2e-14 per cent          -> relative, and a check on the
%                                   quadrature only
%      "19 to 59 times"          -> add the working-resolution figure
%      "confirmed and conservative"
%                                -> confirms the ORDERING; stability and
%                                   accuracy are different quantities
%
%  Added, from results absent in ver10:
%      conclusion 5  exclusion of the mass dipole at 21 sigma
%      conclusion 4  decomposition of the stability ratio; S as the
%                    signal-to-noise of the excess mass
%      conclusion 9  reporting rule keyed to the larger piece
%      closing       the first Richardson condition does not transfer
%
%  All values confirmed. Resolved here: the Arrokoth 395/410 against
%  390/419.7 discrepancy (buried facets retained in the resolution series);
%  the Itokawa-specific envelope, 1.35 %, and the ratio 12.5 against it; the
%  235 kg/m3 run and the three conflicting published density constraints;
%  standard errors on the ellipsoid ratios; the completed residual table and
%  the omega = 0 control. Only the Zenodo DOI is still a placeholder.
%=======================================================================

\section{Conclusions}
\label{sec:conclusions}

This paper takes the interior density inversion of \citet{richardson2014}
and \citet{kanamaru2019} as an object of study rather than as a tool, and
asks what it determines. Five bodies were examined---Itokawa, Eros,
Bennu, 67P/Churyumov--Gerasimenko and Arrokoth---together with eight
synthetic contact binaries constructed so that the strength of the signal
could be varied while everything else was held fixed, and a homogeneous
ellipsoid whose surface gravity has a closed form. The main conclusions
are as follows.

\begin{enumerate}

%---------------------------------------------------------------
\item \emph{Regions must be built by exact plane clipping, and the
natural alternative fails silently.} Closing a region by fanning its
boundary edges to an interior point yields a watertight mesh whose volume
matches the sum of its own tetrahedra to $2.1\times10^{-14}$ in relative
terms, and which encloses a radial cone rather than the plane-cut lobe
intended. On a unit sphere cut at $x = 0.5$ the two solids differ by a
factor $1.600$ in volume and by $0.1125\,R$ in centroid; on Itokawa the
error reversed the sign of the recovered contrast. A volume-closure test
cannot detect it, because it compares the mesh against the same
tetrahedra that built it. Five layers are required, of which three can
fail in ways the others cannot: a predicate test, requiring every vertex
of the region to satisfy the inequality defining it, which the cone fails
outright; a per-component manifold test, which catches a surface torn
along the cut that every volume-based test passes; and an audit of the
identity $\Dmass = (\densbulk-\densrest)\Vtot = \Ddens\VR$ on every
solution, which operates on the output rather than on the geometry and
would have caught the cone error there, the two volumes differing by the
same factor $1.600$. The clipped construction reproduces closed-form cap
and slab volumes to $0.09\%$, limited by the polyhedral approximation of
the sphere itself, and reproduces the independently published lobe volume
of 67P, $26.6\%$ against $27\%$ \citep{jorda2016}.

%---------------------------------------------------------------
\item \emph{The density recovered by the inversion is resolution-stable
over the meshes tested, and the residual sensitivity enters through the
plane detector rather than through the objective function.} This does not
follow from the stability of the dispersion established in Paper~IV,
since the minimum of a stable function need not itself be stable, and was
therefore tested separately. Three meshes of Itokawa sharing a dividing
plane, spanning a $64\%$ increase in facet count, return $\denshead$
unchanged; because no sub-grid refinement is applied, that bounds the
variation at half a scan step, $\SI{12.5}{\kilo\gram\per\cubic\metre}$ or
$0.45\%$, rather than resolving it. Three meshes of Arrokoth at a fixed
plane, spanning $173\%$, likewise return an identical
$\SI{395}{\kilo\gram\per\cubic\metre}$; that series retains the buried
facets of Section~\ref{sec:methods:shapes}, since it compares meshes with
one another rather than reporting an absolute value, which is why its
densities are $395$ and $410$ rather than the $390$ and $419.7$ of
Table~\ref{tab:arrokothdensity}. The informative case is the
coarsest Itokawa mesh, on which the detector moves the plane by
$\SI{9.4}{\metre}$; at the local gradient that displaces $\denshead$ by
about $3.4\%$, comparable to its own curvature interval, and $\densrest$
by $0.14\%$, where it adds nothing to the budget. Resolution sensitivity
therefore matters for the quantity the method does not constrain and not
for the one it does, and a study that fixed the plane externally would
not see it at all. Three meshes are not a convergence sequence, and no
claim of convergence in the analytical sense is made.
% Resolved: the resolution series retains the buried facets and the table
% excludes them; both files now say so (Sections 4.7 and 9).

%---------------------------------------------------------------
\item \emph{Of five bodies, only Itokawa yields a dispersion improvement
above the mesh-resolution sensitivity envelope of Paper~IV.} At the
detected neck the improvement is $16.9\%$ with
$\denshead/\densrest = 1.51 \pm 0.09$, a factor $2.9$ above the upper
edge of the $0.9$--$5.9\%$ envelope, which measures discretization
sensitivity rather than statistical noise. That upper edge is attained on
a different body. Against the value the envelope takes on Itokawa itself,
$1.35\%$ over the ladder from $2{,}000$ to $42{,}000$ facets with quadric
decimation, the improvement is $12.5$ times larger, and that is the
meaningful comparison.
Eros, Bennu, 67P and Arrokoth return $0.418\%$, below $10^{-3}\%$,
$0.574\%$ and $1.272\%$, all within or below that envelope and all
returning contrasts consistent with unity once the curvature interval is
propagated through the mass constraint---Eros, the nearest to a
detection, at $1.018 \pm 0.026$, or $0.7$ intervals from unity. Each is
explicable: Eros is measured to be nearly uniform, Bennu's reported
heterogeneity is radial and invisible to a plane cut, 67P's
centre-of-mass position is stated by \citet{jorda2016} to be inexplicable
by two uniform lobes, and Arrokoth has no measured mass. At the dividing
plane $x = \SI{0.150}{\kilo\metre}$ used by \citet{kanamaru2019} we
recover a contrast of $1.467$ against their $1.460$; the agreement tests
implementations rather than physics, since both minimize a functional of
the surface geopotential, and our own interval on that contrast is
$\pm0.09$, two orders of magnitude wider than the difference. The same
plane gives a centre-of-mass offset of $16.5 \pm \SI{2.9}{\metre}$,
consistent with the $\SI{21+-12}{\metre}$ of \citet{lowry2014}.

%---------------------------------------------------------------
\item \emph{Where the method responds strongly, the quantity it
constrains is the background density of the body---and equivalently the
mass anomaly---not the density of the region; and that holds in the
strong-signal regime rather than of the method as such.} The mass
constraint makes $\Dmass = \Ddens\VR$ equal to $\Mtot-\densrest\Vtot$
exactly, so a statement that one is invariant is the same statement about
the other; these are one result and not two, and the pair
$(\densR,\VR)$ remains degenerate along $\Ddens\VR = \mathrm{const}$.
Across the thirteen-plane Itokawa sweep the head density varies by
$\pm38.3\%$ while the background density varies by $\pm0.434\%$, and in
coefficients of variation the ratio is $82$. Most of that is not a
result. Equation~\eqref{eq:excessamplify} fixes a factor of $11.8$ of it as
algebra, present whatever the objective does, and the subtraction of a
nearly constant background offsets a further $2.4$, which is also algebra.
What remains, the factor of $16.9$ by which $\Dmass$ is steadier than
$\Ddens$, measures compensation between contrast and region volume; over
the eleven planes the composition permits the figures are $32.1$, $11.78$,
$2.77$ and $7.52$, and over the nine below the softer bound $27.3$, $11.88$,
$2.92$ and $6.71$. It is the compensation that is the result. On 67P and
Arrokoth the
order of stability reverses, and the controlled synthetic experiment
resolves the discrepancy through the exponent
$s = \mathrm{d}\log\lvert\Ddens\rvert/\mathrm{d}\log\VR$, which takes the
value $-1$ when the background density is held fixed. Bodies with a weak
response show a density difference that changes sign, so no exponent
exists; those with a moderate response give $s \simeq -2$; only the
strongest gives $-1.22$, against Itokawa's $-1.141 \pm 0.018$ over the
eleven planes that span the same factor $1.907$ in region volume as the
synthetic family. The correlation between signal strength and exponent is
$0.95$, but over the four synthetic bodies that have an exponent, three of
them clustered in signal strength, so the direction of the effect is
supported and its functional form is not determined. Because the
synthetic bodies are uniform by construction, this establishes a property
of the objective function rather than a physical inference. For Itokawa
the constrained quantities are
$\densrest = 1858 \pm \SI{27}{\kilo\gram\per\cubic\metre}$ and
$\Dmass = (2.75 \pm 0.48)\times10^{9}$~kg, or $7.7 \pm 1.3\%$ of the
total, the interval in both cases being dominated by the curvature
sensitivity interval and not by the plane-placement scatter, which is
three times smaller and is not an uncertainty. Equivalently, the
displacement statistic $\Sstat$ of Equation~\eqref{eq:sstat} equals
$\Dmass/\delta(\Dmass)$ identically, so its value of $5.8$ on Itokawa is
the statement that the excess mass is determined to $17\%$.

%---------------------------------------------------------------
\item \emph{The conserved quantity is the excess mass and not the mass
dipole, and the two are separated by twenty-one standard errors.} The
alternative hypothesis---that the objective holds fixed
$\Dipole = \Dmass\,(x_R-x_{\mathrm{cf}})$, which is the quantity the YORP
determination of \citet{lowry2014} constrains---makes a different and
quantitatively separated prediction, because the geometric factor is not
constant along a sweep. Over the eleven admissible planes
$\mathrm{d}\log(x_R-x_{\mathrm{cf}})/\mathrm{d}\log\VR = -0.244$, fixed by
the shape model alone, so dipole conservation requires $s = -0.756$; the
measured $-1.141 \pm 0.018$ lies on the opposite side of $-1$. Stated
directly, $\mathrm{d}\log\Dipole/\mathrm{d}\log\VR = -0.385 \pm 0.018$,
twenty-one standard errors from the zero that conservation requires, and
the conclusion does not depend on which subset of planes is used: the same
calculation gives $10.9\sigma$ over all thirteen and $16.8\sigma$ over the
nine. This matters beyond the internal question, because it makes the
agreement with \citet{lowry2014} a genuinely independent check rather than
two readings of one quantity.

%---------------------------------------------------------------
\item \emph{Within the two-region planar model, dispersion minimization
does not constrain its own dividing plane, and the demonstration requires
no extrapolation.} The dispersion at the minimum falls monotonically as
the plane advances, from $0.07549$ to $0.06373$ across the thirteen
Itokawa planes, an improvement rising from $12.9\%$ to $26.4\%$ without a
single reversal. At the outermost plane the recovered head density is
$\SI{5000}{\kilo\gram\per\cubic\metre}$, which exceeds the LL-chondrite
grain density of $3540 \pm \SI{130}{\kilo\gram\per\cubic\metre}$---a
zero-porosity ceiling---by $11.2$ standard deviations, and the empirical
block-density bound of $3220 \pm \SI{220}{\kilo\gram\per\cubic\metre}$ by
$8.1$. Since the improvement is strictly increasing and the admissible set
is an interval bounded above, the improvement attains its maximum over
that set at a boundary fixed by the meteorite collection rather than by
the minimization, and no statement about the behaviour of $\potvarn$
beyond the sweep is required. Converted through the measured gradient the
two bounds give $\xsplit \le \SI{0.2069+-0.0047}{\kilo\metre}$ and
$\le \SI{0.1935+-0.0110}{\kilo\metre}$. The reason follows from
conclusion~4: with $\densrest$ held nearly fixed,
$\densR = \densrest + \Dmass/\VR$ diverges as $\VR \to 0$, and nothing in
the objective opposes the divergence. Cubic and quartic fits place a
turning point near $\SI{0.254}{\kilo\metre}$, but we do not rest anything
on it: the figure is an extrapolation beyond the last plane computed, the
measured gradient extrapolates instead to $0.262$--$\SI{0.278}{\kilo
\metre}$, the position coincides with the plane at which the scan ceiling
begins to truncate the recovered contrast, and it is not established that
an interior turning point exists at all, the objective having a finite
limit as the region shrinks to a point mass at the tip.

%---------------------------------------------------------------
\item \emph{The plane must come from outside the method, and the two
external criteria available are not equally strong.} The neck detector
returns $\xsplit = \SI{160.8}{\metre}$ from the cross-sectional profile
alone, and the material bound excludes planes beyond
$\SI{194+-11}{\metre}$; the detected neck falls $\SI{33}{\metre}$ inside
the bound, about three times the bound's own uncertainty. That the two do
not contradict one another is reassurance and not corroboration, because
the compositional criterion is one-sided: it excludes planes that are too
far out and says nothing about planes that are too far in, and every plane
below $\SI{194}{\metre}$ satisfies it, which is $58\%$ of the range swept.
The objective's own preference among admissible planes is real but not one
the method can defend: it prefers the block-density boundary to the
detected neck by $4.6$ percentage points of improvement, which lies inside
the mesh-resolution envelope. The saddle prominence, which should be
published so that a reader can judge the provenance of the plane,
separates accepted from refused bodies by more than two orders of
magnitude but carries no information about the strength of the result:
Arrokoth has the highest prominence in the study, $0.490$, and returns
$1.272\%$, while Itokawa has the third, $0.170$, and returns $16.9\%$.

%---------------------------------------------------------------
\item \emph{On a body without a measured mass, the assumed bulk density
can reverse the sign of the inferred contrast, not merely change its
scale.} For Arrokoth the contrast runs from $0.9999$ at an assumed
$\SI{250}{\kilo\gram\per\cubic\metre}$ to $0.9159$ at $500$, and at fixed
assumed mass the plane sweep carries the excess mass from $+2.5\%$ to
$-11.6\%$ of the total, passing through zero at a region fraction of
$61.2\%$---five percentage points from the plane the detector selects. The
method cannot say which lobe is denser. A fourth run at the central published estimate of \citet{keane2022},
$\SI{235}{\kilo\gram\per\cubic\metre}$, returns $1.066$: the sense of the
answer reverses between $235$ and $250$, inside the published one-sigma
range of $155$ to $600$ and below the bound of $290$ that
\citet{spencer2020} derive from the requirement that mutual gravity
exceed the centrifugal acceleration. What the inversion reports about this
body is close to a restatement of what was assumed about it. The converse suggests a diagnostic: on Itokawa a fourfold change
in assumed mass moves the contrast by $11\%$ and never near unity. The
outcome must be read as the position of the resulting interval relative to
unity rather than as its width, since on the raw change Itokawa is the
more sensitive of the two bodies tested and the separation appears only
per e-fold of assumed mass. We put this forward as a candidate rather than
an established test: it rests on two bodies that differ in several
respects besides signal strength, and it is of use precisely where the
mass is unknown and there is nothing to calibrate it against.
% Resolved: 235 with a 1-sigma range 155-600 is Keane et al. (2022); the
% 290 lower bound is Spencer et al. (2020) under the no-tension condition;
% McKinnon et al. (2020) give 250-500. All three are now cited in Sections
% 2.9 and 4.7, and the 235 run has been performed.

%---------------------------------------------------------------
\item \emph{Seven failure modes produce plausible-looking numbers, and no
single diagnostic separates a result from an artefact.} A region too small
flattens the minimum until the density interval reaches $22\%$ of its own
value; a complement too small sharpens it, and the tightest curvature
bound in this study, $\pm\SI{15}{\kilo\gram\per\cubic\metre}$, belongs to
the one run known to be wrong---an artefact of the coordinate scanned,
since propagating it through the mass constraint at a region fraction of
$88.4\%$ gives $\pm\SI{115}{\kilo\gram\per\cubic\metre}$, or $3.83\%$, the
worst determination anywhere in the study. A minimum at a scan edge
reported an apparent $13.6\%$ improvement from a value pinned to the scan
floor. Facets lying inside a shape model delivered as two unjoined shells,
$0.943\%$ of the surface area, produced a crossing of unity on Arrokoth
that is not there. And a plane misplaced onto a flank of Eros returned a
contrast of $0.876 \pm 0.039$, excluding unity by more than three
intervals, with every geometric check passing. Four numerical diagnostics
were tested against that last case and they do not agree: the
amplification factor fails outright, ranking the artefact as the strongest
signal in the study at $\lvert\ampfac\rvert = 9.12$ against Itokawa's
$12.0$, because $\ampfac$ depends on which piece is called the region; the
displacement statistic ranks it correctly at $\Sstat = 2.86$ against
$5.8$, and is the only diagnostic invariant under that relabelling, but it
places the artefact above all four genuine nulls; the envelope comparison
rejects it by a margin of $6\%$; and the decisive argument is the
provenance of the plane, which is not available from the output. A
validity window on region volume fraction was considered and rejected, the
degradation having no threshold in the data. What replaces it is
procedural: place the plane at a true saddle, publish its prominence,
propagate the curvature interval through the mass constraint before
reading it, report the sensitivity of the reported quantity to moving the
plane, and test facet centroids for containment whenever a model has more
than one component.

%---------------------------------------------------------------
\item \emph{The quantity to report is the density of the larger of the two
pieces, together with the excess mass and the amplification factor.}
Differentiating the mass constraint gives
$\lvert\delta\densrest\rvert/\lvert\delta\densR\rvert = f/(1-f)$, so the
density of the larger piece is always the better determined and the
threshold is exactly $f = \tfrac{1}{2}$. On Itokawa, Eros and 67P the
region is the smaller piece and $\densrest$ is the constrained quantity;
on Arrokoth the detector places $66.3\%$ of the volume in the region, so
there $\densrest$ is the density of the smaller piece and is worse
determined by a factor of $1.83$. Which piece is called the region is a
labelling choice and carries no physics: it changes neither uncertainty,
only which density a reporting rule names. Report $\Dmass$ alongside,
since it is the same statement and is the quantity that survives the
relabelling, and report $\ampfac$ with the plane at which it was
evaluated, since a fractional range in $\densrest$ corresponds to one in
$\Dmass$ that is \emph{larger} by that factor---on Itokawa, $2.95\%$
against $0.251\%$---and since a large $\ampfac$ is a warning rather than a
precision. The sensitivity to plane placement must likewise be quoted for
the quantity being reported: at the Itokawa reference plane it is
$\SI{10.0}{\kilo\gram\per\cubic\metre\per\metre}$ for $\denshead$, rising
to $55$ at the outer planes, against
$\SI{0.28}{\kilo\gram\per\cubic\metre\per\metre}$ for $\densrest$ and
$0.18\%$ per metre for $\Dmass$.

%---------------------------------------------------------------
\item \emph{Measured against an analytic ellipsoid, the geopotential
dispersion is more accurate than the mean surface slope at every facet count
tested, by a factor near $8$ at the $49{,}152$-facet resolution at which the
Itokawa and Eros results are computed, and by $19\pm4$ to $59\pm18$ over the
coarser range $2{,}000$ to $32{,}000$ facets.} The working-resolution figure
is the one a summary should quote, and it carries two qualifications: it is a
single shape and spin state rather than the average of nine, and the
dispersion error against which it is formed, $-0.02\%$, is quoted to one
significant figure, so the ratio is determined only to within roughly $6$ to
$12$. This test needs no assumption about which mesh is the
truth: the surface gravity of a homogeneous ellipsoid has a closed form,
and our evaluation of it reproduces the sphere limit to
$2.2\times10^{-14}$ in relative terms---a check on the numerical
quadrature, not on the polyhedral code, which cannot reproduce a sphere to
any such precision. Two qualifications belong with the factor. The
advantage narrows systematically with resolution, the two errors
converging as $N^{-1.26}$ and $N^{-0.91}$ so that their ratio falls as
$N^{-0.35}$; extrapolated to the $49{,}152$-facet meshes on which the
Itokawa and Eros results are computed it is of order $16$ to $20$, and
that is the figure a summary should quote. And the comparison is with
Paper~IV's \emph{ordering} of the two quantities, which it confirms
independently of the assumption Paper~IV could not test; it is not a
confirmation that Paper~IV's factor of $3$ to $14$ was conservative, since
mesh-to-mesh variation measures stability and the ellipsoid measures
accuracy, and a quantity can be stable and biased. The same experiment
shows vertex clustering and quadric edge collapse to be closely comparable
in accuracy, with mean absolute slope errors of $3.92\%$ and $4.05\%$, so
the numbers reported in Papers~III and~IV, which used clustering, are not
in question; quadric collapse is adopted here for topology rather than
accuracy. The ellipsoid also bounds its own use: decimating over the same
range moves the dispersion by about $0.18$ percentage points on these
smooth bodies against the $0.9$--$5.9\%$ Paper~IV measured on real ones,
so surface roughness dominates the envelope and the analytic test
validates its ordering rather than its magnitude.
% Resolved: standard errors are now given in Tables 8.3 and 8.4, and the
% decimator comparison is made on the paired differences, +0.019 +/- 0.012
% per cent on the dispersion and +0.12 +/- 0.20 on the slope, neither
% significant. The 49,152-facet row was abandoned as disproportionately
% expensive; one configuration completed and gives a ratio near 8, quoted
% in Section 8.3 as an indication.

\end{enumerate}

Three qualifications frame these results.

The relaxation premise on which the whole approach rests is unverified for
every body examined, and on Itokawa, where three of the four conditions of
\citet{richardson2014} are supported, the first fails: they tabulate its
mean rotational potential as $2.4\%$ of the mean gravitational potential
and their global bulk-density inference on it returns
$\SI{330}{\kilo\gram\per\cubic\metre}$ against a measured $1950$. The
single strong signal reported here is thus on the body where the premise
is weakest. Two things bear on how much weight that should carry.
Section~\ref{sec:discussion:conditioning} shows that the failing condition
is the condition under which a \emph{bulk-density} inversion is not
identically flat, and that it has no counterpart in a fixed-mass
two-region inversion: substituting the mass constraint into the
geopotential makes the surface potential affine in the single free
parameter, $\Upot_i = \alpha_i + \beta_i\densR$, and $\alpha$ and $\beta$
are proportional only if $\Mtot = 0$, so an interior stationary point
survives even in the limit of no rotation. The mass constraint fixes the
scale that the rotational term fixes in their problem. That is a statement
about conditioning and not about physics: it removes a reason to expect
failure and does not establish that the surface has relaxed. Empirically,
the inference is insensitive to the assumed mass over a factor of four and
agrees with a YORP determination that makes no appeal to relaxation at
all---though that determination is itself only $1.75$ standard deviations
from zero, and its interval contains the centre-of-mass offset at every
plane of the sweep, so it corroborates the sign and order of magnitude of
the anomaly and discriminates neither the plane nor the value. Against
this, Itokawa's residual dispersion after the best fit, $7.2\%$, is large
in absolute terms and larger than Eros's $3.4\%$; that comparison is not
controlled for shape, and the residuals for the other four are $98.7\%$ on Arrokoth,
$99.4\%$ on 67P, $99.6\%$ on Eros and $100.0\%$ on Bennu, so Itokawa's
$83.1\%$ stands apart from a tightly grouped set that includes two other
bilobate bodies. The $\omega = 0$ control of Section~\ref{sec:discussion}
has also been run: with the rotational term removed the scan still returns
an interior minimum, at $\SI{2950}{\kilo\gram\per\cubic\metre}$ with an
improvement of $17.2\%$, so the conditioning argument is confirmed in the
output; the contrast moves by $7.1\%$, which is larger than the error
budget and should be read beside it.
% Both done: Table 8.5 is complete and the omega = 0 control is reported
% (v10_s8_tests.json).

The model admits two uniform regions divided by a plane, a form known to
be insufficient for two of the bodies studied and shown by
\citet{kanamaru2019pss} to become largely unconstrained if relaxed to
three regions on Itokawa. A mass anomaly does not specify the region that
carries it, which is why the head--neck question raised by
\citet{scheeres2008} and revisited by \citet{lowry2014} and
\citet{kanamaru2019pss} remains open in general: structures placing
comparable excess mass at comparable distance from the centre of figure
produce comparable centre-of-mass offsets and comparable dispersion
reductions, and no surface measurement resolves them. The head and neck
models of this paper are not an instance of that degeneracy and should not
be cited as one---they place excess masses differing by a factor of $2.2$
and are separated by the objective by a factor of $7.3$ in improvement and
by $5.8$ against $2.0$ on $\Sstat$, so on the global surface the head
model is preferred and the neck model does not clear the resolution
envelope at all. What the two models do show is that the invariance
established here holds under a change of plane position within a fixed
region shape and not under a change of region shape, the two returning
background densities $\SI{85}{\kilo\gram\per\cubic\metre}$ apart.

The recovered densities are quantized at the
$\SI{25}{\kilo\gram\per\cubic\metre}$ scan step, with no sub-grid
refinement, and the consequences are visible throughout: two planes
$\SI{0.8}{\metre}$ apart returning the same $\denshead$, two non-monotone
entries in the centre-of-mass column, and an improvement that can come out
marginally negative where the grid does not contain the uniform solution,
which is an arithmetic impossibility for the underlying problem. Since
$\potvarn^{2}$ is a ratio of quadratics in $\densR$ once the mass
constraint is substituted, its stationarity condition is linear with a
single root, and solving it in closed form would remove the grid from the
error budget entirely and yield the curvature interval as the roots of a
quadratic rather than as a fitted quantity. We have not done so, and it is
the clearest improvement available to a future implementation.

The natural continuation is a calibrated detection limit. The four
weak-signal cases reported here are individually explicable, but the study
cannot say whether one strong signal in five reflects the rarity of large
interior contrasts or the sensitivity floor of the method. Such a test
cannot be performed by planting a contrast in a body and leaving its shape
alone, since the objective is built from the shape and never reads a field
computed beforehand; the premise requires that the surface have relaxed
onto an equipotential of the true interior, so the shape must be built
before the inversion is applied. Our attempt to do so did not converge
(Section~\ref{sec:discussion:rate}), and its mode of failure is
instructive: the plane holding a fixed volume fraction moved as the
surface relaxed, so the shape and the dividing surface co-evolve, and a
fixed point requires solving for both together---which is the coupling
Section~\ref{sec:plane} reports from the other direction. A stable
relaxation scheme is a piece of work in its own right. It would settle
both the detection limit and the interpretation of the synthetic
experiment, which as performed here varies shape rather than density, and
we regard it as a paper rather than an appendix.

\vspace{1em}
\noindent\textbf{Data availability.} All code, driver scripts and derived
tables, including the five-layer self-test and the mass-balance audit over
all thirteen Itokawa planes, are archived at Zenodo
(\href{https://doi.org/10.5281/zenodo.23003040}{doi:10.5281/zenodo.23003040})
under CC~BY~4.0, and the archived configuration records the scan interval
used at every plane so that the widening described in
Section~\ref{sec:itokawa:sweep} can be audited rather than taken on trust.
Development is tracked at
\url{https://github.com/praponkan/shape-based-surface-gravity}. Shape
models are not redistributed and are identified by their published
references in Section~\ref{sec:methods}; Table~\ref{tab:shapes} gives the
facet counts, volumes, masses and spin periods actually used, and the
deposit records the archive that holds each model.
% DOI reserved 2026-09-28: 10.5281/zenodo.23003040, filled here and in
% Section 2.9. Confirm before submission that the published record is the
% one that produced the numbers.


\appendix
\section{Plane clipping and its verification}
\label{app:clip}

This appendix sets out the region construction of
Section~\ref{sec:methods:clip} in enough detail to be reimplemented, derives
the closed-form quantities used to verify it, and states the tolerance
treatment on which its robustness depends.

\subsection{The algorithm}
\label{app:clip:alg}

Let $(\mathcal{V},\mathcal{F})$ be a closed, outward-oriented triangle mesh and
let the cutting plane be $x = c$. We construct the closed mesh bounding the
solid $\{x > c\}$; the complement and the slab of the neck model follow by
reversing the sense and by clipping twice.

Each vertex is classified by its signed distance $d_i = x_i - c$ into one of
three states,
\begin{equation}
\mathrm{code}(i) =
\begin{cases}
+1, & d_i > \tau, \\
\phantom{+}0, & |d_i| \le \tau, \\
-1, & d_i < -\tau,
\end{cases}
\label{eq:appclassify}
\end{equation}
with $\tau$ the tolerance of Section~\ref{app:clip:tol}. Vertices in state $0$
are snapped exactly onto the plane, $x_i \leftarrow c$, and are thereafter
treated as belonging to the kept side.

A facet is discarded if all three codes are $\le 0$ and retained unchanged if
all three are $\ge 0$. Otherwise it is clipped. Traversing its three edges in
order, the vertex $i$ is emitted whenever $\mathrm{code}(i) \ge 0$, and an
intersection point is emitted whenever an edge changes sign \emph{strictly},
that is when $\mathrm{code}(i)\,\mathrm{code}(j) < 0$, at
\begin{equation}
\mathbf{p} = \mathbf{v}_i + t\,(\mathbf{v}_j - \mathbf{v}_i),
\qquad
t = \frac{d_i}{d_i - d_j},
\label{eq:appintersect}
\end{equation}
with the $x$-coordinate of $\mathbf{p}$ set to $c$ exactly rather than left to
rounding. The requirement that the sign change be strict is what prevents a
vertex lying on the plane from generating two coincident intersection points;
this is developed in Section~\ref{app:clip:tol}. The resulting polygon is
emitted as a fan from its first vertex after consecutive coincident points have
been removed, and is discarded if fewer than three distinct vertices remain.

Intersection points are cached by the edge that produced them. An edge is
identified by the unordered pair of its endpoint indices, so the two facets
sharing a cut edge receive the same new vertex index. Created per facet
instead, they would receive two distinct indices at the same position: the
volume would be unchanged, since it is a sum of signed tetrahedra, but the
surface would be split along the entire cut and would fail the manifold test of
Section~\ref{app:clip:verify}.

\subsection{Closing the opening}
\label{app:clip:cap}

Every edge of the retained surface whose two endpoints both lie on the plane
and which has non-zero length is a boundary edge. These are traced into closed
loops by following, from each edge, the unique unvisited outgoing edge at its
head; the traversal terminates when it returns to its start or runs out. Each
loop is then capped separately, to its own apex: for a loop $\mathcal{L}$ with
vertex set $\mathcal{V}_\mathcal{L}$, the apex $\mathbf{a}$ is the centroid of
$\mathcal{V}_\mathcal{L}$ projected onto the plane, and for each
$(i,j) \in \mathcal{L}$ the triangle $(\mathbf{a}, \mathbf{v}_j, \mathbf{v}_i)$
is added, the reversed traversal orienting the cap outward.

That a fan closes a planar loop correctly, for any coplanar apex, follows from
the signed-area identity: for a closed loop $\{\mathbf{v}_k\}$ in a plane with
unit normal $\nhat$,
\begin{equation}
\sum_k \tfrac{1}{2}\,\nhat \cdot
\big[(\mathbf{v}_k - \mathbf{a}) \times (\mathbf{v}_{k+1} - \mathbf{a})\big]
\;=\; A_{\mathrm{signed}},
\label{eq:appfan}
\end{equation}
independently of $\mathbf{a}$, provided $\mathbf{a}$ lies in the plane.

The identity holds for each loop separately, and this is why the loops must be
traced rather than fanned together. A single apex shared between two disjoint
loops still gives the correct total signed area---so the volume would be
right---but the resulting triangles are not a surface: the two loops become
connected through a vertex that belongs to both, and the mesh is no longer a
manifold. The condition arises whenever a cross-section is multiply connected
or the region falls into more than one piece, and both occur in this study. The
Arrokoth shape model is supplied as two unjoined shells, so a plane beyond its
neck produces a region in two pieces, each with its own cap.

\subsection{Tolerance and degenerate facets}
\label{app:clip:tol}

The tolerance is set to
\begin{equation}
\tau = 10^{-10} \, \lVert \mathbf{v}_{\max} - \mathbf{v}_{\min} \rVert ,
\label{eq:apptol}
\end{equation}
the bounding-box diagonal of the mesh, so that it scales with the body.

Its purpose is to prevent a specific degeneracy. Suppose a vertex $k$ lies
exactly on the plane, $d_k = 0$, and is classified by a strict test as
\emph{outside}. Both edges incident on it from inside then satisfy
$d_i > 0 > d_k$ formally, and Eq.~\eqref{eq:appintersect} returns $t = 1$ for
each, placing both intersection points at $\mathbf{v}_k$ itself. The clipped
polygon acquires two coincident vertices, and the triangles formed with them
have zero area, on the surface and on the cap alike.

The consequence is more severe than a small geometric error. A facet of zero
area has an undefined normal, and in the Werner--Scheeres formulation the
per-facet contribution is assembled from that normal; the evaluation returns
NaN not at points near the facet but at every field point on the body. The
failure is therefore total rather than local, and it is parameter-dependent: at
$\xsplit = 0.1608$~km on Itokawa no such vertex is encountered and the
computation is clean, while at $0.1300$~km four are, and all $49{,}152$ field
values are lost.

Two guards are applied after clipping. Facets with a repeated vertex index are
removed, as are facets whose area falls below $10^{-12}$ times the median
positive area. In practice the classification of Eq.~\eqref{eq:appclassify}
prevents such facets from being created at all: on the runs reported in this
paper the number removed is zero.

Slivers---facets of small but non-zero area, produced where a vertex lies just
outside the tolerance---are not removed, and the polyhedral gravity
implementation flags them as potentially unstable. Their contribution is
proportional to their area and should vanish, but catastrophic cancellation is
conceivable when a facet millimetres across is evaluated at coordinates of
order kilometres. This was therefore measured directly. Raising $\tau$ by four
orders of
magnitude, to $10^{-6}$ of the diagonal, removes them: on 67P the minimum facet
area rose from $2.381\times10^{-12}$ to $3.660\times10^{-8}$~km$^2$ and eight
facets were eliminated. Every reported quantity agreed to the digits printed,
the dispersion at seven planes being identical to eight decimals and the
recovered contrast, improvement and excess mass unchanged in the seventh
significant figure or beyond. On a synthetic body the agreement was exact.

\subsection{Closed-form quantities for verification}
\label{app:clip:analytic}

Two solids with elementary closed forms are used to verify the construction.

For a sphere of radius $R$ cut at $x = c$, write $h = R - c$ for the height of
the resulting cap. Integrating the disc of radius $\sqrt{R^2 - x^2}$ over
$x \in [c, R]$ gives
\begin{equation}
V_{\mathrm{cap}} = \pi \int_{c}^{R} (R^2 - x^2)\,\mathrm{d}x
= \frac{\pi h^{2} (3R - h)}{3},
\qquad
\bar{x}_{\mathrm{cap}} = \frac{3\,(2R - h)^{2}}{4\,(3R - h)} .
\label{eq:appcap}
\end{equation}
For $R = 1$ and $c = 0.5$ these are $0.6544985$ and $0.6750000$.

For the slab $|x| \le a$ of the same sphere,
\begin{equation}
V_{\mathrm{slab}} = \pi \int_{-a}^{a} (R^2 - x^2)\,\mathrm{d}x
= 2\pi\!\left( R^{2} a - \frac{a^{3}}{3} \right),
\label{eq:appslab}
\end{equation}
which for $R = 1$ and $a = 0.2$ gives $1.2398819$.

The corresponding radial cone---the solid produced by the construction of
Appendix~\ref{app:cone}, whose apex is the sphere's centre and whose base is
the spherical cap, and which for a sphere is a spherical sector---has
\begin{equation}
V_{\mathrm{sector}} = \tfrac{1}{3}\,R\,A_{\mathrm{cap}}
= \tfrac{1}{3}\,R\,(2\pi R h)
= \frac{2\pi R^{2} h}{3},
\qquad
\bar{x}_{\mathrm{sector}} = \tfrac{3}{4}\!\left( R - \frac{h}{2} \right),
\label{eq:appcone}
\end{equation}
giving $1.0471976$ and $0.5625000$ for the same cut. The centroid follows from
the fan of tetrahedra sharing the apex, whose centroids lie three quarters of
the way from apex to base, the base being a spherical cap of mean
$x$-coordinate $R - h/2$.

\subsection{Verification}
\label{app:clip:verify}

Table~\ref{tab:appverify} reports what the implementation returns on an
icosphere of $20{,}480$ facets.

\begin{table}[htbp]
\centering
\caption{Verification of the clipping construction against
Eqs.~\eqref{eq:appcap}--\eqref{eq:appcone}, on a unit icosphere of $20{,}480$
facets. The residuals near $0.1\%$ are the polyhedral approximation of the
sphere: at the same subdivision the whole-body volume differs from $4\pi/3$ by
$0.05\%$. The clipping itself is exact for the given polyhedron.}
\label{tab:appverify}
\begin{tabular}{lrrr}
\toprule
Quantity & Measured & Closed form & Difference \\
\midrule
Cap volume, $c = 0.5$        & $0.653931$ & $0.654498$ & $-0.087\%$ \\
Cap centroid, $c = 0.5$      & $0.674936$ & $0.675000$ & $-0.010\%$ \\
Slab volume, $a = 0.2$       & $1.239437$ & $1.239882$ & $-0.036\%$ \\
Sector volume, $c = 0.5$     & $1.047453$ & $1.047198$ & $+0.024\%$ \\
Sector centroid, $c = 0.5$   & $0.562213$ & $0.562500$ & $-0.051\%$ \\
\midrule
Complement closure error     & \multicolumn{3}{r}{$2.0\times10^{-14}\%$} \\
Predicate violation, lobe    & \multicolumn{3}{r}{$0$} \\
Predicate violation, sector  & \multicolumn{3}{r}{$0.5$} \\
Euler characteristic, lobe   & \multicolumn{3}{r}{$2$} \\
Facets failing containment   & \multicolumn{3}{r}{$0$ (one component)} \\
\bottomrule
\end{tabular}
\end{table}

Four checks are applied to every region, and a fifth where the shape model
has more than one component. Each inspects a property the others do not, and
the set was arrived at by failure: each was added after something passed all
the checks then in place.

The \emph{analytic} comparison tests magnitude but would be passed by any
construction returning the right volume for the wrong reason.

The \emph{complement} test, $V(x>c) + V(x<c) = \Vtot$, tests internal
consistency and is passed by the sector as readily as by the lobe, since both
partition the tetrahedral sum.

The \emph{predicate} test inspects the identity of the solid, requiring
$\min_i x_i \ge c$ over the vertices of the region. It is the only one of the
first three that the sector fails.

The \emph{manifold} test inspects topology. Every directed edge must appear
exactly once with its reverse present, and the Euler characteristic
$\chi = V - E + F$ must equal $2$ for each connected component. This is
insensitive to none of the same failures: a surface split along the cut, as
described in Section~\ref{app:clip:alg}, passes the analytic, complement and
predicate tests, because none of them consults connectivity, and fails only
this one. Conversely a sector passes the manifold test and fails the predicate.

The \emph{containment} test inspects the embedding. Where the mesh has more
than one component, each facet centroid is tested by ray casting for
containment in every other component, and those found inside are excluded from
the area-weighted sums. It catches what none of the other four can: a mesh in
two closed pieces is topologically sound whatever their arrangement in space,
so a facet buried inside the opposite shell passes every previous check while
contributing a geopotential evaluated at a point that is not on the surface.
On the Arrokoth model $455$ facets, $0.943\%$ of the area, fail it, and
excluding them moves the assumed bulk density at which the recovered contrast
crosses unity from about $345$ to just below
$\SI{250}{\kilo\gram\per\cubic\metre}$
(Section~\ref{sec:nulls:arrokoth}).

The criterion is stated per connected component rather than as $\chi = 2$,
because a region cut from a body supplied as several unjoined shells is
legitimately in several pieces. Of the shape models used here, Arrokoth has
$\chi = 4$ over two components as supplied, and the 67P SPC~2017 model is not
closed at all, with one repeated and two unmatched directed edges. Both defects
are inherited by any region cut from those models, and neither affects a volume
or a potential.

The same containment machinery bears on the volume. Two components that
overlapped in space would have the material in the overlap counted twice,
since the divergence-theorem volume sums over all facets without regard to
whether they bound the same region. Measured by Monte Carlo over the
overlapping bounding box with $200{,}000$ samples, the intersection of the two
Arrokoth shells is $0.044 \pm 0.001$ per cent of the total volume, consistent
with zero: the shells touch but do not interpenetrate, and no material is
double counted. It is the surface average, not the volume, that the buried
facets corrupt.

A stress case is worth recording. On an icosphere subdivided so that $128$
vertices lie exactly on $x = 0$, cutting at that plane returns a hemisphere of
volume $2.093260$ against the closed-form $2\pi/3 = 2.094395$, a difference of
$-0.054\%$, with $\chi = 2$, zero degenerate facets and a minimum facet area of
$1.92\times10^{-4}$. A strict-inequality classification would produce
coincident points on every one of those $128$ vertices, which is why the
tolerance-based classification of Section~\ref{app:clip:tol} is required.

\subsection{Detecting the dividing plane}
\label{app:clip:neck}

The saddle detector of Section~\ref{sec:methods:neck} operates on the binned
profile $w(x)$, the ninetieth percentile of facet-centroid distance from the
$x$ axis within each of sixty bins. Two aspects of its implementation are worth
stating.

Candidates are identified on \emph{runs} of equal value rather than at single
points. A run $[j,k]$ with $w_j = \dots = w_k$ within tolerance is a local
minimum if $w_{j-1} > w_j$ and $w_{k+1} > w_k$, and its position is taken as
the midpoint. A strict point test, requiring $w_j < w_{j\pm1}$, fails on a
symmetric profile, where two samples straddle the minimum with equal values;
applied to Itokawa it returned no neck at all. It also fails on any plateau.

The prominence of Eq.~\eqref{eq:prominence} is measured against the
\emph{lower} of the two flanking maxima, in the manner of the topographic
quantity of the same name. This has a consequence worth noting: Eros scores
higher than Itokawa, $0.216$ against $0.170$, despite its shallower saddle,
because its two flanking humps are nearly equal ($7.9072$ and $7.8902$~km,
differing by $0.2\%$) while Itokawa's differ by $15\%$ ($0.1437$ and
$0.1215$~km). Prominence measures how well separated the lobes are, not how
deep the neck is, and the two are not the same.

The detector is also mesh-dependent. On Itokawa it returns $0.1702$~km on the
$20{,}000$-facet mesh and $0.1608$~km on the meshes of $30{,}000$, $40{,}000$
and $49{,}152$ facets, a difference of $9.4$~m. Any test that holds the
plane fixed must therefore supply it explicitly rather than rely on the
detector returning the same value twice.

The threshold $p \ge 0.02$ is not a fine judgement. Across the bodies examined
the accepted values run from $0.490$ down to $0.091$, and the single rejected
value is $0.0004$: the gap spans more than two orders of magnitude. Where no
candidate qualifies the routine returns no plane, and the neck model is skipped
rather than run on a guess.



\section{The conical-region error}
\label{app:cone}

This appendix records an error made during the development of this work, why
the verification then in place could not detect it, and how it was eventually
found. It is included because the failure is instructive rather than
idiosyncratic: the construction that produced it is a natural one, the mesh it
produces is valid by every ordinary measure, and the resulting scientific
conclusion was not merely inaccurate but inverted.

\subsection{The construction}
\label{app:cone:construction}

The task is to build a closed mesh bounding the material beyond a plane. A
natural approach, and the one first taken, avoids clipping altogether. Select
the facets whose centroids satisfy $x > c$; these form an open surface patch.
Find the boundary edges of the patch---those belonging to exactly one selected
facet---and close the surface by adding, for each such edge $(i,j)$, the
triangle $(O, \mathbf{v}_j, \mathbf{v}_i)$, where $O$ is a point inside the
body, in our case the mean of the vertices.

The appeal of this is that it requires no geometric predicates beyond a
centroid test, no intersection arithmetic, and no tolerance. The resulting mesh
is closed, consistently oriented, and its volume, computed by the divergence
theorem, agrees with the sum of the signed tetrahedra
$\{(O, \mathbf{v}_a, \mathbf{v}_b, \mathbf{v}_c)\}$ over the selected facets to
machine precision. All of this is true. None of it is relevant.

\subsection{What the mesh actually bounds}
\label{app:cone:solid}

The solid enclosed is not the material beyond the plane. It is the radial cone
with apex at $O$ subtended by the selected surface patch: the union of all
segments joining $O$ to a point of the patch. For a sphere with $O$ at the
centre this is a spherical sector. The two solids share their outer boundary
and nothing else. The plane-cut lobe is bounded below by the plane $x = c$; the
sector is bounded below by a lateral surface running from the rim of the patch
to the interior point $O$, and therefore includes material at every $x$ between
$x_O$ and $c$ that the intended region excludes.

For a sphere of radius $R$ cut at $x = c$, with $O$ at the centre, both solids
admit closed forms, given as Eqs.~\eqref{eq:appcap} and~\eqref{eq:appcone} of
Appendix~\ref{app:clip}. At $R = 1$, $c = 0.5$:
\begin{equation}
\frac{V_{\mathrm{sector}}}{V_{\mathrm{cap}}}
= \frac{2\pi R^{2}h/3}{\pi h^{2}(3R-h)/3}
= \frac{2R^{2}}{h(3R-h)}
= 1.600 ,
\qquad
\bar{x}_{\mathrm{cap}} - \bar{x}_{\mathrm{sector}} = 0.1125\,R .
\label{eq:appconeratio}
\end{equation}
The sector has sixty per cent more volume and its centroid lies more than a
tenth of a radius further in. Figure~\ref{fig:cone} shows the two.

The discrepancy is not a small-parameter effect that vanishes for a thin cap.
As $h \to 0$ the ratio in Eq.~\eqref{eq:appconeratio} diverges as
$2R/(3h)$: the thinner the intended region, the worse the substitution. This
matters directly, because Section~\ref{sec:itokawa:sweep} advances the plane
precisely into that regime.



\subsection{The consequence on Itokawa}
\label{app:cone:consequence}

The density model assigns an elevated density to the region and reduces the
complement to conserve the total mass. With the correct region, raising
$\denshead$ adds mass in the head, displacing the centre of mass toward $+x$.
With the sector, it adds mass along a wedge whose centroid lies at $+0.5625R$
rather than $+0.6750R$ on the sphere analogue, and on Itokawa reaches back past
the neck into the body.

The recovered contrast inverted. Table~\ref{tab:appcone} gives both models as
they were computed with the conical construction, alongside the plane-cut
result at the same plane on the same shape model.

\begin{table}[htbp]
\centering
\small
\caption{The Itokawa inversion computed with the conical region, against the
plane-cut result. The conical runs used a mesh decimated to $41{,}369$ facets
and a body volume of $0.0177848$~km$^{3}$; the plane-cut column is the
reference run of Section~\ref{sec:itokawa:reference} on the unreduced mesh.
Both conical models return a head several times \emph{less} dense than the
body, with the minimum pinned to the lower edge of the density scan.}
\label{tab:appcone}
\begin{tabular}{lrrr}
\toprule
 & \multicolumn{2}{c}{Conical region} & Plane-cut \\
\cmidrule(lr){2-3}\cmidrule(lr){4-4}
Quantity & head & neck & head \\
\midrule
Region facets            & $6{,}412$ & $5{,}730$ & $7{,}950$ \\
Region volume (km$^3$)   & $0.004229$ & $0.001486$ & $0.002922$ \\
Volume fraction          & $23.8\%$ & $8.4\%$ & $16.4\%$ \\
Volume closure error     & $2.1\times10^{-14}\%$ & $0\%$ & $2.1\times10^{-14}\%$ \\
$\densR$ (kg\,m$^{-3}$)  & $300$ & $300$ & $2800$ \\
$\densrest$ (kg\,m$^{-3}$) & $2547$ & $2169$ & $1858$ \\
Contrast                 & $0.118$ & $0.138$ & $1.507$ \\
Interior minimum         & no & no & yes \\
Improvement              & $6.121\%$ & $3.831\%$ & $16.889\%$ \\
\bottomrule
\end{tabular}
\end{table}

Three features of the conical columns are consistent with a region that reaches
into the body. The recovered density sits at $300$~kg\,m$^{-3}$, the lower
limit of the scan, in both models: the objective function is minimized by
removing as much mass as the scan permits from a wedge that spans the neck, and
would have gone lower had it been allowed to. Neither minimum is interior, so
the edge-minimum guard of Section~\ref{sec:failures:edge} would flag both---but
that guard was added later, and at the time the values were reported as
results. And the head region holds $23.8\%$ of the volume against $16.4\%$ for
the plane-cut region at the same plane, the excess being precisely the wedge
between $O$ and the plane.

The volume closure error is the point of the table. At
$2.1\times10^{-14}$ per cent it is as small as floating-point arithmetic allows,
and it is identical for the conical and the plane-cut construction. A check
that returns the same value for a correct solid and for one wrong by a factor
of $1.6$ in volume is measuring something other than correctness.




\subsection{Why the verification did not catch it}
\label{app:cone:verification}

The construction was accompanied by a check described as a volume-closure test.
It compared the volume of the assembled mesh, computed by the divergence
theorem over its facets, against the sum of the signed tetrahedra used to build
it. On the Itokawa head region it agreed to $2\times10^{-14}$ per cent.

The agreement was exact and meaningless. Both quantities describe the same
solid---the sector---and the test verified that the mesh had been assembled
consistently from the tetrahedra, which it had. A test constructed this way
cannot fail for the reason that mattered, and therefore carries no information
about it. Nor would the other checks of Section~\ref{sec:methods:verify} have helped, had
they existed at the time. The complement test is passed by both constructions,
since sectors over complementary patches partition the body just as plane-cut
lobes do. The manifold test is passed by the sector, which is a perfectly good
closed surface. And the containment test finds nothing, because the sector is a
single component and there is no other component for any of its facets to lie
inside. Of the five checks the paper now applies, exactly one distinguishes the
two solids.

What distinguishes them is a statement about their vertices rather than their
volumes, their topology or their embedding. Every vertex of the plane-cut region
satisfies $x \ge c$ by construction. The sector contains $O$, which does not.
Testing
\begin{equation}
\max_i \,(c - x_i) \le \epsilon
\label{eq:apppredicate}
\end{equation}
over the vertices of the region returns $0$ for the lobe and $0.5$ for the
sector in the units of the sphere test, where the plane sits at $x = 0.5$ and
$O$ at the origin. The test costs three lines and inspects the property that
was actually in doubt.

\subsection{How it was found}
\label{app:cone:discovery}

Not by any check. The conical results were internally consistent, passed every
test then applied, and gave a smooth, well-behaved dispersion curve. They were
questioned because the recovered contrast disagreed in \emph{sign} with
\citet{lowry2014} and \citet{kanamaru2019}, both of which report a denser head.

The diagnosis followed from constructing the same region on a unit sphere,
where the closed forms of Appendix~\ref{app:clip:analytic} are available. The
tetrahedral construction returned a volume of $1.047453$ and a centroid of
$0.562213$, matching the sector formulae to five and four decimals
respectively, and matching the spherical cap not at all.

Three observations follow that are worth carrying beyond this particular error.

A disagreement with published work is not a verification method, but it is
sometimes the only signal available. Had Itokawa's true structure been unknown,
or had the sign happened to come out right, the error would have survived.

Analytic test cases are cheap and decisive. A sphere has closed-form caps,
slabs and sectors; the entire diagnosis took one script and no shape model. The
same lesson recurs in Section~\ref{sec:discussion:ellipsoid}, where an analytic
ellipsoid settles in one experiment a question that mesh-to-mesh comparison had
answered wrongly. It is worth constructing such a case before trusting a
numerical routine, not after.

And a verification test should be designed against a specific hypothesis of
failure. The question is not whether a check passes, but what it would look
like if the code were wrong in a way one has not thought of. The volume-closure
test answered the question ``was the mesh assembled consistently?''; nobody had
asked ``is this the solid I meant?'' The same pattern recurred later in the
work, with the buried facets of Section~\ref{sec:failures:buried}: the manifold
test, added precisely because volume comparison had proved insufficient, was
itself insufficient against a failure nobody had yet imagined. Each check
answers the question its author posed, and the list of questions is never
known to be complete.

\subsection{Status of the superseded results}
\label{app:cone:status}

All results computed with the conical construction are superseded, and none
appear in this paper except in Table~\ref{tab:appcone}, where they are shown
for comparison. They are retained in the archived results file
(Section~\ref{sec:conclusions}, data availability) under a key marking them as
superseded, together with the diagnosis, so that the record of the work is
complete and the error can be reproduced by anyone wishing to check this
account.

One further point should be recorded. The pilot study preceding this work,
whose results informed Paper~IV, was internally inconsistent in a related way:
it evaluated the potential on a planar-capped head mesh while computing the
region volume as a sum of origin-anchored tetrahedra, which is a sector volume.
The mismatch is the likely origin of a three per cent volume discrepancy noted
at the time and left unexplained. No conclusion of Paper~IV depends on it,
since that paper reports no density inversion, but the episode is recorded here
for completeness.



\section{The synthetic contact binaries}
\label{app:synth}

This appendix specifies the synthetic bodies used in
Section~\ref{sec:constrains:synthetic}, states the design decisions that make
them a controlled experiment rather than a demonstration, and records two
errors in the first version of that experiment.

\subsection{Why synthetic bodies}
\label{app:synth:why}

The background density, and with it the excess mass that
Eq.~\eqref{eq:excessidentity} ties to it, is the more invariant quantity on
Itokawa and the less invariant one on 67P and Arrokoth. Three real bodies establish that the
behaviour varies; they cannot establish what it varies with. A fourth real body
would not settle it either. We cannot choose in advance whether a given
asteroid carries a detectable anomaly, and no body's interior is known
independently, so agreement would establish only that something had occurred
twice.

What is needed is a family in which the strength of the signal can be set by
the experimenter while everything else is held fixed. That requires synthetic
bodies. The cost is that they are idealised, and Section~\ref{app:synth:limits}
states what that costs.

\subsection{Construction}
\label{app:synth:build}

Each body is the union of prolate lobes of revolution about $x$. The outer
half-width is
\begin{equation}
w(x) = \max_{i}\; A_i \sqrt{1 - \left(\frac{x - c_i}{a_i}\right)^{2}},
\label{eq:appsynthprofile}
\end{equation}
taken as zero where the radicand is negative, with $A_i$, $c_i$ and $a_i$ the
amplitude, centre and half-length of lobe $i$. The surface is meshed directly
by revolving the profile: $100$ stations in $x$, $48$ around, with apex
vertices closing both ends, giving about $9{,}500$ facets. Orientation is fixed
by requiring the total signed volume to be positive.

Using a union of analytic lobes rather than a boolean operation on two meshes
avoids the intersection curve entirely, so no body in the family contains a
seam, a sliver, or a non-manifold edge. It also makes each body a single
connected component, so the containment test of
Section~\ref{app:clip:verify} has nothing to find: the failure mode that
affects the Arrokoth shape model cannot arise here. The profile is trimmed at a half-width
of $10^{-6}$ of its maximum rather than at an absolute threshold; an absolute
one leaves a ring of vertices a few nanometres from the axis near each tip,
distinct by index but coincident in position, which is enough to make the mesh
fail a manifold test. Every body in the family, as built, has
$\chi = 2$ with no repeated or unmatched directed edges, and no clipped region
in the experiment was rejected by the manifold check.

\subsection{The family}
\label{app:synth:family}

The eight bodies fall into three groups, and the grouping is the design.

A near-spherical control stands in for a body with no interior structure. A
\emph{symmetric} family of equal lobes is symmetric about its own neck, so the
density that flattens the geopotential should be close to uniform and the
signal weak by construction, whatever the depth of the neck. An
\emph{asymmetric} family of unequal lobes has a geometric mass asymmetry for
the inversion to find. Within each family the neck is deepened by moving the
lobes apart while the enclosed volume varies by no more than seven per cent, so
that depth of neck and size of body are not confounded.

\begin{table}[htbp]
\centering
\small
\caption{The synthetic family. Lobes are given as $(A_i, c_i, a_i)$ in km.
Prominence is that measured by the detector of
Section~\ref{sec:methods:neck}; the neck ratio is the half-width at the saddle
divided by the largest half-width. All bodies were run at a bulk density of
$2000$~kg\,m$^{-3}$ and a rotation period of $12.13$~h, with seven planes
spanning a factor $1.900$ in region volume, and all have $\chi = 2$.}
\label{tab:appsynth}
\begin{tabular}{llrrr}
\toprule
Body & Lobes $(A, c, a)$ & $\Vtot$ (km$^3$) & Neck ratio & Prominence \\
\midrule
\texttt{sphere} & $(0.15, 0.00, 0.16)$ & $0.01503$ & $0.650$ & none \\
\addlinespace
\texttt{sym\_waist} & $(0.14, \mp0.11, 0.185)$ & $0.02705$ & $0.787$ & $0.170$ \\
\texttt{sym\_neck} & $(0.14, \mp0.15, 0.185)$ & $0.02951$ & $0.587$ & $0.360$ \\
\texttt{sym\_deep} & $(0.14, \mp0.18, 0.185)$ & $0.03025$ & $0.235$ & $0.627$ \\
\addlinespace
\texttt{asym\_waist} & $(0.15, -0.13, 0.20)$, $(0.12, +0.12, 0.17)$ & $0.02707$ & $0.653$ & $0.124$ \\
\texttt{asym\_neck} & $(0.15, -0.13, 0.20)$, $(0.12, +0.18, 0.17)$ & $0.02849$ & $0.485$ & $0.272$ \\
\texttt{asym\_deep} & $(0.15, -0.13, 0.20)$, $(0.12, +0.22, 0.17)$ & $0.02893$ & $0.289$ & $0.449$ \\
\texttt{asym\_extreme} & $(0.15, -0.13, 0.20)$, $(0.10, +0.20, 0.14)$ & $0.02461$ & $0.196$ & $0.362$ \\
\bottomrule
\end{tabular}
\end{table}

The bodies are on the scale of Itokawa, a few hundred metres, and share its
rotation period. The near-spherical control has no saddle and its plane was
imposed, as Bennu's was.

\subsection{Selecting the planes}
\label{app:synth:planes}

Planes are chosen by bisection on the fraction of volume they leave above them,
not by position. For a target fraction $f$ the plane is found by bisecting on
$x$ until $\lvert V(x > \xsplit)/\Vtot - f \rvert$ is below tolerance; thirty
iterations give agreement to four decimal places in the fraction, at a cost of
under a second per plane.

Seven targets are placed geometrically about the fraction at the detected
saddle, spanning a factor
\begin{equation}
\frac{f_{\max}}{f_{\min}} = 1.900 ,
\label{eq:appspan}
\end{equation}
chosen to match the Itokawa sweep. This is the point on which the experiment
turns. A plane moved a fixed distance sweeps very different amounts of volume
on bodies of different shape---on Arrokoth, whose neck is long and thin,
$3.5$~km of travel changes the region volume by only five per cent against
ninety on Itokawa---so equal distances do not probe the same thing. Selecting
by volume fraction makes the sweeps comparable across bodies, and it is also
the variable against which the exponent of Eq.~\eqref{eq:slope} is fitted.

\subsection{What is not being tested}
\label{app:synth:not}

It is worth being explicit, since the natural expectation of a synthetic
experiment is different. We do not impose a known density on a synthetic body
and ask the method to recover it. That test would fail by construction and
would establish nothing.

Potential-dispersion minimization does not fit observed data. It assumes that
the true density is the one that renders the surface geopotential most uniform.
It assumes that the true density is the one that renders the surface
geopotential most uniform, and the objective is built from the shape and the
trial density alone: a field computed beforehand from an imposed interior is
never read. Planting a contrast in a body and leaving its shape unaltered
therefore changes nothing the minimization can see, and the recovered value is
whatever the shape alone would have given. What the experiment examines
is the structure of the objective function: which quantity it holds fixed as
the dividing plane moves, and how that depends on the strength of the response.

Every body in the family is therefore of uniform density. What varies is shape.
A consequence is that the experiment cannot establish that a strong response
implies a physical mass anomaly, only that in the strong-response regime the
quantity held fixed is the background density, and with it $\Dmass$, rather
than $\densR$. A genuine ground-truth test would have to build the body the
premise describes---a shape whose surface has relaxed onto an equipotential of
a planted interior---before inverting it.
Section~\ref{sec:discussion:rate} reports our attempt at that construction and
why its outcome is not presented.

\subsection{Results}
\label{app:synth:results}

\begin{table}[htbp]
\centering
\small
\caption{Results for the synthetic family. Half-ranges are taken over the seven
planes of each sweep, and standard errors on $s$ are those of the
least-squares fit. Where the density difference changes sign no exponent
exists and the fractional range of $\Dmass$ is not meaningful, its mean lying
inside its spread. $\bar{\ampfac}$ is the amplification factor of
Equation~\eqref{eq:excessamplify} averaged over the sweep: it is what
separates the $\denshead$ and $\Dmass$ columns, so those two are not
comparable across rows without it, and it is undefined on the symmetric
bodies, where $\Dmass$ passes through zero. The last column is the
compensation ratio, the half-range of $\Ddens$ divided by that of $\Dmass$,
which is the part of the stability ratio that is not fixed by algebra
(Section~\ref{sec:constrains:observation}). The Itokawa row is the eleven
planes below the grain-density ceiling, the set whose span of $1.907$ matches
the $1.900$ used here; its maximum improvement over that set is $22.4\%$
rather than the $16.9\%$ quoted at the detected neck.}
\label{tab:appsynthres}
\begin{tabular}{lrrrrrr}
\toprule
Body & $\improve_{\max}$ & $s$ & $\bar{\ampfac}$ & $\denshead$ half-range
 & $\Dmass$ half-range & $\Ddens/\Dmass$ \\
\midrule
\texttt{sphere}       & $1.24\%$  & none & $44$ & $3.10\%$  & sign change & --- \\
\texttt{sym\_waist}   & $3.68\%$  & none & ---  & $5.61\%$  & sign change & --- \\
\texttt{sym\_neck}    & $4.35\%$  & none & ---  & $6.22\%$  & sign change & --- \\
\texttt{sym\_deep}    & $4.47\%$  & none & ---  & $5.60\%$  & sign change & --- \\
\addlinespace
\texttt{asym\_waist}  & $16.67\%$ & $-2.020\pm0.109$ & $15.0$ & $9.28\%$  & $31.3\%$ & $1.9$ \\
\texttt{asym\_neck}   & $17.63\%$ & $-2.366\pm0.165$ & $16.2$ & $9.37\%$  & $39.2\%$ & $1.7$ \\
\texttt{asym\_deep}   & $19.50\%$ & $-2.069\pm0.094$ & $14.9$ & $8.81\%$  & $29.1\%$ & $2.0$ \\
\addlinespace
\texttt{asym\_extreme}& $43.71\%$ & $-1.217\pm0.023$ & $8.1$  & $12.78\%$ & $6.2\%$  & $6.1$ \\
\midrule
Itokawa (real)        & $22.4\%$  & $-1.141\pm0.018$ & $11.8$ & $13.20\%$ & $5.12\%$ & $7.1$ \\
\bottomrule
\end{tabular}
\end{table}

The design succeeded in separating the groups. Every symmetric body, however
deep its neck, produces a weak signal and a density difference that changes
sign across its sweep; the deepest, \texttt{sym\_deep}, has a prominence of
$0.627$---the highest in the family, a more sharply necked profile than
\texttt{asym\_extreme}---and still returns only
$4.47\%$. Depth of neck is not signal strength. The asymmetric bodies produce
signals three to ten times larger at comparable neck depths, because they have
a mass asymmetry to find.

The exponent varies monotonically with signal strength across the four bodies
that have one, and the Pearson correlation is $0.95$ ($p = 0.05$ two-tailed, $n = 4$), while
the Spearman rank correlation over the same four is only $0.40$, the three
clustered bodies not being in rank order. Only \texttt{asym\_extreme}
approaches $-1$, the value at which Eq.~\eqref{eq:slopeequiv} makes the
background density the quantity held fixed as the plane moves, and it is the only body in the family that reproduces Itokawa's
behaviour---on the exponent, on the density half-range, and on the excess-mass
half-range simultaneously.

\subsection{Two errors in the first version}
\label{app:synth:errors}

Both are recorded because either would have produced a plausible and wrong
conclusion.

\emph{The sweeps were not comparable.} The first version swept a fixed
$\pm0.05$~km on every body. Because the bodies differ in how sharply their
cross-section changes near the neck, this produced volume-fraction spans of
$1.13$ to $1.54$ against the $1.90$ of the Itokawa sweep. A body whose region
volume barely changes cannot show either quantity varying, so the comparison
measured the sweep design rather than the bodies. The bisection scheme of
Section~\ref{app:synth:planes} was introduced to correct this.

\emph{The diagnostic broke down where it mattered most.} The first version
compared fractional half-ranges directly. That statistic divides by the mean,
and on a body with no anomaly the mean of $\Dmass$ sits inside its own spread,
so the quotient is arbitrary. The symmetric bodies returned apparent ratios of
up to $1623$, which were read at first as an extraordinarily strong result in
the opposite direction. They were artefacts of dividing by nearly zero. The
exponent of Eq.~\eqref{eq:slope} was introduced because it does not have this
failure mode: where the density difference changes sign, the exponent simply
does not exist, and its absence is the correct answer rather than a large
number.

\subsection{Limitations}
\label{app:synth:limits}

Four bodies in the family carry a measurable exponent, and three of those
cluster between $16.7\%$ and $19.5\%$ in signal strength while the fourth sits
at $43.71\%$. The correlation therefore rests on a single high-leverage point,
which is what the Spearman value of $0.40$ records.
The direction of the effect is supported; its functional form is not
determined, and a threshold anywhere between roughly twenty and forty-four per
cent fits the data as well as a gradual trend does. Filling that gap would
require bodies designed to land in it, which the present family does not
attempt.

The numerical threshold does not transfer to real bodies. \texttt{asym\_extreme}
requires $43.71\%$ to reach the regime in which the background density is
pinned, while Itokawa reaches it at $22.4\%$ over the eleven planes below
the grain-density ceiling, or $26.4\%$ if the compositionally excluded
planes are counted. The synthetic surfaces are smooth surfaces of revolution; real
bodies carry roughness at every resolved scale, which Paper~IV found
contributes an error floor of order two percentage points to the dispersion. No
threshold value from this experiment should be quoted as universal, and none is
quoted in this paper.

The experiment uses the same forward model in both directions and therefore
examines the objective function, not the accuracy of the gravity calculation.
That accuracy is tested separately in
Section~\ref{sec:discussion:ellipsoid} against a homogeneous ellipsoid with a
closed-form surface gravity.

Finally, the family varies one geometric parameter---the separation of two
lobes of fixed shape. Real contact binaries differ in more ways than that, and
whether the relation between signal strength and exponent survives variation in
lobe elongation, in the ratio of lobe sizes, or in spin rate has not been
tested here.



\section{The mesh-resolution sensitivity envelope}
\label{app:envelope}

The envelope quoted throughout this paper as $0.9$--$5.9\%$ is the range over
which the normalized geopotential dispersion $\potvarn$ moves when the same
body is meshed at different resolutions. It is used as a scale against which
an improvement may be judged large or small, and it is not a statistical
threshold: nothing about it licenses a significance statement. Because
several sections compare an improvement against the value the envelope takes
on one particular body, this appendix records those values, how they were
obtained, and how far they depend on choices that are not forced.

\subsection{Definition and provenance}
\label{app:envelope:def}

Paper~IV decimated four shape models---Eros, Itokawa, Gaspra and Ida---over a
ladder from roughly $2{,}000$ to $42{,}000$ facets and recorded $\potvarn$ and
the area-weighted mean slope at each rung. For each body the envelope is the
full range of $\potvarn$ across its ladder, as a fraction of the smallest
value:
\begin{equation}
E \;=\; \frac{\max_i \potvarn^{(i)} - \min_i \potvarn^{(i)}}
             {\min_i \potvarn^{(i)}} ,
\label{eq:envelope}
\end{equation}
the index running over the rungs. The published figures are $0.9\%$ at one
extreme and $5.9\%$ at the other, and the paper quotes that pair as a range
across the test set rather than as a property of any one body.

Two things follow that matter here. The envelope is a property of a body and
its ladder, not of the method, so comparing an improvement on Itokawa against
the largest value in a four-body set is the weaker of the two available
comparisons; and the ladder itself is a choice, so the envelope moves when the
decimator does.

\subsection{Recomputed values}
\label{app:envelope:values}

Table~\ref{tab:envelope} gives Eq.~\eqref{eq:envelope} recomputed for this
paper on the same four bodies, with both decimators, on ladders of six rungs
each. The dispersion column is the envelope; the slope column is the same
quantity for the mean surface slope, included because the comparison between
the two is the result of Paper~IV; and the plateau column is the change in
$\potvarn$ between the two finest meshes alone.

\begin{table}[htbp]
\centering
\small
\setlength{\tabcolsep}{5pt}
\caption{The envelope of Eq.~\eqref{eq:envelope} recomputed on the four
Paper~IV resolution bodies, six rungs each. The slope column is the same
range for the area-weighted mean slope. The plateau is the change in
$\potvarn$ between the two finest meshes, which is the quantity Paper~IV used
to argue that the dispersion has settled. Values in bold are those quoted
elsewhere in this paper: the Itokawa envelope in
Sections~\ref{sec:itokawa} and~\ref{sec:conclusions}, and the Eros envelope
in Sections~\ref{sec:nulls} and~\ref{sec:failures}.}
\label{tab:envelope}
\begin{tabular}{llrrrr}
\toprule
Decimator & Body & Facets & Dispersion $E$ & Slope range & Plateau \\
\midrule
Vertex clustering & Eros    & $2{,}097$--$42{,}907$ & $5.40\%$ & $21.91\%$ & $0.083\%$ \\
                  & Itokawa & $2{,}183$--$48{,}785$ & $0.95\%$ & $14.24\%$ & $0.019\%$ \\
                  & Gaspra  & $2{,}821$--$32{,}040$ & $1.48\%$ & $1.84\%$  & $0.000\%$ \\
                  & Ida     & $2{,}141$--$32{,}040$ & $4.40\%$ & $5.79\%$  & $0.009\%$ \\
\addlinespace
Quadric collapse  & Eros    & $2{,}000$--$42{,}000$ & $\mathbf{6.03\%}$ & $3.50\%$ & $0.098\%$ \\
                  & Itokawa & $2{,}000$--$42{,}000$ & $\mathbf{1.35\%}$ & $4.50\%$ & $0.008\%$ \\
                  & Gaspra  & $2{,}000$--$32{,}040$ & $2.92\%$ & $3.15\%$  & $0.000\%$ \\
                  & Ida     & $2{,}000$--$32{,}040$ & $1.91\%$ & $6.19\%$  & $0.000\%$ \\
\bottomrule
\end{tabular}
\end{table}

\subsection{What the table shows}
\label{app:envelope:reading}

\emph{The published range is reproduced at its lower end and not at its
upper.} With vertex clustering, which Paper~IV used, the four dispersion
envelopes run from $0.95\%$ to $5.40\%$ against the published
$0.9$--$5.9\%$. The lower end agrees; the upper is $0.5$ percentage points
short, and the difference is the ladder, whose rungs are set by the
decimator's own grid and do not fall at the same facet counts as Paper~IV's.
We quote $0.9$--$5.9\%$ throughout because it is the published figure, and
note here that the recomputed range is slightly narrower.

\emph{The envelope depends on the decimator, and not weakly.} On Eros it is
$5.40\%$ with clustering and $6.03\%$ with quadric collapse; on Itokawa,
$0.95\%$ and $1.35\%$. That dependence is the reason
Sections~\ref{sec:nulls:misplaced} and~\ref{sec:failures:plane} state both
values when the envelope is used as a test: the misplaced Eros plane returns
an improvement of $5.526\%$, which is $0.92$ times the quadric envelope and
rejected by it, and $1.02$ times the clustering envelope and not rejected.
A diagnostic whose verdict turns on a choice of decimation is worth less than
its arithmetic suggests, and we rely on plane provenance instead.

\emph{The slope column behaves differently under the two decimations.} With
clustering the mean slope varies by $1.8\%$ to $21.9\%$ across the four
bodies and rises monotonically with facet count on every one; with quadric
collapse it varies by $3.1\%$ to $6.2\%$ and shows no monotone trend. The
analytic ellipsoid of Section~\ref{sec:discussion:ellipsoid} finds the two
decimations equally accurate on a smooth surface, so the monotone drift under
clustering is a property of rough surfaces that the ellipsoid cannot
reproduce, and we make no claim about which decimation is closer to the truth
on a real body. What survives either way is the comparison Paper~IV drew: the
dispersion moves by a few per cent across a ladder that moves the slope by
several times as much, and the ellipsoid puts that ratio at $19\pm4$ to
$59\pm18$ depending on resolution.

\emph{The plateau is tighter than the figure the paper uses.} Between the two
finest meshes the dispersion moves by $0.008\%$ on Itokawa and by at most
$0.098\%$ on any of the four bodies, with either decimator. Paper~IV reports
$0.60\%$ for this quantity, and the curvature sensitivity intervals of
Section~\ref{sec:itokawa} are evaluated at that level. Retaining the
published figure makes those intervals conservative by roughly a factor of
six, which is the direction in which an interval should err.

\subsection{Scope}
\label{app:envelope:scope}

Four bodies, six rungs, two decimators. The envelope is not known for Bennu,
67P or Arrokoth, whose meshes were decimated once each rather than laddered,
so the comparisons in Section~\ref{sec:nulls} for those three bodies use the
cross-body range and not a body-specific value. Extending the ladder to them
would cost one batch of runs and would sharpen three of the four
weak-signal cases; we did not do it, and a reader who wants a body-specific
envelope for those bodies should treat its absence as a gap rather than as
evidence that the cross-body range applies.


\begin{thebibliography}{99}

% ---- companion papers in this series ----
\bibitem[Kanjanatarayont(2026a)]{paper1}
Kanjanatarayont, P., 2026. Surface gravity of small bodies from the
area-equivalent radius: an $e^{2}$-cancellation theorem and an
{$\eta \approx 1.05$} domain of validity. Companion paper, Paper I.
doi:10.5281/zenodo.21327611.

\bibitem[Kanjanatarayont(2026b)]{paper2}
Kanjanatarayont, P., 2026. The optimal area--volume radius blend for the
mean surface gravity of small bodies: a non-analytic $\Pthree/\Ptwo$ law.
Companion paper, Paper II. doi:10.5281/zenodo.21328178.

\bibitem[Kanjanatarayont(2026c)]{paper3}
Kanjanatarayont, P., 2026. Shape-based estimation of mean surface gravity
for small bodies: the accuracy of the blended radius and the limits of
convex models. Companion paper, Paper III. doi:10.5281/zenodo.21622962.

\bibitem[Kanjanatarayont(2026d)]{paper4}
Kanjanatarayont, P., 2026. Spatial surface gravity of small bodies from
polyhedral shape models: slope, geopotential, and the resolution
robustness of geopotential dispersion. Companion paper, Paper IV.
doi:10.5281/zenodo.21793575.

% ---- origin of the relaxation argument and of the technique ----
\bibitem[Richardson \& Bowling(2014)]{richardson2014}
Richardson, J.~E., Bowling, T.~J., 2014. Investigating the combined
effects of shape, density, and rotation on small body surface slopes and
erosion rates. \emph{Icarus} 234, 53--65.
doi:10.1016/j.icarus.2014.02.015.

\bibitem[Richardson et~al.(2019)]{richardson2019}
Richardson, J.~E., Graves, K.~J., Harris, A.~W., Bowling, T.~J., 2019.
Small body shapes and spins reveal a prevailing state of maximum
topographic stability. \emph{Icarus} 329, 207--221.
doi:10.1016/j.icarus.2019.03.027.

% ---- extension to interior density distribution ----
\bibitem[Kanamaru \& Sasaki(2019)]{kanamaru2019}
Kanamaru, M., Sasaki, S., 2019. Estimation of interior density
distribution for small bodies: the case of asteroid Itokawa.
\emph{Transactions of the Japan Society for Aeronautical and Space
Sciences, Aerospace Technology Japan} 17, 270--275.
doi:10.2322/tastj.17.270.

\bibitem[Kanamaru et~al.(2019)]{kanamaru2019pss}
Kanamaru, M., Sasaki, S., Wieczorek, M., 2019. Density distribution of
asteroid 25143 Itokawa based on smooth terrain shape. \emph{Planetary and
Space Science} 174, 32--42. doi:10.1016/j.pss.2019.05.002.

\bibitem[Motooka \& Kawaguchi(2012)]{motooka2012}
Motooka, N., Kawaguchi, J., 2012. Study on the gravity field on the
surface of Itokawa. Proceedings of the Japan Society for Aeronautical and
Space Sciences. Full text not available to the present author; described
here only as reported by \citet{kanamaru2019}.

% ---- polyhedral gravity method ----
\bibitem[Werner \& Scheeres(1996)]{werner1996}
Werner, R.~A., Scheeres, D.~J., 1996. Exterior gravitation of a
polyhedron derived and compared with harmonic and mascon gravitation
representations of asteroid 4769 Castalia. \emph{Celestial Mechanics
and Dynamical Astronomy} 65, 313--344. doi:10.1007/BF00053511.

\bibitem[Tsoulis(2012)]{tsoulis2012}
Tsoulis, D., 2012. Analytical computation of the full gravity tensor
of a homogeneous arbitrarily shaped polyhedral source. \emph{Geophysics}
77, F1--F11. doi:10.1190/geo2010-0334.1.

\bibitem[Schuhmacher et~al.(2024)]{gravitylib}
Schuhmacher, J., Blazquez, E., Gratl, F., Izzo, D., G\'omez, P., 2024.
Efficient polyhedral gravity modeling in modern C++ and Python.
\emph{Journal of Open Source Software} 9(98), 6384.
doi:10.21105/joss.06384. Software: \texttt{polyhedral-gravity}, version~3.3.1.

\bibitem[Sutherland \& Hodgman(1974)]{sutherland1974}
Sutherland, I.~E., Hodgman, G.~W., 1974. Reentrant polygon clipping.
\emph{Communications of the ACM} 17, 32--42. doi:10.1145/360767.360802.

\bibitem[Gaskell et~al.(2008)]{gaskell2008}
Gaskell, R.~W., Barnouin-Jha, O.~S., Scheeres, D.~J., et~al., 2008.
Characterizing and navigating small bodies with imaging data.
\emph{Meteoritics \& Planetary Science} 43, 1049--1061.
doi:10.1111/j.1945-5100.2008.tb00692.x.

\bibitem[Hergenrother et~al.(2019)]{hergenrother2019}
Hergenrother, C.~W., Maleszewski, C.~K., Nolan, M.~C., et~al., 2019. The
operational environment and rotational acceleration of asteroid (101955)
Bennu from OSIRIS-REx observations. \emph{Nature Communications} 10, 1291.
doi:10.1038/s41467-019-09213-x.

\bibitem[Kaasalainen et~al.(2003)]{kaasalainen2003}
Kaasalainen, M., Kwiatkowski, T., Abe, M., et~al., 2003. CCD photometry and
model of MUSES-C target (25143) 1998~SF36. \emph{Astronomy \& Astrophysics}
405, L29--L32. doi:10.1051/0004-6361:20030844.

\bibitem[Miller et~al.(2002)]{miller2002}
Miller, J.~K., Konopliv, A.~S., Antreasian, P.~G., et~al., 2002.
Determination of shape, gravity, and rotational state of asteroid 433 Eros.
\emph{Icarus} 155, 3--17. doi:10.1006/icar.2001.6753.

\bibitem[Nishihara et~al.(2019)]{nishihara2019}
Nishihara, S., Abe, M., Hasegawa, S., et~al., 2019. Ground-based lightcurve
observation campaign of (25143) Itokawa between 2001 and 2004.
arXiv:1802.02742.

\bibitem[Nolan et~al.(2013)]{nolan2013}
Nolan, M.~C., Magri, C., Howell, E.~S., et~al., 2013. Shape model and surface
properties of the OSIRIS-REx target asteroid (101955) Bennu from radar and
lightcurve observations. \emph{Icarus} 226, 629--640.
doi:10.1016/j.icarus.2013.05.028.

\bibitem[Keane et~al.(2022)]{keane2022}
Keane, J.~T., Porter, S.~B., Beyer, R.~A., et~al., 2022. The geophysical
environment of (486958) Arrokoth---a small Kuiper Belt object explored by
New Horizons. \emph{Journal of Geophysical Research: Planets} 127,
e2021JE007068. doi:10.1029/2021JE007068.

\bibitem[Scheeres et~al.(2019)]{scheeres2019}
Scheeres, D.~J., McMahon, J.~W., French, A.~S., et~al., 2019. The dynamic
geophysical environment of (101955) Bennu based on OSIRIS-REx measurements.
\emph{Nature Astronomy} 3, 352--361. doi:10.1038/s41550-019-0721-3.

\bibitem[Nakamura et~al.(2011)]{nakamura2011}
Nakamura, T., Noguchi, T., Tanaka, M., et~al., 2011. Itokawa dust particles:
a direct link between S-type asteroids and ordinary chondrites.
\emph{Science} 333, 1113--1116. doi:10.1126/science.1207758.

\bibitem[MacMillan(1930)]{macmillan1930}
MacMillan, W.~D., 1930. \emph{The Theory of the Potential}.
McGraw-Hill, New York.

% ---- Itokawa ----
\bibitem[Fujiwara et~al.(2006)]{fujiwara2006}
Fujiwara, A., et al., 2006. The rubble-pile asteroid Itokawa as observed
by Hayabusa. \emph{Science} 312, 1330--1334. doi:10.1126/science.1125841.

\bibitem[Abe et~al.(2006)]{abe2006}
Abe, S., et al., 2006. Mass and local topography measurements of Itokawa
by Hayabusa. \emph{Science} 312, 1344--1349. doi:10.1126/science.1126272.

\bibitem[Gaskell et~al.(2008)]{gaskell2008}
Gaskell, R.~W., et al., 2008. Characterizing and navigating small bodies
with imaging data. \emph{Meteoritics \& Planetary Science} 43, 1049--1061.
doi:10.1111/j.1945-5100.2008.tb00692.x.

\bibitem[Lowry et~al.(2014)]{lowry2014}
Lowry, S.~C., et al., 2014. The internal structure of asteroid (25143)
Itokawa as revealed by detection of YORP spin-up. \emph{Astronomy \&
Astrophysics} 562, A48. doi:10.1051/0004-6361/201322602.

\bibitem[Scheeres \& Gaskell(2008)]{scheeres2008}
Scheeres, D.~J., Gaskell, R.~W., 2008. Effect of density inhomogeneity on YORP: The case of Itokawa. \emph{Icarus} 198, 125--129.
doi:10.1016/j.icarus.2008.07.010.

\bibitem[Tatsumi \& Sugita(2018)]{tatsumi2018}
Tatsumi, E., Sugita, S., 2018. Cratering efficiency on coarse-grain
targets: implications for the dynamical evolution of asteroid 25143
Itokawa. \emph{Icarus} 300, 227--248. doi:10.1016/j.icarus.2017.09.004.

% ---- Eros ----
\bibitem[Yeomans et~al.(2000)]{yeomans2000}
Yeomans, D.~K., et al., 2000. Radio science results during the
NEAR-Shoemaker spacecraft rendezvous with Eros. \emph{Science}
289, 2085--2088. doi:10.1126/science.289.5487.2085.

\bibitem[Zuber et~al.(2000)]{zuber2000}
Zuber, M.~T., et al., 2000. The shape of 433 Eros from the NEAR-Shoemaker
laser rangefinder. \emph{Science} 289, 2097--2101.
doi:10.1126/science.289.5487.2097.

% ---- Bennu ----
\bibitem[Scheeres et~al.(2020)]{scheeres2020}
Scheeres, D.~J., et al., 2020. Heterogeneous mass distribution of the
rubble-pile asteroid (101955) Bennu. \emph{Science Advances} 6,
eabc3350. doi:10.1126/sciadv.abc3350.

\bibitem[Barnouin et~al.(2019)]{barnouin2019}
Barnouin, O.~S., et al., 2019. Shape of (101955) Bennu indicative of a
rubble pile with internal stiffness. \emph{Nature Geoscience} 12,
247--252. doi:10.1038/s41561-019-0330-x.

% ---- 67P ----
\bibitem[P\"atzold et~al.(2016)]{patzold2016}
P\"atzold, M., et al., 2016. A homogeneous nucleus for comet
67P/Churyumov--Gerasimenko from its gravity field. \emph{Nature} 530,
63--65. doi:10.1038/nature16535.

\bibitem[Preusker et~al.(2017)]{preusker2017}
Preusker, F., et al., 2017. The global meter-level shape model of comet
67P/Churyumov--Gerasimenko. \emph{Astronomy \& Astrophysics} 607, L1.
doi:10.1051/0004-6361/201731798.

\bibitem[Jorda et~al.(2016)]{jorda2016}
Jorda, L., et al., 2016. The global shape, density and rotation of comet
67P/Churyumov--Gerasimenko from preperihelion Rosetta/OSIRIS observations.
\emph{Icarus} 277, 257--278. doi:10.1016/j.icarus.2016.05.002.

% ---- Arrokoth ----
\bibitem[Porter et~al.(2024)]{porter2024}
Porter, S.~B., et al., 2024. The shape and pole of (486958) Arrokoth.
\emph{The Planetary Science Journal} 5, 260. doi:10.3847/PSJ/ad7de6.

\bibitem[Spencer et~al.(2020)]{spencer2020}
Spencer, J.~R., et al., 2020. The geology and geophysics of Kuiper Belt
object (486958) Arrokoth. \emph{Science} 367, eaay3999.
doi:10.1126/science.aay3999.

\bibitem[McKinnon et~al.(2020)]{mckinnon2020}
McKinnon, W.~B., et al., 2020. The solar nebula origin of (486958)
Arrokoth, a primordial contact binary in the Kuiper Belt. \emph{Science}
367, eaay6620. doi:10.1126/science.aay6620.

% ---- meteorite densities ----
\bibitem[Consolmagno et~al.(2008)]{consolmagno2008}
Consolmagno, G.~J., Britt, D.~T., Macke, R.~J., 2008. The significance of
meteorite density and porosity. \emph{Chemie der Erde} 68, 1--29.
doi:10.1016/j.chemer.2008.01.003.

% ---- alternative approaches to interior structure ----
\bibitem[Izzo \& G\'omez(2022)]{izzo2022}
Izzo, D., G\'omez, P., 2022. Geodesy of irregular small bodies via neural
density fields. \emph{Communications Engineering} 1, 48.
doi:10.1038/s44172-022-00050-3.

\bibitem[Takahashi \& Scheeres(2014)]{takahashi2014}
Takahashi, Y., Scheeres, D.~J., 2014. Morphology driven density
distribution estimation for small bodies. \emph{Icarus} 233, 179--193.
doi:10.1016/j.icarus.2014.02.004.

% ---- rubble piles / general ----
\bibitem[Walsh(2018)]{walsh2018}
Walsh, K.~J., 2018. Rubble pile asteroids. \emph{Annual Review of
Astronomy and Astrophysics} 56, 593--624.
doi:10.1146/annurev-astro-081817-052013.

\bibitem[Dobrovolskis(2019)]{dobrovolskis2019}
Dobrovolskis, A.~R., 2019. Classification of ellipsoids by shape
and surface gravity. \emph{Icarus} 321, 891--928.
doi:10.1016/j.icarus.2018.11.023.

\end{thebibliography}

\end{document}


